\documentclass{notes}
\usepackage[utf8]{inputenc}
\usepackage{mathtools}
\usepackage{mathrsfs}
\usepackage{enumitem}
\usepackage{adjustbox}
\usepackage{pdflscape}
\usepackage{longtable}
\usepackage{float}
\usepackage{needspace}
\usepackage{microtype}
\usepackage{booktabs}
\usepackage{array}
\usepackage{tikz}
\usetikzlibrary{positioning,arrows.meta,shapes.geometric,fit,calc,matrix}
\usepackage{mathmongo-macros}
\hypersetup{colorlinks=true,linkcolor=blue!55!black,citecolor=teal!55!black,urlcolor=blue!55!black}

\newcommand{\A}{\mathcal A}
\newcommand{\B}{\mathcal B}
\newcommand{\K}{\mathcal K}
\newcommand{\Calkin}{\mathcal Q}
\newcommand{\GL}{\operatorname{GL}}
\newcommand{\spn}{\operatorname{span}}
\newcommand{\alg}{\operatorname{alg}}
\newcommand{\Id}{I}
\newcommand{\spec}{\sigma}
\newcommand{\rad}{r}
\newcommand{\bbD}{\mathbb D}
\newcommand{\bbT}{\mathbb T}
\newcommand{\SO}{\operatorname{SO}}
\newcommand{\PSO}{\operatorname{PSO}}
\newcommand{\ran}{\operatorname{Ran}}
\newcommand{\Ker}{\operatorname{Ker}}
\newcommand{\Ind}{\operatorname{Ind}}
\newcommand{\clos}{\overline}

\newcommand{\origen}[1]{%
  \par\smallskip{\footnotesize\color{DIgray}\textbf{Correspondencia con las notas base:} #1.}\par\smallskip
}

\newtcolorbox{dependencias}[1][]{
  notesbase,
  title={Dependencias lógicas},
  colback=DIlightblue,
  colframe=DIblue,
  #1
}
\newtcolorbox{investigacion}[1][]{
  notesbase,
  title={Ejemplo tomado de la línea de investigación Haseman/SIO},
  colback=DIlightgreen,
  colframe=DIgreen,
  #1
}
\newtcolorbox{puente}[1][]{
  notesbase,
  title={Puente hacia la Sección 2},
  colback=DIlightpurple,
  colframe=DIpurple,
  #1
}
\newtcolorbox{criterio}[1][]{
  notesbase,
  title={Qué conviene recordar},
  colback=DIlightorange,
  colframe=DIorange,
  #1
}

\tikzset{
  flowbox/.style={draw=DIblue, rounded corners=3pt, very thick, fill=DIlightblue,
    align=center, text width=3.0cm, minimum height=1.15cm, inner sep=6pt},
  flowbox2/.style={draw=DIteal, rounded corners=3pt, very thick, fill=BootTealLight,
    align=center, text width=3.0cm, minimum height=1.15cm, inner sep=6pt},
  flowbox3/.style={draw=DIpurple, rounded corners=3pt, very thick, fill=DIlightpurple,
    align=center, text width=3.0cm, minimum height=1.15cm, inner sep=6pt},
  flowend/.style={draw=DIgreen, rounded corners=3pt, very thick, fill=DIlightgreen,
    align=center, text width=3.0cm, minimum height=1.15cm, inner sep=6pt},
  taskbox/.style={draw=DIorange, rounded corners=3pt, thick, fill=DIlightorange,
    align=center, text width=3.1cm, minimum height=0.95cm, inner sep=4pt},
  theorybox/.style={draw=DIblue, rounded corners=3pt, thick, fill=DIlightblue,
    align=center, text width=3.15cm, minimum height=0.95cm, inner sep=4pt},
  researchbox/.style={draw=DIgreen, rounded corners=3pt, thick, fill=DIlightgreen,
    align=center, text width=3.2cm, minimum height=0.95cm, inner sep=4pt},
  arr/.style={-{Latex[length=2.6mm]},thick,draw=DIgray},
  mergearr/.style={-{Latex[length=2.6mm]},thick,draw=DIpurple}
}

\newcommand{\proofdiagram}[4]{%
\begin{center}
\begin{tikzpicture}[node distance=7mm]
  \node[flowbox] (h) {\textbf{Hipótesis / datos}\\#1};
  \node[flowbox2,right=of h] (r) {\textbf{Resultados previos}\\#2};
  \node[flowbox3,right=of r] (m) {\textbf{Mecanismo central}\\#3};
  \node[flowend,right=of m] (c) {\textbf{Conclusión}\\#4};
  \draw[arr] (h) -- (r);
  \draw[arr] (r) -- (m);
  \draw[arr] (m) -- (c);
\end{tikzpicture}
\end{center}
}

\newcommand{\proofdiagrambranch}[5]{%
\begin{center}
\begin{tikzpicture}[node distance=7mm]
  \node[flowbox] (h1) {\textbf{Dato 1}\\#1};
  \node[flowbox,below=of h1] (h2) {\textbf{Dato 2}\\#2};
  \node[flowbox2,anchor=west] (r) at ($(h1.east)!0.5!(h2.east)+(9mm,0)$)
    {\textbf{Herramienta}\\#3};
  \node[flowbox3,right=of r] (m) {\textbf{Paso decisivo}\\#4};
  \node[flowend,right=of m] (c) {\textbf{Conclusión}\\#5};
  \draw[arr] (h1.east) -- (r.west);
  \draw[arr] (h2.east) -- (r.west);
  \draw[arr] (r) -- (m);
  \draw[arr] (m) -- (c);
\end{tikzpicture}
\end{center}
}

\newcommand{\examplediagram}[5]{%
\begin{center}
\begin{tikzpicture}[node distance=6mm and 7mm]
  \node[theorybox] (g) {\textbf{Concepto / resultado general}\\#1};
  \node[taskbox,right=of g] (k) {\textbf{Qué se concreta}\\#2};
  \node[researchbox,right=of k] (t) {\textbf{Expresión teórica}\\#3};
  \node[theorybox,below=of t] (s) {\textbf{Expresión simbólica}\\#4};
  \node[taskbox,left=of s] (i) {\textbf{Interpretación}\\#5};
  \draw[arr] (g) -- (k);
  \draw[arr] (k) -- (t);
  \draw[arr] (t) -- (s);
  \draw[arr] (s) -- (i);
\end{tikzpicture}
\end{center}
}

\newcommand{\tasksource}[1]{\par\smallskip\textbf{Resultados usados:} #1\par\smallskip}
\newcommand{\researchcite}[1]{\textcolor{DIgreen}{\textbf{Conexión de investigación:} #1}}

\title{Notas de estudio: teoría espectral en álgebras de Banach y $C^*$-álgebras}
\author{Enrique Díaz Ocampo}
\date{Septiembre de 2026}

\begin{document}
\maketitle
\tableofcontents

\begin{context}
Estas notas reúnen la \emph{Parte 1: Teoría espectral en álgebras de Banach}
\cite{CastillejosParte1}, la \emph{Sección 2: $C^*$-álgebras}
\cite{CastillejosCap2}, la \emph{Tarea 1} y la \emph{Tarea 2}. Cada resultado se lee desde
sus hipótesis hasta la demostración y el diagrama que la resume. Los ejemplos
elementales preceden las conexiones con la literatura: cocientes por compactos,
funciones lentamente oscilantes y componentes de Cauchy del problema de
Haseman. Cada conexión declara su espacio de acción y las hipótesis adicionales
que necesita; así se puede volver desde una aplicación a los resultados del curso.
\end{context}

\chapter{Mapa lógico de los dos temas}

\section{La idea central}

Sea $\A$ una álgebra de Banach compleja unital y $a\in\A$. La Parte 1 parte de una pregunta:
¿cuándo existe $(\lambda1-a)^{-1}$? La serie de Neumann responde cerca de un invertible; el
resolvente organiza la respuesta para todos los escalares; el espectro registra los valores donde
falla la inversión. La Sección 2 añade la identidad $C^*$ y obtiene consecuencias que no valen en
una álgebra de Banach arbitraria.

\section{Mapa de la Parte 1}

\begin{center}
\begin{tikzpicture}[node distance=9mm and 8mm]
  \node[theorybox,text width=3.25cm] (ban) {Álgebra de Banach unital\\completitud y producto};
  \node[theorybox,text width=3.25cm,right=of ban] (neu) {Serie de Neumann\\$\|x\|<1\Rightarrow(1-x)^{-1}$};
  \node[theorybox,text width=3.25cm,right=of neu] (open) {$\GL(\A)$ abierto\\inversión continua};
  \node[theorybox,text width=3.25cm,below=of neu] (bound) {Cota para $|\lambda|>\|a\|$\\espectro compacto};
  \node[theorybox,text width=3.25cm,right=of bound] (res) {Resolvente\\identidad y holomorfía};
  \node[theorybox,text width=3.25cm,below=of bound] (non) {Espectro no vacío\\Liouville + Hahn--Banach};
  \draw[arr] (ban)--(neu);
  \draw[arr] (neu)--(open);
  \draw[arr] (open)--(res);
  \draw[arr] (neu)--(bound);
  \draw[arr] (open.south west)--(bound.north east);
  \draw[arr] (res.south) |- (non.east);
  \draw[arr] (bound)--(non);
\end{tikzpicture}
\end{center}

La flecha que llega a la no vacuidad necesita dos ingredientes distintos: la holomorfía del
resolvente y su decaimiento en infinito, obtenido con la serie de Neumann. La compacticidad usa
la apertura de $\GL(\A)$ para el cierre y la misma serie para la acotación.

\begin{center}
\begin{tikzpicture}[node distance=5mm and 8mm]
  \node[theorybox,text width=3.25cm] (non) {Espectro no vacío};
  \node[theorybox,text width=3.25cm,right=of non] (gm) {Gelfand--Mazur\\álgebras de división};
  \node[theorybox,text width=3.25cm,right=of gm] (max) {Ideales maximales\\$\A/J\cong\mathbb C$ si $\A$ es conmutativa};
  \node[theorybox,text width=3.25cm,below=of max] (char) {Caracteres\\$\Omega(\A)$};
  \node[theorybox,text width=3.25cm,left=of char] (gel) {Transformada de Gelfand\\$a\mapsto\widehat a$};
  \node[theorybox,text width=3.25cm,left=of gel] (rad) {Radio espectral\\resolvente + potencias};
  \draw[arr] (non)--(gm);
  \draw[arr] (gm)--(max);
  \draw[arr] (max)--(char);
  \draw[arr] (char)--(gel);
  \draw[arr] (non)--(rad);
\end{tikzpicture}
\end{center}

El radio espectral ocupa una rama propia: su fórmula utiliza la holomorfía del resolvente,
la relación espectral para potencias y el principio de acotación uniforme; no se deduce de la
transformada de Gelfand. La identificación entre caracteres e ideales maximales requiere
conmutatividad, mientras que la teoría inicial del espectro no la requiere.

\section{Mapa de la Sección 2}

\begin{center}
\begin{tikzpicture}[node distance=5mm and 8mm]
  \node[theorybox,text width=3.25cm] (id) {Identidad $C^*$\\$\|a^*a\|=\|a\|^2$};
  \node[theorybox,text width=3.25cm,right=of id] (nr) {Para $a$ normal\\$\|a\|=r(a)$};
  \node[theorybox,text width=3.25cm,right=of nr] (uq) {Unicidad de la norma $C^*$\\por el radio de $a^*a$};
  \draw[arr] (id)--(nr);
  \draw[arr] (nr)--(uq);
\end{tikzpicture}
\end{center}

\begin{center}
\begin{tikzpicture}[node distance=5mm and 8mm]
  \node[theorybox,text width=3.25cm] (a) {$a$ normal en una $C^*$-álgebra unital};
  \node[theorybox,text width=3.25cm,right=of a] (comm) {$C^*(a,1)$ conmutativa\\por $aa^*=a^*a$};
  \node[theorybox,text width=3.25cm,right=of comm] (gn) {Gelfand--Naimark\\conmutativo};
  \node[theorybox,text width=3.25cm,below=of gn] (fc) {Cálculo funcional continuo\\$f\mapsto f(a)$};
  \node[theorybox,text width=3.25cm,left=of fc] (class) {Clasificación espectral\\y raíces positivas};
  \draw[arr] (a)--(comm);
  \draw[arr] (comm)--(gn);
  \draw[arr] (gn)--(fc);
  \draw[arr] (fc)--(class);
\end{tikzpicture}
\end{center}

La igualdad $\|a\|=r(a)$ requiere normalidad y la identidad $C^*$.
La conmutatividad de $C^*(a,1)$ se deduce directamente de
$aa^*=a^*a$. La transformada de Gelfand de la Parte 1 y la igualdad de
normas permiten la representación isométrica conmutativa. La permanencia
espectral identifica el espectro de $a$ en $C^*(a,1)$ con el espectro
ambiente. La raíz positiva y la descomposición polar usan después el cálculo
funcional, con las hipótesis de sus enunciados.

\section{Tres rutas de lectura}

\noindent\textbf{Invertibilidad y espectro.} Serie de Neumann $\Rightarrow$ estabilidad de los invertibles
$\Rightarrow$ resolvente holomorfo; junto con la estimación en infinito y Liouville, se obtiene
$\sigma_{\A}(a)\neq\varnothing$. En operadores $A\in\mathcal B(X)$, el mismo resolvente es
$(\lambda I_X-A)^{-1}$.

\noindent\textbf{Conmutatividad y funciones.} Para $\A$ conmutativa y unital, los ideales maximales
producen caracteres; éstos identifican $\sigma_{\A}(a)$ con el rango de $\widehat a$. Si $\A$
es además una $C^*$-álgebra, la representación conmutativa permite el cálculo funcional.

\noindent\textbf{Cocientes y operadores.} Para un espacio de Banach $X$, el ideal cerrado
$\mathcal K(X)\triangleleft\mathcal B(X)$ conduce al cociente de Calkin. La invertibilidad de
$A+\mathcal K(X)$ equivale a que $A$ sea Fredholm; esta equivalencia es un resultado de teoría de
operadores que se usa después de construir el cociente, no un teorema elemental sobre toda
álgebra de Banach. En el capítulo de investigación se muestran los operadores de Haseman.

\chapter{Teoría espectral en álgebras de Banach}

En los resultados unitarios de este capítulo, $1$ abrevia la unidad $1_\A$ de
la álgebra ambiente $\A$; se escribe $1_\A$ cuando se necesita fijar el
espacio al que pertenece una expresión.

\section{Álgebras de Banach, involución y ejemplos básicos}

\begin{cthdefinicion}{(Álgebra de Banach)}
Una \emph{álgebra de Banach} es un álgebra asociativa $\A$ sobre $\mathbb C$, equipada con una
norma tal que $(\A,\|\cdot\|)$ es un espacio de Banach y
\[
  \|ab\|\le \|a\|\,\|b\|,\qquad a,b\in\A.
\]
Decimos que $\A$ es \emph{unital} si existe $1_\A$ con $\|1_\A\|=1$, y \emph{conmutativa} si
$ab=ba$ para todo $a,b\in\A$.
\end{cthdefinicion}
\origen{Definición 1.1}

\begin{cthejemplo}{(Funciones acotadas como álgebra de Banach)}
Buscamos un modelo conmutativo de la definición. Sea $X$ no vacío,
localmente compacto y Hausdorff, y consideremos
\[
 C_b(X)=\{f:X\to\mathbb C: f\text{ continua y acotada}\},\qquad
 \|f\|_\infty=\sup_{x\in X}|f(x)|,
\]
con las operaciones puntuales. Para $f,g\in C_b(X)$,
$|f(x)g(x)|\le\|f\|_\infty\|g\|_\infty$ en cada $x$, y por tanto
$\|fg\|_\infty\le\|f\|_\infty\|g\|_\infty$.
El límite uniforme de funciones continuas y acotadas sigue en $C_b(X)$,
lo que da completitud; la función constante $1$ es la unidad.
Así $C_b(X)$ cumple la definición de álgebra de Banach y es conmutativa.
Si $X$ es compacto, toda función continua es acotada y $C_b(X)=C(X)$.
\end{cthejemplo}

\begin{cthejemplo}{(Composición de operadores acotados)}
Para contrastar con el producto puntual, sea $X$ un espacio de Banach no nulo.
Estudiamos $\mathcal B(X)$ (también denotada $\mathcal L(X)$), el álgebra
de los operadores lineales acotados $X\to X$, con composición y norma de
operador. La desigualdad $\|STx\|\le\|S\|\|T\|\|x\|$ implica
$\|ST\|\le\|S\|\|T\|$; además, la identidad $I_X$ es unidad y
el límite en norma de operadores acotados permanece en $\mathcal B(X)$.
El producto puede cambiar al alterar el orden: en
$\mathcal B(\mathbb C^2)=M_2(\mathbb C)$, si $E_{ij}$ denota la matriz
con un $1$ en la posición $(i,j)$ y ceros en las demás,
$E_{12}E_{21}=E_{11}$, mientras $E_{21}E_{12}=E_{22}$.
Este ejemplo muestra por qué las pruebas posteriores no pueden suponer
conmutatividad de toda el álgebra.
\end{cthejemplo}

\begin{investigacion}
\textbf{Conexión con la literatura: ejemplo directo.}
Bajo las hipótesis de Karlovich para una curva estrellada $\Gamma$ y $1<p<\infty$, el problema
de frontera conduce a un elemento de $\mathcal B(L^p(\Gamma))$:
\[
 T=V_\alpha P_+ + M_G P_-\in \mathcal B(L^p(\Gamma)),
\]
donde $V_\alpha f=f\circ\alpha$, $M_Gf=Gf$ y $P_\pm=(I\pm S_\Gamma)/2$.
La fórmula muestra por qué se estudian sumas y composiciones en una misma álgebra. En esta
geometría no se presupone $S_\Gamma^2=I$ ni que $P_\pm$ sean proyecciones algebraicas;
la reducción por Sokhotski--Plemelj se explica al final de las notas
\cite[pp.~463--464, ecs.~(1.1)--(1.2)]{Karlovich2017}.
\textbf{Qué reconocer:} $\mathcal B(L^p(\Gamma))$ concreta la álgebra abstracta de
operadores acotados; la fórmula de $T$ combina elementos de esa álgebra.
\end{investigacion}

\begin{cthdefinicion}{(Involución)}
Una \emph{involución} en un álgebra compleja $\A$ es un mapa $a\mapsto a^*$ tal que
\[
(\lambda a+b)^*=\overline\lambda a^*+b^*,\qquad
(a^*)^*=a,\qquad
(ab)^*=b^*a^*.
\]
Estas identidades se exigen para todo $a,b\in\A$ y $\lambda\in\mathbb C$.
Si además $\A$ es una álgebra de Banach y $\|a^*\|=\|a\|$ para
todo $a\in\A$, llamamos a $\A$ una \emph{álgebra de Banach involutiva}
en el sentido de estas notas.
\end{cthdefinicion}
\origen{Definición 1.4}

\begin{cthejemplo}{(Conjugación puntual como involución)}
En el álgebra $C_b(X)$ del ejemplo anterior queremos identificar una
involución concreta. Para $f\in C_b(X)$ definimos
$f^*(x)=\overline{f(x)}$. Si $f,g\in C_b(X)$, entonces
$(fg)^*(x)=\overline{f(x)g(x)}=g^*(x)f^*(x)$;
también $(f^*)^*=f$ y $\|f^*\|_\infty=\|f\|_\infty$.
La conjugación puntual satisface así la definición de involución
isométrica; la igualdad $g^*f^*=f^*g^*$ refleja que este modelo es conmutativo.
\end{cthejemplo}

\section{Invertibilidad y serie de Neumann}

\begin{cthdefinicion}{(Grupo de invertibles)}
Sea $\A$ una álgebra de Banach unital, con unidad $1_\A$. Su grupo de invertibles es
\[
 \GL(\A)=\{a\in\A: \exists b\in\A,\ ab=ba=1_\A\}.
\]
Para $a\in\GL(\A)$, ese $b$ es único y se escribe $a^{-1}$.
\end{cthdefinicion}
\origen{Definición 1.7}

\begin{cthproposicion}{(Serie de Neumann)}
Sea $\A$ una álgebra de Banach unital, con unidad $1_\A$, y sea $a\in\A$.
Si $\|1_\A-a\|<1$, entonces $a\in\GL(\A)$ y
\[
 a^{-1}=\sum_{n=0}^{\infty}(1_\A-a)^n.
\]
\end{cthproposicion}
\origen{Proposición 1.8}

\begin{proof}[Demostración]
Escribamos $q=1_\A-a$. La hipótesis da $\|q\|<1$ y, por submultiplicatividad,
$\|q^n\|\le\|q\|^n$ para todo $n\ge0$. Por ello
\[
 \sum_{n=0}^{\infty}\|q^n\|
 \le \sum_{n=0}^{\infty}\|q\|^n<\infty.
\]
La completitud de $\A$ asegura que las sumas parciales
$s_N=\sum_{n=0}^{N}q^n$ convergen en norma a un elemento $b\in\A$.
Como $a=1_\A-q$, el producto por la izquierda telescopa:
\[
 as_N=(1_\A-q)\sum_{n=0}^{N}q^n
 =\sum_{n=0}^{N}q^n-\sum_{n=1}^{N+1}q^n
 =1_\A-q^{N+1}.
\]
También debe comprobarse el producto por la derecha:
\[
 s_Na=\sum_{n=0}^{N}q^n(1_\A-q)
 =\sum_{n=0}^{N}q^n-\sum_{n=1}^{N+1}q^n
 =1_\A-q^{N+1}.
\]
Aquí sólo se ha usado que $q$ conmuta con sus propias potencias, sin exigir que
$\A$ sea conmutativa. Finalmente,
$\|q^{N+1}\|\le\|q\|^{N+1}\to0$; por la continuidad de la
multiplicación al pasar al límite en ambas identidades se obtiene
$ab=ba=1_\A$. Así $a\in\GL(\A)$ y $b=a^{-1}$.
\end{proof}

\proofdiagram
{$\A$ Banach unital, $a\in\A$; $\|1_\A-a\|<1$}
{Completitud de $\A$; $\|(1-a)^n\|\le\|1-a\|^n$}
{La serie geométrica converge y el producto telescopa}
{$a^{-1}=\sum_{n\ge0}(1_\A-a)^n$}

\begin{cthejemplo}{(Serie de Neumann para una matriz perturbada)}
Queremos invertir una matriz cercana a la identidad. En
$M_2(\mathbb C)$ con la norma operatorial, sea $A=I_2-E$ con
$\|E\|<1$. La proposición de Neumann se aplica porque
$\|I_2-A\|=\|E\|<1$, y proporciona
\[
 A^{-1}=I_2+E+E^2+\cdots.
\]
Por ejemplo, si
$E=\begin{pmatrix}0&t\\0&0\end{pmatrix}$ con $|t|<1$,
entonces $E^2=0$ y la serie termina en
$A^{-1}=I_2+E$; en efecto, $(I_2-E)(I_2+E)=I_2$.
La condición de norma asegura la inversión aun cuando la serie no termine.
\end{cthejemplo}

\begin{cthejemplo}{(Invertibilidad estable alrededor de un operador)}
Sea $X$ un espacio de Banach no nulo y supongamos conocido un operador
invertible $A_0\in\mathcal B(X)$. Queremos invertir $A_0+E$ cuando
$E\in\mathcal B(X)$ es pequeño respecto de $A_0$, en el sentido de que
$\|A_0^{-1}E\|<1$. La factorización
$A_0+E=A_0(I_X+A_0^{-1}E)$ permite aplicar Neumann al segundo factor:
\[
 (A_0+E)^{-1}=(I_X+A_0^{-1}E)^{-1}A_0^{-1}
 =\sum_{n=0}^{\infty}(-A_0^{-1}E)^n A_0^{-1}.
\]
La fórmula calcula una inversa exacta y muestra que la escala relevante de
la perturbación es $\|A_0^{-1}E\|$. La pregunta de Fredholm se formulará más
adelante mediante invertibilidad de una clase en un cociente.
\end{cthejemplo}

\begin{cthnota}{(Dónde se usan las hipótesis en la serie de Neumann)}
La desigualdad $\|1-a\|<1$ hace sumable la serie de normas; la completitud de $\A$ produce
su suma dentro de $\A$; la unidad permite escribir $a=1-(1-a)$ y efectuar el telescopaje.
No se usa conmutatividad de toda la álgebra.
\end{cthnota}

\begin{cthproposicion}{(Estabilidad local de la invertibilidad)}
Sea $\A$ una álgebra de Banach unital. El conjunto $\GL(\A)$ es abierto.
Más precisamente, si $a\in\GL(\A)$ y $b\in\A$ satisfacen
\[
 \|b-a\|<\|a^{-1}\|^{-1},
\]
entonces $b\in\GL(\A)$.
\end{cthproposicion}
\origen{Observación 1.9}

\begin{proof}[Demostración]
Pongamos $c=a^{-1}(a-b)$. La hipótesis y la submultiplicatividad dan
\[
 \|c\|\le \|a^{-1}\|\,\|a-b\|<1,
 \qquad b=a(1_\A-c).
\]
La proposición anterior da una inversa bilateral
$(1_\A-c)^{-1}=\sum_{n\ge0}c^n$ dentro de $\A$. El producto de
los dos factores invertibles es invertible, con
\[
 b^{-1}=(1_\A-c)^{-1}a^{-1}
 =\sum_{n=0}^{\infty}\bigl[a^{-1}(a-b)\bigr]^n a^{-1}.
\]
Por tanto toda bola $B(a,\|a^{-1}\|^{-1})$ está contenida en
$\GL(\A)$, que es abierto.
\end{proof}

\proofdiagram
{$\A$ unital, $a\in\GL(\A)$ y $b\in\A$ cercano a $a$}
{Serie de Neumann}
{Factorizar $b=a\,[1-a^{-1}(a-b)]$}
{$\|a^{-1}(a-b)\|<1\Rightarrow b$ invertible}

\begin{cthejemplo}{(Radio de estabilidad alrededor de $2I_n$)}
Queremos hacer concreta la bola de invertibilidad de la proposición anterior.
En $M_n(\mathbb C)$ con la norma operatorial inducida, tomamos
$a=2I_n$; entonces $a^{-1}=\frac12I_n$ y
$\|a^{-1}\|^{-1}=2$. Si una matriz $b$ satisface
$\|b-2I_n\|<2$, la condición de estabilidad local se cumple y
$b$ es invertible. El número $2$ mide aquí cuánto puede variar la
matriz antes de que este criterio deje de garantizar una inversa.
\end{cthejemplo}

\begin{cthcorolario}{(Invertibilidad fuera del disco de radio $\|a\|$)}
Sea $\A$ una álgebra de Banach compleja unital, con unidad $1_\A$,
$a\in\A$ y $\lambda\in\mathbb C$. Si $\|a\|<|\lambda|$, entonces
$\lambda1_\A-a\in\GL(\A)$.
\end{cthcorolario}
\origen{Corolario 1.10}

\begin{proof}[Demostración]
La desigualdad $\|a\|<|\lambda|$ implica $\lambda\ne0$. Para
$c=\lambda^{-1}a$ se tiene $\|c\|=\|a\|/|\lambda|<1$;
por la proposición de Neumann, $1_\A-c$ tiene la inversa bilateral
$\sum_{n\ge0}c^n$. Como
$\lambda1_\A-a=\lambda(1_\A-c)$ y el escalar
$\lambda1_\A$ es invertible, también lo es $\lambda1_\A-a$.
Más precisamente, su inversa, por ambos lados, es
\[
 (\lambda1_\A-a)^{-1}
 =\lambda^{-1}\sum_{n=0}^{\infty}(\lambda^{-1}a)^n.
\]
La serie converge en norma por $\|\lambda^{-1}a\|<1$.
\end{proof}

\proofdiagram
{$a\in\A$, $\lambda\in\mathbb C$; $|\lambda|>\|a\|$}
{$\lambda\ne0$ y $\|\lambda^{-1}a\|<1$}
{Neumann: $1_\A-\lambda^{-1}a\in\GL(\A)$}
{$\lambda1_\A-a\in\GL(\A)$}

\begin{cthejemplo}{(Invertibilidad de $2I_2-A$ fuera del disco normado)}
Comprobemos el corolario en una matriz para la que la inversa se ve
directamente. Sea
$A=\operatorname{diag}(1,-1)\in M_2(\mathbb C)$ con la norma operatorial;
$\|A\|=1$. Para $\lambda=2$ se tiene $|\lambda|>\|A\|$ y
\[
  (2I_2-A)^{-1}=\operatorname{diag}(1,1/3).
\]
La multiplicación diagonal verifica que éste es un inverso bilateral.
El cálculo anticipa que fuera del disco cerrado de radio $\|a\|$ no hay
valores espectrales, lectura que se formaliza tras definir el espectro.
\end{cthejemplo}

\section{Espectro y resolvente}

\begin{cthdefinicion}{(Espectro)}
Sea $\A$ una álgebra de Banach compleja unital, con unidad $1_\A$,
y sea $a\in\A$. El \emph{espectro de $a$ en $\A$} es
\[
 \spec_\A(a)=\{\lambda\in\mathbb C:\lambda1_\A-a\notin\GL(\A)\}.
\]
Cuando no hay ambigüedad escribimos $\spec(a)$.
\end{cthdefinicion}
\origen{Definición 1.11}

\begin{cthobservacion}{(Lectura espectral del corolario anterior)}
Sea $\A$ una álgebra de Banach compleja unital, $a\in\A$ y
$\lambda\in\mathbb C$. Con la notación recién definida, si
$|\lambda|>\|a\|$, entonces $\lambda\notin\sigma_\A(a)$.
\end{cthobservacion}

\begin{cthejemplo}{(Espectro como imagen de una función continua)}
Para relacionar la definición de espectro con un objeto familiar, sea
$X$ compacto Hausdorff no vacío y $f\in C(X)$. Queremos decidir cuándo
$\lambda1-f$ es invertible en $C(X)$. Si
$f(x)\ne\lambda$ para todo $x\in X$, la función
$x\mapsto(\lambda-f(x))^{-1}$ es continua y acotada por compacidad,
y es su inversa puntual. Si $f(x_0)=\lambda$, ningún producto por una
función de $C(X)$ puede valer $1$ en $x_0$. Así
$\spec_{C(X)}(f)=f(X)$: el espectro coincide con los valores
efectivamente alcanzados por la función.
\end{cthejemplo}

\pagebreak
\begin{cthproposicion}{(Compactación del espectro)}
Sea $\A$ una álgebra de Banach compleja unital y $a\in\A$.
Entonces $\spec_\A(a)$ es compacto y
\[
 \spec_\A(a)\subset\{\lambda\in\mathbb C:|\lambda|\le\|a\|\}.
\]
\end{cthproposicion}
\origen{Proposición 1.12}

\begin{proof}[Demostración]
Si $|\lambda|>\|a\|$, el corolario anterior da
$\lambda1_\A-a\in\GL(\A)$. Por la definición de espectro, ningún
tal $\lambda$ pertenece a $\spec_\A(a)$; éste está contenido en el
disco cerrado de radio $\|a\|$ y es, por tanto, acotado.
Además, el mapa $F:\mathbb C\to\A$, $F(\lambda)=\lambda1_\A-a$,
es continuo, pues
$\|F(\lambda)-F(\mu)\|=|\lambda-\mu|\|1_\A\|$.
La proposición de estabilidad muestra que $\GL(\A)$ es abierto, de modo que
$F^{-1}(\GL(\A))=\mathbb C\setminus\spec_\A(a)$ es abierto.
El espectro es cerrado y acotado en $\mathbb C$; por Heine--Borel, es compacto.
\end{proof}

\proofdiagram
{$\A$ unital, $a\in\A$; definición de $\spec_\A(a)$}
{Escalares grandes resolventes; $\GL(\A)$ abierto}
{El complemento del espectro es abierto y fuera del disco todo es resolvente}
{$\spec(a)$ cerrado y acotado, luego compacto}

\begin{cthejemplo}{(Espectro de una matriz triangular)}
Verifiquemos la cota espectral en una matriz que no es diagonal.
Tomamos $A=\begin{pmatrix}2&1\\0&3\end{pmatrix}\in M_2(\mathbb C)$.
La matriz $\lambda I_2-A$ es triangular y su determinante es
$(\lambda-2)(\lambda-3)$; por tanto es invertible exactamente cuando
$\lambda\notin\{2,3\}$ y $\spec(A)=\{2,3\}$.
La proposición de compactación asegura además $|\lambda|\le\|A\|$
para ambos valores: la entrada fuera de la diagonal puede cambiar la
norma, pero no estos valores espectrales.
\end{cthejemplo}

\begin{cthdefinicion}{(Conjunto y función resolvente)}
Sea $\A$ una álgebra de Banach compleja unital y $a\in\A$.
El \emph{conjunto resolvente} es $\rho_\A(a)=\mathbb C\setminus\spec_\A(a)$.
La \emph{función resolvente} es
\[
 R_a:\rho_\A(a)\to\A,\qquad R_a(z)=(z1_\A-a)^{-1}.
\]
Cuando el álgebra ambiente está clara se escribe $\rho(a)$.
\end{cthdefinicion}
\origen{Definición 1.13}

\noindent La conclusión del corolario sobre invertibilidad fuera del disco
también puede escribirse ahora como
$|\lambda|>\|a\|\Rightarrow\lambda\in\rho_\A(a)$, porque
$\rho_\A(a)=\mathbb C\setminus\spec_\A(a)$.

\begin{cthejemplo}{(Resolvente de una matriz diagonal)}
Queremos calcular la función resolvente de
$A=\operatorname{diag}(1,2)\in M_2(\mathbb C)$.
Para $z\notin\{1,2\}$, la matriz $zI_2-A$ tiene entradas diagonales
no nulas y su inversa se obtiene invirtiéndolas por separado:
\[
 R_A(z)=\operatorname{diag}\bigl((z-1)^{-1},(z-2)^{-1}\bigr).
\]
La fórmula deja de estar definida precisamente en los valores
espectrales $1$ y $2$; el resolvente codifica así dónde falla la
invertibilidad.
\end{cthejemplo}

\begin{cthejemplo}{(Resolvente visto en el álgebra de operadores)}
Situemos la misma pregunta en el álgebra $\mathcal B(X)$ de un espacio
de Banach no nulo $X$. Su unidad es $I_X$, de modo que para
$A\in\mathcal B(X)$ las definiciones anteriores se leen
\[
 \spec_{\mathcal B(X)}(A)
 =\{\lambda\in\mathbb C:\lambda I_X-A\notin\GL(\mathcal B(X))\},
 \qquad R_A(\lambda)=(\lambda I_X-A)^{-1}.
\]
Invertibilidad aquí significa la existencia de una inversa lineal
\emph{acotada} $X\to X$. Si $X=\mathbb C^2$ y
$A=\operatorname{diag}(1,2)$, recuperamos el cálculo del ejemplo
anterior. La notación abstracta permite aplicar la misma definición
a operadores en espacios de dimensión infinita.
\end{cthejemplo}

\begin{cthlema}{(Identidad y holomorfía del resolvente)}
Sea $\A$ una álgebra de Banach compleja unital y $a\in\A$.
La función $R_a:\rho_\A(a)\to\A$ es holomorfa y
\[
 R_a'(z)=-(z1-a)^{-2}.
\]
Además, para $z,w\in\rho(a)$,
\[
 R_a(z)-R_a(w)=(w-z)R_a(w)R_a(z).
\]
\end{cthlema}
\origen{Lema 1.14}

\begin{proof}[Demostración]
Escribamos $b_z=z1_\A-a$ y $b_w=w1_\A-a$. Ambos son invertibles
para $z,w\in\rho_\A(a)$, y $b_z-b_w=(z-w)1_\A$.
La identidad general para dos invertibles,
$b_z^{-1}-b_w^{-1}=b_z^{-1}(b_w-b_z)b_w^{-1}$,
se comprueba porque el lado derecho es
$b_z^{-1}b_wb_w^{-1}-b_z^{-1}b_zb_w^{-1}
=b_z^{-1}-b_w^{-1}$. Además, $b_z$
y $b_w$ conmutan, y por ello también conmutan sus inversos. Así,
\[
 R_a(z)-R_a(w)=(w-z)R_a(z)R_a(w)
                =(w-z)R_a(w)R_a(z).
\]
Para justificar también la holomorfía, fijemos $w\in\rho_\A(a)$.
Si $|z-w|\|R_a(w)\|<1$, la factorización
$b_z=b_w[1_\A+(z-w)R_a(w)]$ y Neumann dan
\[
 R_a(z)=\sum_{n=0}^{\infty}(-1)^n(z-w)^nR_a(w)^{n+1}.
\]
Esta serie converge en norma localmente alrededor de $w$, así que $R_a$
es holomorfa allí. Al dividir la identidad anterior por $z-w$ y dejar
$z\to w$ en la serie se obtiene
$R_a'(w)=-R_a(w)^2=-(w1_\A-a)^{-2}$.
\end{proof}

\proofdiagram
{$\A$ Banach unital, $a\in\A$; $z,w\in\rho_\A(a)$}
{Identidad de inversos y serie local de Neumann}
{Obtener la identidad resolvente y derivar la serie local}
{$R_a'(w)=-R_a(w)^2$}

\begin{cthejemplo}{(Derivada del resolvente escalar)}
Para comprobar el signo de la fórmula del lema, tomemos
$\A=\mathbb C$ y $a=\alpha\in\mathbb C$.
En $z\ne\alpha$, $R_a(z)=(z-\alpha)^{-1}$ y
\[
 \frac{R_a(z+h)-R_a(z)}{h}
 =-\frac{1}{(z+h-\alpha)(z-\alpha)}
 \longrightarrow -(z-\alpha)^{-2}.
\]
Es la derivada ordinaria de una función recíproca; el lema extiende
este cálculo a inversos dentro de cualquier álgebra de Banach unital.
\end{cthejemplo}

\begin{cthteorema}{(No vacuidad del espectro)}
Sea $\A$ una álgebra de Banach compleja unital y $a\in\A$. Entonces
\[
 \spec(a)\neq\varnothing.
\]
\end{cthteorema}
\origen{Teorema 1.15}

\begin{proof}[Demostración]
Supongamos $\spec_\A(a)=\varnothing$. Entonces
$R_a:\mathbb C\to\A$ es entera por el lema anterior. En particular,
$R_a(0)=(-a)^{-1}\ne0$, porque es inverso de un elemento de la
álgebra no nula. Hahn--Banach proporciona
$\varphi\in\A^*$ con $\varphi(R_a(0))\ne0$.
Para $|z|>2\|a\|$, la serie de Neumann da
\[
 R_a(z)=z^{-1}(1_\A-z^{-1}a)^{-1},\qquad
 \|R_a(z)\|\le\frac{1}{|z|}
 \sum_{n=0}^{\infty}\left(\frac{\|a\|}{|z|}\right)^n
 \le\frac{2}{|z|}.
\]
La función escalar $f=\varphi\circ R_a$ es entera y tiende a cero
cuando $|z|\to\infty$. Es acotada fuera de un disco por la estimación,
y dentro de él por continuidad y compacidad. Liouville obliga a que
$f$ sea constante; el límite en infinito hace esa constante igual a cero.
Esto contradice $f(0)=\varphi(R_a(0))\ne0$. Luego
$\spec_\A(a)\ne\varnothing$.
\end{proof}

\proofdiagrambranch
{$\A$ Banach compleja unital, $a\in\A$; suponer $\spec(a)=\varnothing$}
{$R_a$ sería entera y $R_a(z)\to0$ cuando $|z|\to\infty$}
{Hahn--Banach + Liouville aplicados a $\varphi\circ R_a$}
{Elegir $\varphi$ con $\varphi(R_a(0))\ne0$; Liouville da $\varphi\circ R_a=0$}
{$\spec(a)$ debe ser no vacío}

\begin{cthejemplo}{(Espectro sin recurrir a un determinante)}
Como caso familiar, consideremos
$A=\begin{pmatrix}1&1\\0&1\end{pmatrix}\in M_2(\mathbb C)$.
El determinante de $\lambda I_2-A$ es $(\lambda-1)^2$, de modo que
$\spec(A)=\{1\}$ y el espectro no está vacío.
El teorema de no vacuidad obtiene la misma conclusión general sin
usar determinantes, que no están disponibles para un operador
arbitrario de dimensión infinita.
\end{cthejemplo}

\begin{cthnota}{(Dónde se usan las hipótesis de la no vacuidad)}
La complejidad del cuerpo permite aplicar Liouville a $\varphi\circ R_a$; la completitud
garantiza la serie de Neumann y la holomorfía con valores en $\A$; Hahn--Banach separa un
valor no nulo de $R_a$. La unidad interviene en $\lambda1-a$. No se necesita
conmutatividad.
\end{cthnota}

\section{Caracteres, ideales maximales y transformada de Gelfand}

\begin{cthteorema}{(Rigidez de Gelfand--Mazur)}
Si $\A$ es una álgebra de Banach compleja unital y todo elemento no nulo
es invertible, entonces
$\A=\mathbb C1_\A$.
\end{cthteorema}
\origen{Teorema 1.16}

\begin{proof}[Demostración]
Fijemos $a\in\A$. La no vacuidad del espectro da
$\lambda\in\spec_\A(a)$ para algún $\lambda\in\mathbb C$.
Por definición, $\lambda1_\A-a$ no es invertible. Como todo elemento
no nulo de $\A$ sí es invertible por hipótesis, necesariamente
$\lambda1_\A-a=0$, esto es, $a=\lambda1_\A$.
Así, el homomorfismo unital $\mathbb C\to\A$,
$\lambda\mapsto\lambda1_\A$, es sobreyectivo; es inyectivo porque
$1_\A\ne0$. Identifica $\A$ con $\mathbb C1_\A$.
\end{proof}

\proofdiagram
{$\A$ compleja unital; todo $a\in\A\setminus\{0\}$ es invertible}
{No vacuidad de $\spec(a)$}
{Elegir $\lambda\in\spec(a)$; entonces $\lambda1-a$ no es invertible}
{$\lambda1-a=0$, así $a=\lambda1$}

\begin{cthejemplo}{(Los escalares como álgebra de división)}
Probemos la hipótesis del teorema en el modelo más simple:
$\A=\mathbb C$ con el valor absoluto. Cada $z\ne0$ tiene el inverso
$z^{-1}=1/z$, y el espectro de $\alpha\in\mathbb C$ es
$\{\alpha\}$ porque $\lambda-\alpha$ deja de ser invertible sólo
cuando vale cero. Gelfand--Mazur afirma que cualquier álgebra de
Banach compleja unital en la que todo elemento no nulo sea invertible
es, algebraicamente, este mismo modelo escalar.
\end{cthejemplo}

\begin{cthdefinicion}{(Carácter y espectro de Gelfand)}
Un \emph{carácter} de una álgebra de Banach conmutativa unital $\A$ es un funcional lineal
multiplicativo no nulo $\varphi:\A\to\mathbb C$. El conjunto de caracteres se denota
\[
 \Omega(\A)=\{\varphi:\A\to\mathbb C:\varphi\text{ carácter}\},
\]
como conjunto. Después de probar continuidad de los caracteres se le dará
la topología débil-$*$ heredada de $\A^*$.
\end{cthdefinicion}
\origen{Definición 1.17}

\begin{cthejemplo}{(La evaluación puntual es un carácter)}
Sea $X$ compacto Hausdorff no vacío y fijemos $x\in X$.
Queremos comprobar que evaluar en $x$ cumple la definición de carácter
de $C(X)$. El funcional $\operatorname{ev}_x(f)=f(x)$ es lineal y
\[
 \operatorname{ev}_x(fg)=f(x)g(x)
 =\operatorname{ev}_x(f)\operatorname{ev}_x(g),\qquad
 \operatorname{ev}_x(1)=1.
\]
Por tanto es multiplicativo y no nulo. Cada punto produce así un
carácter; la Tarea 2 demostrará que en $C(X)$ no hay otros.
\end{cthejemplo}

\begin{cthproposicion}{(Los caracteres tienen norma uno)}
Sea $\A$ una álgebra de Banach compleja conmutativa unital y
$\varphi\in\Omega(\A)$. Entonces $\varphi$ es continuo y
\[
 \|\varphi\|=\varphi(1_\A)=1.
\]
\end{cthproposicion}
\origen{Proposición 1.18}

\begin{proof}[Demostración]
Como $\varphi\ne0$, existe $b\in\A$ con $\varphi(b)\ne0$.
La igualdad $\varphi(b)=\varphi(1_\A b)=\varphi(1_\A)\varphi(b)$
obliga a $\varphi(1_\A)=1$. Fijemos $a\in\A$.
Si $|\varphi(a)|>\|a\|$, el corolario de Neumann muestra que
$u=\varphi(a)1_\A-a$ es invertible. Pero
\[
 \varphi(u)=\varphi(a)\varphi(1_\A)-\varphi(a)=0.
\]
Si $v=u^{-1}$, la multiplicatividad daría
$1=\varphi(1_\A)=\varphi(uv)=\varphi(u)\varphi(v)=0$,
una contradicción. Por tanto $|\varphi(a)|\le\|a\|$ para todo
$a\in\A$. Esta cota prueba la continuidad de $\varphi$ y
$\|\varphi\|\le1$; como $\|1_\A\|=\varphi(1_\A)=1$,
se obtiene $\|\varphi\|\ge1$ y finalmente $\|\varphi\|=1$.
\end{proof}

\proofdiagram
{$\A$ conmutativa unital, $\varphi\in\Omega(\A)$, $a\in\A$}
{$\varphi(1)=1$; escalares grandes son resolventes}
{Si $|\varphi(a)|>\|a\|$, entonces $\varphi(a)1-a$ sería invertible pero su imagen por $\varphi$ es $0$}
{$|\varphi(a)|\le\|a\|$ y $\|\varphi\|=1$}

\begin{cthejemplo}{(Norma del funcional de evaluación en $C(X)$)}
Sea $X$ compacto Hausdorff no vacío, fijemos $x\in X$ y consideremos
$\operatorname{ev}_x:C(X)\to\mathbb C$, $\operatorname{ev}_x(f)=f(x)$.
Queremos calcular su norma. Para cualquier $f\in C(X)$,
\[
 |\operatorname{ev}_x(f)|=|f(x)|\le\|f\|_\infty,
\]
de donde $\|\operatorname{ev}_x\|\le1$. La función constante $1$
tiene norma uniforme uno y $\operatorname{ev}_x(1)=1$, por lo que
$\|\operatorname{ev}_x\|\ge1$. Concluimos
$\|\operatorname{ev}_x\|=1$, en acuerdo con la proposición general
sobre caracteres: evaluar en un punto es un funcional continuo
normalizado.
\end{cthejemplo}

\begin{cthdefinicion}{(Topología del espacio de caracteres)}
Sea $\A$ una álgebra de Banach compleja conmutativa unital.
Como todo $\varphi\in\Omega(\A)$ es continuo, $\Omega(\A)\subset\A^*$.
Su \emph{topología de Gelfand} es la topología débil-$*$
restringida desde $\A^*$: una red $\varphi_i$ converge a $\varphi$ si
$\varphi_i(a)\to\varphi(a)$ para todo $a\in\A$.
Con esta topología, $\Omega(\A)$ es compacto y Hausdorff.
\end{cthdefinicion}
\origen{Definición 1.19}

\begin{proof}[Justificación de compacidad]
Por la proposición anterior, todo carácter tiene norma uno, así que
$\Omega(\A)$ está dentro de la bola unidad de $\A^*$. Si una red de
caracteres $\varphi_i$ converge débil-$*$ a $\varphi$ en esa bola,
entonces $\varphi(1_\A)=\lim_i\varphi_i(1_\A)=1$, de modo que
$\varphi\ne0$. El límite puntual conserva la linealidad y, para
$a,b\in\A$,
\[
 \varphi(ab)=\lim_i\varphi_i(ab)
 =\lim_i\bigl(\varphi_i(a)\varphi_i(b)\bigr)
 =\varphi(a)\varphi(b).
\]
Así $\Omega(\A)$ es cerrado en la bola débil-$*$ compacta por
Banach--Alaoglu. La topología débil-$*$ es Hausdorff porque dos
funcionales distintos difieren en algún elemento de $\A$.
\end{proof}

\begin{cthdefinicion}{(Ideal bilateral y álgebra cociente)}
Sea $\A$ una álgebra de Banach. Un subespacio $J\subset\A$ es un
\emph{ideal bilateral} si $aJ\subset J$ y
$Ja\subset J$ para todo $a\in\A$. Sus clases $a+J$ forman la álgebra cociente
$\A/J$, con $(a+J)(b+J)=ab+J$. Si $J$ es cerrado, la norma
\[
 \|a+J\|_{\A/J}=\inf_{j\in J}\|a+j\|
\]
hace de $\A/J$ una álgebra de Banach; la demostración de completitud se desarrolla en la
Tarea 1, Problema 4. Un ideal maximal es un ideal propio que no está contenido en otro
ideal propio más grande.
\end{cthdefinicion}

\begin{cthejemplo}{(Clases y norma en un cociente sencillo)}[unbreakable]
En aritmética, $\mathbb Z/n\mathbb Z$ enseña a identificar elementos
que difieren en un ideal. Veamos la versión normada con
$X=\{p,q\}$ discreto, $\A=C(X)\cong\mathbb C^2$ y
$J=\{f:f(p)=0\}$. Dos funciones representan la misma clase
exactamente cuando tienen el mismo valor en $p$. Además,
\[
 \|f+J\|_{\A/J}
 =\inf_{h\in J}\max\{|f(p)|,|f(q)+h(q)|\}
 =|f(p)|,
\]
pues podemos elegir $h(q)=-f(q)$. El cociente conserva sólo la
coordenada $p$; la norma cociente mide el tamaño de esa información
que permanece.
\end{cthejemplo}

\begin{cthproposicion}{(Ideales maximales y cocientes escalares)}
Sea $J$ un ideal maximal de una álgebra de Banach unital $\A$.
Entonces $J$ es cerrado. Si $\A$ es conmutativa, además
\[
 \A/J\cong\mathbb C.
\]
\end{cthproposicion}
\origen{Proposición 1.20}

\begin{proof}[Demostración]
La clausura $\overline J$ sigue siendo un ideal bilateral: si
$j_n\in J$ converge a $j$ y $a\in\A$, la continuidad de la
multiplicación da $aj_n\to aj$ y $j_na\to ja$, ambos en
$\overline J$. Si $J$ no fuese cerrado, la maximalidad exigiría
$\overline J=\A$. Entonces existiría $j\in J$ con
$\|1_\A-j\|<1$. Neumann haría invertible a $j$, y
$1_\A=j^{-1}j\in J$, contradiciendo que $J$ es propio.
Luego $J$ es cerrado y $\A/J$ es una álgebra de Banach unital no nula.
Su unidad es $1_\A+J$ y tiene norma uno: la norma cociente es a lo más
$\|1_\A\|=1$, mientras una distancia menor que uno de $1_\A$ a $J$
produciría, por Neumann, un elemento invertible de $J$ y forzaría
$J=\A$.

Supongamos ahora que $\A$ es conmutativa. Los ideales de $\A/J$
se corresponden, mediante la proyección cociente, con los ideales
de $\A$ que contienen $J$; por maximalidad, $\A/J$ no tiene
ideales distintos de $\{0\}$ y de sí misma. Sea $a+J\ne0$.
El ideal $(\A/J)(a+J)$ es no nulo y, por tanto, es todo $\A/J$.
Existe $b+J$ con $(b+J)(a+J)=1_\A+J$. La conmutatividad da
también $(a+J)(b+J)=1_\A+J$, así que todo elemento no nulo
del cociente es invertible. Gelfand--Mazur proporciona
$\A/J\cong\mathbb C$ como álgebras complejas unitarias.
\end{proof}

\proofdiagrambranch
{$\A$ Banach unital, $J\subset\A$ maximal}
{En la segunda afirmación, $\A$ conmutativa}
{La clausura de $J$ sigue siendo ideal; $\A/J$ no tiene ideales propios}
{Todo elemento no nulo de $\A/J$ es invertible; aplicar Gelfand--Mazur}
{$J$ cerrado y $\A/J\cong\mathbb C$}

\begin{cthejemplo}{(Un ideal maximal definido por evaluación)}[unbreakable]
Sea $X$ compacto Hausdorff no vacío y fijemos $x\in X$.
Queremos describir el cociente que olvida las funciones nulas en $x$.
El conjunto $J_x=\{f\in C(X):f(x)=0\}$ es el núcleo del carácter
$\operatorname{ev}_x$, y la aplicación
\[
 C(X)/J_x\longrightarrow\mathbb C,\qquad f+J_x\longmapsto f(x)
\]
está bien definida, es multiplicativa y sobreyectiva porque incluye
las funciones constantes. Su núcleo es cero, así que es un
isomorfismo. Como el cociente es un cuerpo, $J_x$ es maximal:
el valor en un punto es toda la información que sobrevive.
\end{cthejemplo}

\begin{investigacion}
\textbf{Ejemplo directo de construcción.}
Sea $X$ un espacio de Banach complejo infinito-dimensional.
Un operador $K\in\mathcal B(X)$ es \emph{compacto} si
$K$ envía la bola unidad a un conjunto relativamente compacto. El conjunto
$\mathcal K(X)$ de tales operadores es un ideal bilateral cerrado de $\mathcal B(X)$:
la composición con un operador acotado preserva compacidad y el límite en norma
de compactos es compacto. El cociente
\[
 \mathcal Q(X):=\mathcal B(X)/\mathcal K(X)
\]
se llama \emph{álgebra de Calkin} \cite[p.~2]{Karlovich2025}.
Si $X$ es Hilbert, este cociente es también una
$C^*$-álgebra; en un Banach general se usa aquí sólo su estructura de álgebra de Banach.
\textbf{Qué reconocer:} $\mathcal K(X)$ cumple exactamente el papel del ideal
cerrado $J$ de la construcción $\A/J$; $\mathcal Q(X)$ es el cociente concreto.
\end{investigacion}

\begin{investigacion}
\textbf{Conexión con la literatura: uso de un cociente.}
Sea $X$ un espacio de Banach complejo infinito-dimensional.
Un operador $A\in\mathcal B(X)$ es \emph{Fredholm} si su núcleo es finito-dimensional,
su rango es cerrado y la codimensión de su rango es finita. La equivalencia de la
teoría de operadores
\[
 A\text{ Fredholm}\quad\Longleftrightarrow\quad
 A+\mathcal K(X)\in\GL(\mathcal Q(X))
\]
identifica una pregunta operatorial con una pregunta de invertibilidad en un
cociente \cite[p.~2]{Karlovich2025}. En $X=\ell^2$, sea $P$ la proyección de rango
uno sobre $\mathbb Ce_1$ y $A=I_X-P$. Entonces $A$ no es invertible, pero
$A+\mathcal K(X)=I_X+\mathcal K(X)$ sí lo es: la diferencia $P$ desaparece
en el cociente. Ésta es una elaboración docente a partir del criterio citado.
\textbf{Qué reconocer:} el cociente $\A/J$ de la Tarea 1.4 reaparece con
$\A=\mathcal B(X)$ y $J=\mathcal K(X)$; la equivalencia con Fredholmness es
un resultado adicional de teoría de operadores.
\end{investigacion}

\pagebreak
\begin{cthproposicion}{(Caracteres e ideales maximales)}
Sea $\A$ una álgebra de Banach compleja conmutativa unital.
La asignación
\[
 \Omega(\A)\longrightarrow\{\text{ideales maximales de }\A\},\qquad
 \varphi\longmapsto\ker\varphi,
\]
es una biyección.
\end{cthproposicion}
\origen{Proposición 1.21}

\begin{proof}[Demostración]
Si $\varphi\in\Omega(\A)$, entonces $\varphi(1_\A)=1$ y
$\varphi$ es sobreyectivo porque $\varphi(\lambda1_\A)=\lambda$
para todo $\lambda\in\mathbb C$. Su núcleo es un ideal propio y el
primer teorema de isomorfía da
$\A/\ker\varphi\cong\mathbb C$; como $\mathbb C$ no tiene
ideales propios no nulos, $\ker\varphi$ es maximal.

Recíprocamente, si $J$ es maximal, la proposición anterior permite
elegir un isomorfismo unital complejo
$\theta:\A/J\to\mathbb C$. Sea
$q:\A\to\A/J$ el mapa cociente. La composición
$\varphi=\theta\circ q$ es lineal, multiplicativa y unital;
es un carácter con $\ker\varphi=J$.
Para comprobar la unicidad, si $\psi$ es otro carácter con el mismo
núcleo, para cada $a\in\A$ se tiene
$a-\varphi(a)1_\A\in J=\ker\psi$. Por tanto
$\psi(a)=\varphi(a)\psi(1_\A)=\varphi(a)$. Las dos asignaciones
son así inversas una de otra.
\end{proof}

\proofdiagram
{$\A$ conmutativa unital; $\varphi\in\Omega(\A)$ o $J\subset\A$ maximal}
{Cocientes $\A/J\cong\mathbb C$}
{Núcleo de carácter tiene codimensión $1$; recíprocamente usar el mapa cociente}
{$\varphi\leftrightarrow\ker\varphi$}

\begin{cthejemplo}{(Recuperar dos puntos a partir de ideales)}
Para ver la biyección entre caracteres e ideales maximales,
tomemos $X=\{p,q\}$ y $C(X)\cong\mathbb C^2$.
Las evaluaciones dan los ideales distintos
$J_p=\{f:f(p)=0\}$ y $J_q=\{f:f(q)=0\}$.
Si $\varphi$ es cualquier carácter, las funciones indicadoras
$e_p,e_q$ satisfacen $e_p+e_q=1$ y $e_pe_q=0$;
sus imágenes son idempotentes escalares y exactamente una vale $1$.
Como $f=f(p)e_p+f(q)e_q$, se sigue que $\varphi$ es
$\operatorname{ev}_p$ o $\operatorname{ev}_q$.
En este caso, los dos puntos del espacio se recuperan como los
dos ideales maximales del álgebra.
\end{cthejemplo}

\begin{cthteorema}{(El espectro se lee con caracteres)}
Sea $\A$ una álgebra de Banach compleja conmutativa unital.
Entonces para cada $a\in\A$,
\[
 \spec(a)=\{\varphi(a):\varphi\in\Omega(\A)\}.
\]
\end{cthteorema}
\origen{Teorema 1.22}

\begin{proof}[Demostración]
Si $\lambda=\varphi(a)$ para un carácter $\varphi$, entonces
$\varphi(\lambda1_\A-a)=\lambda-\varphi(a)=0$.
Un invertible no puede tener imagen cero por un carácter unital;
por ello $\lambda1_\A-a$ no es invertible y
$\lambda\in\spec_\A(a)$.

En sentido contrario, sea $\lambda\in\spec_\A(a)$ y pongamos
$b=\lambda1_\A-a$. Como $b$ no es invertible, el ideal principal
$\A b$ es propio: si $1_\A=cb$ para algún $c\in\A$, la
conmutatividad también daría $bc=1_\A$. Por el lema de Zorn,
$\A b$ está contenido en un ideal maximal $J$: en efecto, la unión
de una cadena de ideales propios que contienen $\A b$ sigue siendo
propia, pues si contuviera $1_\A$, éste pertenecería ya a uno de
sus miembros. Esa unión es una cota superior para la cadena.
La proposición
anterior da un carácter $\varphi$ con $J=\ker\varphi$.
Como $b\in J$, se obtiene
$0=\varphi(b)=\lambda-\varphi(a)$, de donde
$\lambda=\varphi(a)$. Quedan probadas ambas inclusiones.
\end{proof}

\proofdiagrambranch
{$a\in\A$, $\varphi\in\Omega(\A)$ y $\lambda=\varphi(a)$}
{$\lambda\in\spec_\A(a)$}
{Caracteres $\leftrightarrow$ ideales maximales}
{$\lambda1-a$ cae en un ideal maximal exactamente cuando algún carácter lo anula}
{$\lambda\in\spec(a)\Longleftrightarrow\exists\varphi\in\Omega(\A):\lambda=\varphi(a)$}

\begin{cthejemplo}{(Espectro de la función cuadrado en $[0,1]$)}
Apliquemos el teorema espectral por caracteres a
$X=[0,1]$ y $f(x)=x^2\in C(X)$.
Cada evaluación $\operatorname{ev}_x$ toma en $f$ el valor $x^2$,
y el criterio de invertibilidad puntual ya comprobado para $C(X)$
descarta otros valores espectrales. Por tanto
\[
 \spec_{C(X)}(f)
 =\{\operatorname{ev}_x(f):x\in[0,1]\}
 =\{x^2:x\in[0,1]\}=[0,1].
\]
El espectro recupera exactamente el rango geométrico de la función.
\end{cthejemplo}

\begin{cthproposicion}{(Lema de Jacobson para $1-ab$)}
Sea $\A$ una álgebra de Banach unital con unidad $1_\A$ y
$a,b\in\A$. Entonces
\[
 1_\A-ab\in\GL(\A)\quad\Longleftrightarrow\quad 1_\A-ba\in\GL(\A).
\]
Si $1_\A-ab$ es invertible,
\[
 (1_\A-ba)^{-1}=1_\A+b(1_\A-ab)^{-1}a.
\]
\end{cthproposicion}
\origen{Proposición 1.23}

\begin{proof}[Demostración]
Supongamos que $1_\A-ab$ es invertible y escribamos
$c=(1_\A-ab)^{-1}$. Las dos identidades de inverso dan
$c-abc=1_\A$ y $c-cab=1_\A$.
Definamos $d=1_\A+bca$. Por la izquierda,
\[
 (1_\A-ba)d
 =1_\A-ba+bca-babca
 =1_\A+b(c-1_\A-abc)a=1_\A.
\]
Por la derecha se usa la otra identidad de $c$:
\[
 d(1_\A-ba)
 =1_\A-ba+bca-bcaba
 =1_\A+b(c-1_\A-cab)a=1_\A.
\]
Así $d$ es la inversa bilateral de $1_\A-ba$, lo que prueba
la fórmula del enunciado. Si se parte de la invertibilidad de
$1_\A-ba$, el mismo razonamiento, con $a$ y $b$ intercambiados,
demuestra la implicación contraria.
\end{proof}

\proofdiagram
{$\A$ unital, $a,b\in\A$ y $1_\A-ab$ invertible}
{Asociatividad}
{Construir explícitamente $1+b(1-ab)^{-1}a$}
{Es inverso bilateral de $1-ba$}

\begin{cthejemplo}{(La identidad de Jacobson en matrices que no conmutan)}[unbreakable]
Queremos ver la proposición sin suponer que los factores conmuten.
En $M_2(\mathbb C)$ tomamos $A=E_{12}$ y
$B=\frac12E_{21}$, donde $E_{ij}$ son las matrices con una
única entrada $1$ en la posición $(i,j)$. Entonces
$AB=\frac12E_{11}$ y $BA=\frac12E_{22}$ son distintos, pero
\[
 (I_2-AB)^{-1}=\operatorname{diag}(2,1),\qquad
 (I_2-BA)^{-1}=\operatorname{diag}(1,2).
\]
La fórmula de Jacobson produce el segundo inverso desde el primero:
$I_2+B(I_2-AB)^{-1}A=I_2+E_{22}$.
La invertibilidad simultánea proviene de la estructura del producto,
sin requerir $AB=BA$.
\end{cthejemplo}

\begin{cthcorolario}{(Espectros no nulos de $ab$ y $ba$)}
Sea $\A$ una álgebra de Banach compleja unital y $a,b\in\A$.
Entonces
\[
 \spec_\A(ab)\cup\{0\}=\spec_\A(ba)\cup\{0\},
 \qquad\text{equivalentemente}\qquad
 \spec_\A(ab)\setminus\{0\}=\spec_\A(ba)\setminus\{0\}.
\]
\end{cthcorolario}
\origen{Corolario 1.24}

\begin{proof}[Demostración]
Sea $\lambda\in\mathbb C\setminus\{0\}$ y usemos
la factorización
$\lambda1_\A-ab=\lambda(1_\A-\lambda^{-1}ab)$.
El factor escalar $\lambda1_\A$ es invertible, de modo que
\[
\lambda1-ab\in\GL(\A)
\Longleftrightarrow
1-\lambda^{-1}ab\in\GL(\A)
\Longleftrightarrow
1-\lambda^{-1}ba\in\GL(\A)
\Longleftrightarrow
\lambda1-ba\in\GL(\A).
\]
La equivalencia central es el lema anterior aplicado a
$\lambda^{-1}a$ y $b$; la última vuelve a usar el factor escalar
$\lambda\ne0$. Por la definición de espectro, $\lambda$ pertenece
a $\spec_\A(ab)$ exactamente cuando pertenece a
$\spec_\A(ba)$. La igualdad de los espectros fuera de cero se
sigue para todo $\lambda\ne0$; al agregar $\{0\}$ a ambos
conjuntos se obtiene la primera igualdad del enunciado.
\end{proof}

\proofdiagram
{$a,b\in\A$, $\lambda\in\mathbb C\setminus\{0\}$}
{Lema de Jacobson}
{$\lambda1_\A-ab=\lambda(1_\A-\lambda^{-1}ab)$}
{Coinciden los espectros no nulos; añadir $\{0\}$ da la igualdad}

\begin{cthejemplo}{(Cambio de orden entre dimensiones diferentes)}
Sea $A:\mathbb C^m\to\mathbb C^n$ y
$B:\mathbb C^n\to\mathbb C^m$. Queremos comparar los productos
$AB\in M_n(\mathbb C)$ y $BA\in M_m(\mathbb C)$.
La identidad de Jacobson se adapta a las dos unidades y, para
$\lambda\ne0$, da
\[
 (I_m-\lambda^{-1}BA)^{-1}
 =I_m+\lambda^{-1}B(I_n-\lambda^{-1}AB)^{-1}A
\]
 si $I_n-\lambda^{-1}AB$ es invertible. Como caso concreto, tomemos
$A(t)=(t,0)$ de $\mathbb C$ a $\mathbb C^2$ y
$B(x,y)=x$ de $\mathbb C^2$ a $\mathbb C$. Entonces
$AB=\operatorname{diag}(1,0)$ y $BA=(1)$:
sus espectros no nulos son ambos $\{1\}$, aunque sólo el primero
contiene $0$. Esto explica por qué el corolario agrega $\{0\}$.
\end{cthejemplo}

\begin{cthdefinicion}{(Radio espectral)}
Sea $\A$ una álgebra de Banach compleja unital y $a\in\A$.
El \emph{radio espectral} de $a$ es
\[
 r(a)=\max\{|\lambda|:\lambda\in\spec_\A(a)\}.
\]
\end{cthdefinicion}
\origen{Definición 1.25}

\begin{cthejemplo}{(Radio espectral y norma de matrices)}
Para $A=\operatorname{diag}(2,-3,i)\in M_3(\mathbb C)$,
los valores espectrales son sus entradas diagonales, así que
$r(A)=\max\{2,3,1\}=3$. Con la norma operatorial euclidiana también
$\|A\|=3$, por ser diagonal. La igualdad no es automática para matrices
no normales: si
$B=\begin{pmatrix}2&10\\0&-3\end{pmatrix}$,
su espectro sigue siendo $\{2,-3\}$ y $r(B)=3$, mientras
$\|B\|\ge\|Be_2\|=\sqrt{109}>3$.
El radio mide los valores espectrales y puede ser menor que la norma.
\end{cthejemplo}

\begin{cthteorema}{(Fórmula del radio espectral)}
Sea $\A$ una álgebra de Banach compleja unital y $a\in\A$.
Entonces
\[
 r(a)=\lim_{n\to\infty}\|a^n\|^{1/n}
     =\inf_{n\ge1}\|a^n\|^{1/n}.
\]
\end{cthteorema}
\origen{Teorema 1.26}

\begin{proof}[Demostración]
Comencemos por comprobar que el límite de normas existe. Pongamos
$\beta_n=\|a^n\|$. La submultiplicatividad da
$\beta_{m+n}\le\beta_m\beta_n$. Si $\beta_m=0$ para algún
$m\ge1$, entonces $a^n=0$ para todo $n\ge m$ y
$\lim_n\beta_n^{1/n}=\inf_{n\ge1}\beta_n^{1/n}=0$.
En otro caso, fijemos $m\ge1$ y escribamos $n=qm+s$,
$0\le s<m$. Como
$\beta_n\le\beta_m^q\beta_s$ y sólo hay un número finito de
posibles $s$, al tomar raíces $n$-ésimas y dejar $n\to\infty$
se obtiene
\[
 \limsup_{n\to\infty}\beta_n^{1/n}\le\beta_m^{1/m}.
\]
Esto vale para todo $m$. Por otra parte, cada término
$\beta_n^{1/n}$ es al menos $\inf_{m\ge1}\beta_m^{1/m}$.
En consecuencia,
\[
 L:=\lim_{n\to\infty}\|a^n\|^{1/n}
   =\inf_{n\ge1}\|a^n\|^{1/n}.
\]

Veamos que $r(a)\le L$. Si $\lambda\in\spec_\A(a)$ y
$n\ge1$, la factorización polinómica
\[
 \lambda^n1_\A-a^n
 =(\lambda1_\A-a)
   \sum_{k=0}^{n-1}\lambda^{n-1-k}a^k
 \]
tiene factores que conmutan. Si el producto fuese invertible, cada
factor lo sería: en efecto, si $x$ y $y$ conmutan y $xy$ tiene
inversa $c$, la igualdad $x(xy)=(xy)x$ implica $xc=cx$ al
multiplicar por $c$. Entonces
$x(yc)=(xy)c=1_\A$ y
$(yc)x=y(cx)=yx c=(xy)c=1_\A$, de modo que $yc$ es
una inversa bilateral de $x$.
Esto contradice que $\lambda1_\A-a$ no es invertible. Así
$\lambda^n\in\spec_\A(a^n)$ y la cota espectral ya demostrada da
$|\lambda|^n\le\|a^n\|$. Tomando el máximo sobre el espectro
y luego raíces se obtiene $r(a)\le\|a^n\|^{1/n}$ para todo
$n$, y por tanto $r(a)\le L$.

Para la otra desigualdad, sea $R=1/r(a)$ si $r(a)>0$ y
$R=\infty$ si $r(a)=0$. La función
$g(z)=(1_\A-za)^{-1}$ está definida y es holomorfa en
$|z|<R$: para $z\ne0$ en ese disco,
$|z^{-1}|>r(a)$ y
$1_\A-za=z(z^{-1}1_\A-a)$ es invertible; para $z=0$
es la unidad. Si $z\ne0$, además,
$g(z)=z^{-1}R_a(z^{-1})$, que es holomorfa por el lema del
resolvente y la holomorfía de $z\mapsto z^{-1}$ fuera de cero.
En un entorno de cero, Neumann da
$g(z)=\sum_{n\ge0}a^nz^n$.
Esta serie converge en norma para $|z|\|a\|<1$ y prueba la
holomorfía también en cero.

Fijemos $t$ real con $0<t<R$. Para cada $\varphi\in\A^*$,
$\varphi\circ g$ es holomorfa en $|z|<R$ y su serie de Taylor
en cero tiene coeficientes $\varphi(a^n)$. Elegimos $s$ con
$t<s<R$. La cota de Cauchy en el círculo $|z|=s$ da
$|\varphi(a^n)|\le M_{\varphi,s}s^{-n}$, donde
$M_{\varphi,s}=\max_{|z|=s}|\varphi(g(z))|<\infty$.
Por tanto $t^n|\varphi(a^n)|\le M_{\varphi,s}(t/s)^n$ está
acotada para cada $\varphi$. Apliquemos Banach--Steinhaus en el espacio de Banach
$\A^*$ a los operadores continuos
$T_n:\A^*\to\mathbb C$,
$T_n(\varphi)=t^n\varphi(a^n)$. La acotación puntual implica
$\sup_n\|T_n\|<\infty$. Hahn--Banach da
$\|T_n\|=t^n\|a^n\|$, así que existe $C_t<\infty$ con
$t^n\|a^n\|\le C_t$ para todo $n$. Tomando raíces y límites,
$L\le t^{-1}$. Si $r(a)>0$, hacemos $t\uparrow1/r(a)$ y
obtenemos $L\le r(a)$; si $r(a)=0$, podemos tomar $t$
arbitrariamente grande y obtenemos $L=0$. Junto con la primera
desigualdad y la existencia de $L$, esto prueba las dos igualdades.
\end{proof}

\proofdiagrambranch
{$\A$ unital, $a\in\A$; $\lambda\in\spec(a)\Rightarrow\lambda^n\in\spec(a^n)$}
{Estudiar $g(z)=(1-za)^{-1}$ en $|z|<R$, $R=1/r(a)$ o $\infty$}
{Holomorfía + Taylor + Banach--Steinhaus}
{Acotar $t^n\|a^n\|$ para cada $t<R$ y dejar $t\uparrow R$}
{$r(a)=\lim\|a^n\|^{1/n}$}

\begin{cthejemplo}{(El operador de Volterra es cuasinilpotente)}[unbreakable]
Estudiemos un operador no nulo cuyo radio espectral resulta cero.
En $C[0,1]$ con la norma uniforme, definimos
$(Vf)(x)=\int_0^x f(t)\,dt$. La iteración de la integral da, para
$n\ge1$,
\[
 (V^nf)(x)=\frac1{(n-1)!}\int_0^x(x-t)^{n-1}f(t)\,dt,
 \qquad \|V^n\|\le\frac1{n!}.
\]
La cota se obtiene al reemplazar $|f(t)|$ por $\|f\|_\infty$
e integrar la potencia. 
Por la fórmula del radio espectral,
$r(V)=\lim_n\|V^n\|^{1/n}=0$.
Como el espectro no está vacío y todos sus puntos tienen módulo
a lo más cero, $\spec(V)=\{0\}$.
Este caso distingue un operador cuasinilpotente de un operador nulo:
$V\ne0$ aunque su espectro sea $\{0\}$.
\end{cthejemplo}

\begin{cthnota}{(Interpretación de la fórmula para iteraciones)}
Sea $X$ un espacio de Banach y $A\in\mathcal B(X)$. Entonces $\|A^n\|$
mide el efecto máximo de $n$ aplicaciones
sucesivas. El límite $\lim_n\|A^n\|^{1/n}$ mide su tasa asintótica exponencial y
coincide con $r(A)$. El ejemplo de Volterra muestra que esta tasa puede ser cero
aunque $A\ne0$.
\end{cthnota}

\begin{cthdefinicion}{(Transformada de Gelfand)}
Sea $\A$ una álgebra de Banach compleja conmutativa unital.
La \emph{transformada de Gelfand} es
\[
 \Gamma:\A\to C(\Omega(\A)),\qquad
 a\mapsto\widehat a,\qquad \widehat a(\varphi)=\varphi(a).
\]
\end{cthdefinicion}
\origen{Definición 1.27}

\begin{cthejemplo}{(La transformada recupera una función de $C(X)$)}
Tomemos $X=[0,1]$ y $f(t)=t\in C(X)$ para concretar la definición.
Cada punto $x\in X$ proporciona el carácter de evaluación
$\operatorname{ev}_x$, y
\[
 \widehat f(\operatorname{ev}_x)
 =\operatorname{ev}_x(f)=x=f(x).
\]
Bajo la identificación $\Omega(C(X))\cong X$, demostrada en la Tarea 2,
la transformada de Gelfand es la propia función $f$.
La fórmula enseña cómo una función se lee a través de todos los
caracteres del álgebra.
\end{cthejemplo}

\begin{cthproposicion}{(La transformada de Gelfand es contractiva)}
Sea $\A$ una álgebra de Banach compleja conmutativa unital y
$a\in\A$. La transformada de Gelfand es lineal, multiplicativa,
unital y
\[
 \|\widehat a\|_\infty=r(a)\le\|a\|.
\]
\end{cthproposicion}
\origen{Observación 1.28}

\begin{proof}[Demostración]
Para cada $\varphi\in\Omega(\A)$, la linealidad y
multiplicatividad del carácter dan
\[
 \widehat{\alpha a+b}(\varphi)
   =\alpha\widehat a(\varphi)+\widehat b(\varphi),\qquad
 \widehat{ab}(\varphi)=\widehat a(\varphi)\widehat b(\varphi),
 \qquad \widehat{1_\A}(\varphi)=1.
\]
Cada $\widehat a$ es continua en la topología débil-$*$ por la
definición de ésta. El teorema que describe el espectro con
caracteres identifica el conjunto de valores de $\widehat a$
con $\spec_\A(a)$. Por ello
\[
 \|\widehat a\|_\infty=\sup_{\varphi\in\Omega(\A)}|\varphi(a)|
 =\max_{\lambda\in\spec(a)}|\lambda|=r(a)\le\|a\|.
\]
El máximo existe porque el espectro es compacto y no vacío;
la última desigualdad es la cota espectral previamente probada.
\end{proof}

\proofdiagram
{Definición $\widehat a(\varphi)=\varphi(a)$}
{Caracteres lineales y multiplicativos; espectro vía caracteres}
{$\|\widehat a\|_\infty=\sup_{\varphi}|\varphi(a)|$}
{$\|\widehat a\|_\infty=r(a)\le\|a\|$}

\begin{investigacion}[unbreakable]
\textbf{Conexión con la literatura: ejemplo directo de álgebra conmutativa.}
Karlovich considera $\SO(\mathbb R_+)$, la álgebra de funciones continuas acotadas cuya
oscilación en $[r,2r]$ tiende a cero cuando $r\to0$ y cuando $r\to\infty$
\cite[pp.~465--466, §2]{Karlovich2017}. Es una $C^*$-subálgebra conmutativa de
$C_b(\mathbb R_+)$: contiene constantes, es estable bajo suma, producto y conjugación, y
la condición de oscilación se conserva bajo límites uniformes. La identidad
$C^*$ se reconocerá formalmente en la Sección 2; en este punto sólo se
necesita que la álgebra es conmutativa y unital. Para $a$ en esta álgebra,
$\widehat a(\xi)=\xi(a)$ evalúa $a$ en un carácter $\xi\in\Omega(\SO(\mathbb R_+))$.
Además de las evaluaciones en $r>0$, el espacio maximal tiene fibras sobre $0$ e
$\infty$. Estas fibras organizan comportamientos límite que un solo valor $a(0)$ o
$a(\infty)$ no describiría. La transformada de Gelfand de esta álgebra de coeficientes
es un paso conceptual hacia los símbolos de Fredholm, no el símbolo de Fredholm completo.
\textbf{Qué reconocer:} los caracteres de $\Omega(\A)$ estudiados en el curso
se vuelven aquí evaluaciones ordinarias y caracteres de fibras de
$\Omega(\SO(\mathbb R_+))$. Para leer la construcción precisa, véase
\cite[§2, proposiciones~2.1--2.2]{Karlovich2017}.
\end{investigacion}

\begin{cthdefinicion}{(Ideal modular)}
Un ideal $I$ de una álgebra $\A$ se llama \emph{modular} si existe $u\in\A$ tal que
\[
 a-ua\in I,\qquad a-au\in I,\qquad a\in\A.
\]
\end{cthdefinicion}
\origen{Definición 1.29}

\begin{cthejemplo}{(Unidad local para un ideal de evaluación)}
Sea $X$ localmente compacto Hausdorff y fijemos $x\in X$.
Queremos comprobar la modularidad del ideal
\[
 I_x=\{f\in C_0(X):f(x)=0\}.
\]
Elegimos una función continua $u$ de soporte compacto con
$u(x)=1$, disponible por local compacidad y la separación de Hausdorff;
en particular, $u\in C_0(X)$.
Para todo $f\in C_0(X)$ se tiene
$(f-uf)(x)=(f-fu)(x)=0$; por tanto ambas diferencias pertenecen a
$I_x$. La clase $u+I_x$ actúa como unidad del cociente, aunque
$C_0(X)$ puede carecer de unidad global.
\end{cthejemplo}

\begin{cthproposicion}{(Criterio de modularidad mediante el cociente)}
Sea $\A$ una álgebra y $I\subsetneq\A$ un ideal bilateral propio.
Entonces $I$ es modular si y sólo si $\A/I$ tiene unidad.
\end{cthproposicion}
\origen{Proposición 1.31}

\begin{proof}[Demostración]
Si $I$ es modular, existe $u\in\A$ tal que para todo
$a\in\A$ se tiene $a-au\in I$ y $a-ua\in I$. En el
cociente, esto equivale a
\[
 (a+I)(u+I)=au+I=a+I=ua+I=(u+I)(a+I).
\]
Por tanto $u+I$ es una unidad bilateral de $\A/I$.
Recíprocamente, si $\A/I$ tiene unidad, ésta es una clase
$u+I$ para algún $u\in\A$. Para cada $a\in\A$, las
identidades de unidad dan $au+I=a+I=ua+I$, es decir,
$a-au\in I$ y $a-ua\in I$. Ese $u$ prueba que $I$ es
modular.
\end{proof}

\proofdiagram
{$I\subsetneq\A$ ideal bilateral propio}
{Definición de producto en $\A/I$}
{$u$ módulo $I$ actúa como identidad exactamente cuando $a-au,a-ua\in I$}
{$I$ modular $\Longleftrightarrow \A/I$ unital}

\begin{cthejemplo}{(Un cociente unital de $C_0(\mathbb R)$)}
Para ver la proposición en una álgebra sin unidad, tomemos
$\A=C_0(\mathbb R)$ e
$I_0=\{f\in\A:f(0)=0\}$. La función constante $1$ no pertenece
a $\A$, pero $u(t)=\max\{1-|t|,0\}$ sí pertenece y cumple $u(0)=1$.
La evaluación induce
\[
 \A/I_0\longrightarrow\mathbb C,\qquad f+I_0\longmapsto f(0),
\]
un isomorfismo sobreyectivo porque las funciones $\lambda u$ realizan
cualquier valor $\lambda$. Así $\A/I_0$ tiene unidad $u+I_0$:
la modularidad detecta una unidad que aparece sólo después del cociente.
\end{cthejemplo}

\chapter{$C^*$-álgebras: rigidez espectral y cálculo funcional}

\section{La identidad $C^*$ y clases especiales de elementos}

\begin{cthdefinicion}{($C^*$-álgebra)}
Una $C^*$-álgebra es una álgebra de Banach involutiva $\A$ que satisface
\[
 \|a^*a\|=\|a\|^2,
 \qquad a\in\A.
\]
\end{cthdefinicion}
\origen{Definición 2.1}

\begin{cthejemplo}{(Operadores en Hilbert y la identidad $C^*$)}
Queremos reconocer un modelo concreto de la definición. Si $H$ es un
espacio de Hilbert, $\mathcal B(H)$ es el álgebra de operadores lineales
acotados $H\to H$, con involución dada por el adjunto $T\mapsto T^*$.
Para $T\in\mathcal B(H)$, la desigualdad
$\|T^*T\|\le\|T^*\|\|T\|=\|T\|^2$ y la identidad
\[
 \|T\|^2=\sup_{\|x\|=1}\langle T^*Tx,x\rangle
 \le\|T^*T\|
\]
dan $\|T^*T\|=\|T\|^2$. Así la norma operatorial y el adjunto
satisfacen exactamente la identidad $C^*$.
\end{cthejemplo}

\begin{cthobservacion}{(Unidad compatible con la involución)}
Si una $C^*$-álgebra $\A$ tiene unidad $1_\A$, entonces
$1_\A^*=1_\A$ y $\|1_\A\|=1$.
\end{cthobservacion}
\origen{Observación 2.2}

\begin{proof}[Demostración]
Para todo $a\in\A$, la identidad $(a1_\A)^*=1_\A^*a^*=a^*$ muestra que
$1_\A^*$ actúa como unidad izquierda. Del mismo modo, actúa como unidad
derecha; por unicidad de la unidad, $1_\A^*=1_\A$. La identidad $C^*$ da
$\|1_\A\|^2=\|1_\A^*1_\A\|=\|1_\A\|$, y $1_\A\ne0$ implica
$\|1_\A\|=1$.
\end{proof}

\begin{cthejemplo}{(Unidad autoadjunta de norma uno en matrices)}
Ilustremos la observación anterior con la unidad de
$M_n(\mathbb C)$, $n\ge1$, dotada de la norma operatorial y del
adjunto conjugado-transpuesto. La matriz $I_n$ cumple
$I_n^*=I_n$; además, $\|I_nx\|=\|x\|$ para todo vector, de donde
$\|I_n\|=1$. La afirmación abstracta sobre $1_\A$ se convierte aquí
en una propiedad visible de la matriz identidad.
\end{cthejemplo}

\begin{investigacion}
\textbf{Ejemplo directo.} Para la estrella logarítmica abierta $\Gamma$ del problema de
Haseman, $C_b(\Gamma)$, con la norma uniforme y la conjugación puntual, es una
$C^*$-álgebra unital; la subálgebra $\SO(\Gamma)$ reúne las funciones que oscilan
lentamente en $0$ e $\infty$ \cite[§1, p.~464]{Karlovich2017}.
\textbf{Lo que conviene reconocer:} ambas son álgebras de funciones del tipo recién
definido; esta observación sólo ejemplifica la identidad $C^*$, sin intervenir en
las demostraciones del curso.
\end{investigacion}

\begin{cthdefinicion}{($*$-homomorfismo)}
Sean $\A,\B$ $C^*$-álgebras. Un $*$-homomorfismo es un homomorfismo lineal y multiplicativo
$\Phi:\A\to\B$ que satisface
\[
 \Phi(a^*)=\Phi(a)^*\qquad(a\in\A).
\]
Si ambas álgebras son unitales, diremos que $\Phi$ es unital si $\Phi(1_\A)=1_\B$.
\end{cthdefinicion}
\origen{Definición 2.4}

\begin{cthejemplo}{(Funciones como operadores de multiplicación)}
Sea $X$ localmente compacto y Hausdorff. Queremos representar el álgebra
$C_b(X)$ por operadores en $\ell^2(X)$, el espacio de familias complejas
cuadrado sumables indexadas por $X$. Para $f\in C_b(X)$ definimos
$(M_f\xi)(x)=f(x)\xi(x)$. La estimación
$\|M_f\xi\|_2\le\|f\|_\infty\|\xi\|_2$ prueba que $M_f$ es acotado.
Además, $M_fM_g=M_{fg}$ y, al comparar productos internos,
$M_f^*=M_{\overline f}$; también $M_1=I$.
Por tanto $f\mapsto M_f$ es un $*$-homomorfismo unital.
La Tarea 2.1 comprueba que preserva incluso la norma: un álgebra de
funciones puede verse fielmente como un álgebra de operadores.
\end{cthejemplo}

\begin{cthdefinicion}{(Normal, autoadjunto, unitario y proyección)}
Sea $\A$ una $C^*$-álgebra y $a\in\A$.
\begin{enumerate}[label=(\roman*)]
\item $a$ es \emph{normal} si $a^*a=aa^*$;
\item $a$ es \emph{autoadjunto} si $a^*=a$; escribimos
$\A_{\mathrm{sa}}=\{x\in\A:x=x^*\}$;
\item si $\A$ tiene unidad $1_\A$, $a$ es \emph{unitario} si
$a^*a=aa^*=1_\A$; escribimos $\mathcal U(\A)$ para el conjunto de unitarios;
\item $a$ es una \emph{proyección} si $a=a^*=a^2$; escribimos
$\mathcal P(\A)$ para el conjunto de proyecciones.
\end{enumerate}
\end{cthdefinicion}
\origen{Definición 2.5}

\begin{cthejemplo}{(Normalidad, autoadjunción, unitariedad y proyección en $M_2$)}
Busquemos un representante de cada clase de la definición en
$M_2(\mathbb C)$:
\[
\begin{array}{c|c}
\text{tipo}&\text{ejemplo}\\ \hline
\text{normal}&\begin{pmatrix}1&0\\0&i\end{pmatrix}\\[2mm]
\text{autoadjunto}&\begin{pmatrix}1&0\\0&-2\end{pmatrix}\\[2mm]
\text{unitario}&\begin{pmatrix}0&1\\1&0\end{pmatrix}\\[2mm]
\text{proyección}&\begin{pmatrix}1&0\\0&0\end{pmatrix}
\end{array}
\]
Las dos matrices diagonales conmutan con sus adjuntos; la segunda
coincide con su adjunto. Para la tercera, $U^*=U$ y $U^2=I_2$,
de modo que $U^*U=UU^*=I_2$. La cuarta satisface
$P=P^*=P^2$. Estas verificaciones muestran por qué todo
autoadjunto, unitario o proyección es normal, aunque la primera
matriz normal no sea autoadjunta.
\end{cthejemplo}

\begin{investigacion}
\textbf{Ejemplo directo en $L^2$; analogía de lectura para Haseman.}
Para $\Gamma_0=\mathbb R$ o la circunferencia
$\mathbb T=\{z\in\mathbb C:|z|=1\}$, con la medida de longitud en $L^2(\Gamma_0)$,
el operador de Cauchy satisface $S_{\Gamma_0}^2=I$ y define
\[
 P_{\Gamma_0}^\pm=\frac12(I\pm S_{\Gamma_0}).
\]
La identidad $S_{\Gamma_0}^2=I$ da $(P_{\Gamma_0}^\pm)^2=P_{\Gamma_0}^\pm$ por cálculo
algebraico. En este modelo, además, $S_{\Gamma_0}$ es autoadjunto y
$(P_{\Gamma_0}^\pm)^*=P_{\Gamma_0}^\pm$: son proyecciones ortogonales
\cite[Sec.~2, p.~22]{BFK2006}. La fórmula
$P_\pm=(I\pm S_\Gamma)/2$ reaparece en el operador
de Haseman sobre $L^p(\Gamma)$ de una estrella logarítmica abierta $\Gamma$, con
$1<p<\infty$ y los datos bajo las hipótesis de \cite{Karlovich2017},
\[
 T=V_\alpha P_+ + GP_-\in\mathcal B(L^p(\Gamma)),
\]
que aplica el shift $V_\alpha f=f\circ\alpha$ al término $P_+f$ y la
multiplicación por $G$ al término $P_-f$. En la estrella abierta,
$S_\Gamma^2=I$ y la ortogonalidad no se deducen del caso $\Gamma_0$
\cite[(1.2)]{Karlovich2017}.
\textbf{Lo que conviene reconocer:} $S^2=I$ asegura idempotencia en el modelo
Hilbertiano y $S^*=S$ añade ortogonalidad; la fórmula de Haseman se lee término a
término sin trasladar esas propiedades a la estrella abierta.
\end{investigacion}

\pagebreak
\section{Norma y espectro de elementos normales}

\begin{cthteorema}{(La norma de un normal es su radio espectral)}
Sea $\A$ una $C^*$-álgebra unital y $a\in\A$ normal. Entonces
\[
 \|a\|=r(a).
\]
\end{cthteorema}
\origen{Teorema 2.6}

\begin{proof}[Demostración]
La normalidad permite reagrupar:
\[
\|a^2\|^2
=\|(a^2)^*a^2\|
=\|(a^*)^2a^2\|
=\|(a^*a)^2\|
=\|a^*a\|^2
=\|a\|^4.
\]
Así $\|a^2\|=\|a\|^2$. Como $a^2$ también es normal, por inducción
$\|a^{2^n}\|=\|a\|^{2^n}$. Aplicando la fórmula del radio espectral sobre la subsucesión
$2^n$,
\[
 r(a)=\lim_{n\to\infty}\|a^{2^n}\|^{1/2^n}=\|a\|.
\]
\end{proof}

\proofdiagram
{$\A$ es $C^*$ unital; $a\in\A$ normal}
{Identidad $C^*$; fórmula del radio espectral}
{Probar inductivamente $\|a^{2^n}\|=\|a\|^{2^n}$}
{$r(a)=\lim_n\|a^{2^n}\|^{1/2^n}=\|a\|$}

\begin{cthejemplo}{(Norma y radio de una matriz normal)}
Comprobemos el teorema en $A=\operatorname{diag}(2,-i)$ con la
norma operatorial euclidiana. Como $A$ es diagonal, conmuta con
$A^*$ y es normal. Sus valores espectrales son $2$ y $-i$;
por tanto $r(A)=\max\{2,1\}=2$. La norma de una matriz diagonal
en esta norma es el mayor módulo de sus entradas, así que
$\|A\|=2=r(A)$. El cálculo concreta la igualdad que el teorema
garantiza para todo elemento normal.
\end{cthejemplo}

\begin{cthnota}{(Dónde intervienen las hipótesis)}
La normalidad permite pasar de $(a^*)^2a^2$ a $(a^*a)^2$.
La identidad $C^*$ convierte esa igualdad algebraica en la igualdad de normas;
la fórmula del radio espectral, demostrada en la Parte 1 para álgebras de Banach,
cierra el argumento.
\end{cthnota}


\begin{cthcorolario}{(La norma $C^*$ está determinada algebraicamente)}
Sea $\A$ una $C^*$-álgebra unital. Para cada $a\in\A$,
\[
 \|a\|=r(a^*a)^{1/2}.
\]
En particular, dos normas que hagan de la misma álgebra involutiva unital una
$C^*$-álgebra coinciden.
\end{cthcorolario}
\origen{Corolario 2.7}

\begin{proof}[Demostración]
El elemento $a^*a$ es autoadjunto y por tanto normal. Entonces
$r(a^*a)=\|a^*a\|=\|a\|^2$ por el teorema precedente y la identidad $C^*$.
Al tomar raíces no negativas obtenemos la fórmula del enunciado.
Si $\|\cdot\|_1$ y $\|\cdot\|_2$ hacen $C^*$ a la misma álgebra involutiva
unital, la invertibilidad de $a^*a-\lambda1_\A$ depende únicamente del producto
y de la unidad. Así $r(a^*a)$ es el mismo para ambas normas, y la fórmula da
$\|a\|_1=\|a\|_2$ para cada $a\in\A$.
\end{proof}

\proofdiagram
{$\A$ es $C^*$ unital; $a\in\A$; $a^*a$ es normal}
{Teorema anterior + identidad $C^*$}
{$r(a^*a)=\|a^*a\|=\|a\|^2$}
{$\|a\|=r(a^*a)^{1/2}$; unicidad de norma}

\begin{cthejemplo}{(La norma uniforme está forzada en $C(X)$)}
Queremos aplicar el corolario de unicidad de la norma a un álgebra
de funciones. Sea $X$ compacto Hausdorff y $f\in C(X)$, con la
conjugación puntual como involución. El elemento
$f^*f=|f|^2$ tiene espectro $|f|^2(X)$, por el criterio de
invertibilidad puntual de la Parte 1. La fórmula del corolario da
\[
 \|f\|=r(f^*f)^{1/2}
 =\bigl(\max_{x\in X}|f(x)|^2\bigr)^{1/2}
 =\|f\|_\infty.
\]
Cualquier norma $C^*$ compatible con esas operaciones debe,
por tanto, coincidir con la norma uniforme.
\end{cthejemplo}

\begin{cthteorema}{($*$-homomorfismos unitarios son contractivos)}
Sean $\A,\B$ $C^*$-álgebras unitales y $\Phi:\A\to\B$ un $*$-homomorfismo unital. Entonces
\[
 \|\Phi(a)\|\le\|a\|\quad(a\in\A),
 \qquad \|\Phi\|=1.
\]
\end{cthteorema}
\origen{Teorema 2.8}

\begin{proof}[Demostración]
Para cada $b\in\A$, el homomorfismo unital satisface
$\spec(\Phi(b))\subset\spec(b)$ y por tanto
$r(\Phi(b))\le r(b)$. Tomando $b=a^*a$,
\[
\|\Phi(a)\|^2
=\|\Phi(a)^*\Phi(a)\|
=\|\Phi(a^*a)\|
=r(\Phi(a^*a))
\le r(a^*a)=\|a\|^2.
\]
Como $\Phi(1_\A)=1_\B$ y ambas unidades tienen norma uno, se tiene
$\|\Phi\|\ge1$; la contractividad da la igualdad.
\end{proof}

\proofdiagram
{$\Phi$ unital y $*$-homomorfismo}
{Inclusión espectral de homomorfismos; teorema normal $\|x\|=r(x)$}
{Aplicar a $a^*a$: $\Phi(a^*a)=\Phi(a)^*\Phi(a)$}
{$\|\Phi(a)\|^2\le\|a\|^2$ y $\|\Phi\|=1$}

\begin{cthejemplo}{(Restricción como $*$-homomorfismo contractivo)}
Sean $X$ compacto Hausdorff e $Y\subset X$ cerrado y no vacío.
Queremos comparar una función definida en $X$ con su restricción a
$Y$. El mapa $\Phi:C(X)\to C(Y)$,
$\Phi(f)=f|_Y$, preserva sumas, productos, conjugación puntual y la
función constante $1$, así que es un $*$-homomorfismo unital.
El teorema de contractividad implica
$\|f|_Y\|_\infty\le\|f\|_\infty$, desigualdad que también se ve al
tomar el supremo sobre un conjunto menor. Restringir el dominio no
puede aumentar la norma uniforme.
\end{cthejemplo}

\begin{cthobservacion}{(Parte real e imaginaria)}
Sea $\A$ una $C^*$-álgebra. Todo $a\in\A$ se escribe de manera única
\[
 a=x+iy,
 \qquad
 x=\frac{a+a^*}{2},\quad y=\frac{a-a^*}{2i},
\]
con $x=x^*$ y $y=y^*$.
\end{cthobservacion}
\origen{Observación 2.9}

\begin{proof}[Demostración]
Las fórmulas dadas cumplen $x^*=x$, $y^*=y$ y $x+iy=a$. Si también
$a=u+iv$ con $u=u^*$ y $v=v^*$, entonces $a^*=u-iv$; al sumar y restar
estas dos identidades se obtiene $u=(a+a^*)/2=x$ y
$v=(a-a^*)/(2i)=y$. Esto prueba la unicidad.
\end{proof}

\begin{cthejemplo}{(Partes autoadjuntas de una matriz no hermitiana)}
Para ver la descomposición de la observación en un caso no trivial,
tomemos
$A=\begin{pmatrix}1&i\\0&2\end{pmatrix}$.
El cálculo de las partes real e imaginaria da
\[
 X=\frac{A+A^*}{2}
 =\begin{pmatrix}1&i/2\\-i/2&2\end{pmatrix},
 \qquad
 Y=\frac{A-A^*}{2i}
 =\begin{pmatrix}0&1/2\\1/2&0\end{pmatrix}.
\]
Se verifica $X^*=X$, $Y^*=Y$ y $A=X+iY$.
Así incluso una matriz que no es autoadjunta se reconstruye a
partir de dos matrices que sí lo son.
\end{cthejemplo}

\begin{cthlema}{(Estimación vertical para autoadjuntos)}
Sea $\A$ una $C^*$-álgebra unital, $x\in\A$ autoadjunto y $t\in\mathbb R$. Entonces
\[
 \|x+it1_\A\|^2\le \|x\|^2+t^2.
\]
\end{cthlema}
\origen{Lema 2.10}

\begin{proof}[Demostración]
\[
\|x+it1_\A\|^2
=\|(x-it1_\A)(x+it1_\A)\|
=\|x^2+t^2 1_\A\|
\le\|x^2\|+t^2
=\|x\|^2+t^2.
\]
\end{proof}

\proofdiagram
{$x\in\A$ autoadjunto; $t\in\mathbb R$}
{Identidad $C^*$}
{$\|x+it1_\A\|^2=\|x^2+t^2 1_\A\|$}
{$\le\|x\|^2+t^2$}

\begin{cthejemplo}{(La cota vertical es exacta para escalares reales)}
Comprobemos si la desigualdad del lema puede mejorarse en general.
En la $C^*$-álgebra $\mathbb C$, tomamos el autoadjunto
$s\in\mathbb R$ y $t\in\mathbb R$. Entonces
$|s+it|^2=(s-it)(s+it)=s^2+t^2$ y $\|s\|^2=s^2$.
La desigualdad del lema es igualdad en este caso; el término
$t^2$ y el coeficiente de $\|s\|^2$ son, por ello, naturales.
\end{cthejemplo}

\begin{cthproposicion}{(Una cota vertical fuerza realidad)}
Si $z\in\mathbb C$ y existe $C\ge0$ tal que
\[
 |z+it|^2\le C+t^2\qquad(t\in\mathbb R),
\]
entonces $z\in\mathbb R$.
\end{cthproposicion}
\origen{Proposición 2.11}

\begin{proof}[Demostración]
Si $z=x+iy$ con $x,y\in\mathbb R$, entonces
$|z+it|^2-t^2=x^2+y^2+2yt$. Si $y>0$, el lado derecho tiende a $+\infty$ cuando
$t\to+\infty$; si $y<0$, ocurre cuando $t\to-\infty$. La existencia de una cota uniforme obliga a
$y=0$.
\end{proof}

\proofdiagram
{$z\in\mathbb C$; $C\ge0$; escribir $z=x+iy$}
{Expandir $|z+it|^2-t^2=x^2+y^2+2yt$}
{Si $y\neq0$, elegir $t\to\pm\infty$ con el signo de $y$}
{La cota falla; luego $y=0$}

\begin{cthejemplo}{(Una cota que falla para $1+i$)}
Para ver cómo la Proposición~2.11 detecta una parte imaginaria, probemos su hipótesis con
$z=1+i$. El cálculo da
\[
 |z+it|^2-t^2=|1+i(1+t)|^2-t^2=2+2t.
\]
Al hacer $t\to+\infty$, esta expresión supera cualquier constante $C$. Por tanto no existe
una cota $|1+i+it|^2\le C+t^2$ válida para todo $t\in\mathbb R$.
El término lineal en $t$ es la señal concreta de que $z$ no es real.
\end{cthejemplo}

\begin{cthproposicion}{(Funcionales que alcanzan su norma en la unidad son reales en autoadjuntos)}
Sea $\A$ una $C^*$-álgebra unital y $\varphi:\A\to\mathbb C$ un funcional
lineal continuo con $\|\varphi\|=\varphi(1_\A)$. Entonces
\[
 \varphi(x)\in\mathbb R\qquad\text{para todo }x\in\A\text{ con }x=x^*.
\]
\end{cthproposicion}
\origen{Proposición 2.12}

\begin{proof}[Demostración]
Si $\varphi=0$, entonces $\varphi(x)=0\in\mathbb R$ para todo autoadjunto $x$.
En otro caso, $\varphi(1_\A)=\|\varphi\|>0$.
Definimos $\psi=\varphi/\varphi(1_\A)$; entonces $\|\psi\|=\psi(1_\A)=1$.
Para $t\in\mathbb R$,
\[
 |\psi(x)+it|^2
 =|\psi(x+it1_\A)|^2
 \le\|x+it1_\A\|^2
 \le\|x\|^2+t^2.
\]
La proposición anterior, con $z=\psi(x)$, implica que $\psi(x)$ es real; como
$\varphi(1_\A)>0$, también $\varphi(x)$ lo es.
\end{proof}

\proofdiagram
{$\varphi:\A\to\mathbb C$ continuo; $\|\varphi\|=\varphi(1_\A)$; $x=x^*$}
{Lema de estimación vertical + proposición escalar anterior}
{Si $\varphi\ne0$: $\psi=\varphi/\varphi(1_\A)$ y $|\psi(x)+it|^2\le\|x+it1_\A\|^2$}
{$\varphi=0$ o $\psi(x)\in\mathbb R$; luego $\varphi(x)\in\mathbb R$}

\begin{cthejemplo}{(Una evaluación cumple la hipótesis de realidad)}
Tomemos $X=[0,1]$ y $x_0\in X$. La evaluación
$\varphi=\operatorname{ev}_{x_0}:C(X)\to\mathbb C$ satisface
$\|\varphi\|=\varphi(1)=1$, como se comprobó en la Sección~1.
Para la función autoadjunta $h(t)=t^2$, se obtiene
$\varphi(h)=x_0^2\in\mathbb R$.
La Proposición~2.12 asegura la misma conclusión para \emph{todo} autoadjunto
de $C(X)$ y, más generalmente, para funcionales con esa condición de norma en
cualquier $C^*$-álgebra unital.
\end{cthejemplo}

\begin{cthproposicion}{(El espectro de un autoadjunto es real)}
Sea $\A$ una $C^*$-álgebra unital y $a\in\A$ autoadjunto. Entonces
\[
 \spec(a)\subset\mathbb R.
\]
\end{cthproposicion}
\origen{Proposición 2.13}

\begin{proof}[Demostración]
Por la Tarea 1.9 existe una subálgebra de Banach conmutativa, involutiva y unital
$\mathcal C\subset\A$ que contiene a $a$ y preserva su espectro. En ella,
\[
 \spec_\A(a)=\spec_{\mathcal C}(a)=\{\varphi(a):\varphi\in\Omega(\mathcal C)\}.
\]
Como $\mathcal C$ es cerrada y hereda la identidad $C^*$, también es una
$C^*$-álgebra. Cada carácter de $\mathcal C$ tiene norma uno y
$\varphi(1_{\mathcal C})=1$, por lo que
la proposición anterior da
$\varphi(a)\in\mathbb R$. La descripción espectral escrita arriba muestra
entonces que cada $\lambda\in\spec_\A(a)$ es real.
\end{proof}

\proofdiagrambranch
{$\A$ es $C^*$ unital; $a\in\A$, $a=a^*$}
{Reducir a una subálgebra conmutativa que preserve el espectro (Tarea 1.9)}
{Espectro vía caracteres + funcionales reales en autoadjuntos}
{$\lambda\in\spec(a)\Rightarrow\lambda=\varphi(a)\in\mathbb R$}
{$\spec(a)\subset\mathbb R$}

\begin{cthejemplo}{(Valores propios reales de una matriz hermitiana)}
Consideremos $H=\begin{pmatrix}2&i\\-i&3\end{pmatrix}\in M_2(\mathbb C)$.
La transposición conjugada da $H^*=H$, así que la Proposición~2.13 predice
$\spec(H)\subset\mathbb R$. En efecto,
\[
 \det(\lambda I-H)=(\lambda-2)(\lambda-3)-1
 =\lambda^2-5\lambda+5,
 \qquad
 \spec(H)=\left\{\frac{5-\sqrt5}{2},\frac{5+\sqrt5}{2}\right\}.
\]
El cálculo finito confirma la conclusión espectral que la proposición extiende
a elementos autoadjuntos sin una matriz asociada.
\end{cthejemplo}

\begin{cthnota}{(Dónde intervienen las hipótesis)}
La igualdad $a=a^*$ permite aplicar la Tarea 1.9 en su caso involutivo y después
la proposición sobre funcionales reales. La unidad interviene en la descripción
del espectro mediante caracteres normalizados.
\end{cthnota}


\begin{cthdefinicion}{(Elemento positivo y orden)}
Sea $\A$ una $C^*$-álgebra unital. Un elemento $x\in\A$ es \emph{positivo}, y escribimos
$x\ge0$, si $x=x^*$ y $\spec_\A(x)\subset[0,\infty)$. Para $x,y\in\A$ autoadjuntos,
$x\le y$ significa que $y-x\ge0$.
\end{cthdefinicion}

\begin{puente}
Éste es uno de los cruces más claros entre las tareas y la Sección 2: la demostración original de la
realidad del espectro usa explícitamente el ejercicio de la Tarea 1 que construye una subálgebra
conmutativa que contiene al elemento y preserva su espectro. El ejercicio no es accesorio: es una
pieza de la prueba.
\end{puente}

\section{Caracteres, permanencia espectral y representación}

\begin{cthcorolario}{(Funcionales de norma en la unidad preservan la involución)}
Sea $\A$ una $C^*$-álgebra unital y $\varphi:\A\to\mathbb C$ un funcional lineal
continuo con $\|\varphi\|=\varphi(1_\A)$. Entonces, para todo $a\in\A$,
\[
 \varphi(a^*)=\overline{\varphi(a)}.
\]
En particular, esto vale para los caracteres de $\A$.
\end{cthcorolario}
\origen{Corolario 2.14}

\begin{proof}[Demostración]
Sea $a\in\A$ y escribamos $a=x+iy$, donde
$x=(a+a^*)/2$ y $y=(a-a^*)/(2i)$ son autoadjuntos. Por la proposición
anterior, $\varphi(x),\varphi(y)\in\mathbb R$, y
\[
 \varphi(a^*)=\varphi(x-iy)=\varphi(x)-i\varphi(y)
 =\overline{\varphi(x)+i\varphi(y)}=\overline{\varphi(a)}.
\]
Si $\varphi$ es un carácter, entonces es un funcional continuo con
$\varphi(1_\A)=\|\varphi\|=1$, de modo que cumple la hipótesis usada.
\end{proof}

\proofdiagram
{$\varphi:\A\to\mathbb C$ continuo; $\|\varphi\|=\varphi(1_\A)$; $a=x+iy$ con $x,y$ autoadjuntos}
{Proposición 2.12: $\varphi(x),\varphi(y)\in\mathbb R$}
{$\varphi(a^*)=\varphi(x)-i\varphi(y)$}
{$\varphi(a^*)=\overline{\varphi(a)}$}

\begin{cthejemplo}{(Una evaluación preserva la conjugación)}[unbreakable]
En $C([0,1])$ elijamos $f(t)=1+it$ y $x\in[0,1]$.
La involución es conjugación puntual, por lo que
$f^*(t)=1-it$. Para el carácter $\operatorname{ev}_x$,
\[
 \operatorname{ev}_x(f^*)=1-ix
 =\overline{1+ix}
 =\overline{\operatorname{ev}_x(f)}.
\]
Esta igualdad ilustra el Corolario~2.14 en un carácter concreto; el corolario
también la garantiza para funcionales de norma alcanzada en la unidad.
\end{cthejemplo}

\begin{cthcorolario}{(La transformada de Gelfand es un $*$-homomorfismo)}
Si $\A$ es una $C^*$-álgebra conmutativa unital, entonces
\[
 \Gamma:\A\to C(\Omega(\A)),\qquad a\mapsto\widehat a,
\]
es un $*$-homomorfismo unital.
\end{cthcorolario}
\origen{Corolario 2.15}

\begin{proof}[Demostración]
Para $a,b\in\A$ y $\varphi\in\Omega(\A)$,
$\widehat{a+b}(\varphi)=\widehat a(\varphi)+\widehat b(\varphi)$,
$\widehat{ab}(\varphi)=\widehat a(\varphi)\widehat b(\varphi)$ y
$\widehat{1_\A}(\varphi)=1$, por linealidad, multiplicatividad y unidad del
carácter. También $\widehat{\lambda a}(\varphi)=\lambda\widehat a(\varphi)$
para $\lambda\in\mathbb C$. Por tanto la transformada de Gelfand es un
homomorfismo lineal, multiplicativo y unital. Cada
$\varphi\in\Omega(\A)$ satisface $\|\varphi\|=\varphi(1_\A)=1$, así que el corolario
anterior da, para $a\in\A$,
\[
 \widehat{a^*}(\varphi)=\varphi(a^*)
 =\overline{\varphi(a)}=(\widehat a)^*(\varphi).
\]
Como la igualdad vale para todo $\varphi\in\Omega(\A)$, la transformada preserva $*$.
\end{proof}

\proofdiagram
{Gelfand es homomorfismo unital}
{Los caracteres preservan $*$}
{$\widehat{a^*}(\varphi)=\varphi(a^*)=\overline{\varphi(a)}$}
{$\widehat{a^*}=(\widehat a)^*$}

\begin{cthejemplo}{(Conjugación bajo la transformada de Gelfand)}
Tomemos $\A=C([0,1])$ y $f(t)=1+it$. Como los caracteres son las evaluaciones,
$\widehat f(\operatorname{ev}_x)=1+ix$ para cada $x\in[0,1]$. Asimismo,
\[
 \widehat{f^*}(\operatorname{ev}_x)=1-ix
 =\overline{\widehat f(\operatorname{ev}_x)}.
\]
El Corolario~2.15 expresa esta identidad simultáneamente para todos los
elementos y caracteres: la transformada de Gelfand respeta la involución.
\end{cthejemplo}

\begin{cthproposicion}{(Invertibilidad detectada por $a^*a$ y $aa^*$)}
Para $a$ en una $C^*$-álgebra unital $\A$,
\[
 a\in\GL(\A)
 \quad\Longleftrightarrow\quad
 a^*a\in\GL(\A)\text{ y }aa^*\in\GL(\A).
\]
\end{cthproposicion}
\origen{Proposición 2.16}

\begin{proof}[Demostración]
Si $a\in\GL(\A)$, la identidad
$(a^{-1})^*a^*=a^*(a^{-1})^*=1_\A$ muestra que $a^*$ es invertible.
Por tanto, sus productos $a^*a$ y
$aa^*$ son invertibles. Para la recíproca, tomamos
$y=(a^*a)^{-1}$ y $z=(aa^*)^{-1}$, de modo que
$y(a^*a)=1_\A$ y $(aa^*)z=1_\A$. Entonces $ya^*$ es inverso izquierdo de $a$ y $a^*z$ es inverso
derecho. Además,
\[
 ya^*=ya^*(aa^*z)=(ya^*a)a^*z=a^*z.
\]
Por tanto existe un inverso bilateral.
\end{proof}

\proofdiagrambranch
{$a\in\A$; $a^*a\in\GL(\A)$}
{$aa^*\in\GL(\A)$}
{$y=(a^*a)^{-1}$ y $z=(aa^*)^{-1}$}
{$ya^*$ y $a^*z$ coinciden por asociatividad}
{$a\in\GL(\A)$}

\begin{cthejemplo}{(Inversos de una matriz detectados por sus productos)}
En $M_2(\mathbb C)$, sea
$A=\begin{pmatrix}1&1\\0&1\end{pmatrix}$. Tenemos
\[
 A^*A=\begin{pmatrix}1&1\\1&2\end{pmatrix},
 \qquad
 AA^*=\begin{pmatrix}2&1\\1&1\end{pmatrix},
\]
y ambos determinantes son $1$. Por la Proposición~2.16, $A$ es invertible;
directamente, $A^{-1}=\begin{pmatrix}1&-1\\0&1\end{pmatrix}$.
La comprobación en matrices cuadradas muestra cómo los dos productos
proporcionan inversos izquierdo y derecho en el argumento abstracto.
\end{cthejemplo}

\begin{cthproposicion}{(Permanencia espectral en $C^*$-subálgebras)}
Sea $\A$ una $C^*$-álgebra unital y $\B\subset\A$ una $C^*$-subálgebra unital
con $1_\B=1_\A$. Para todo $b\in\B$,
\[
 b\in\GL(\B)\Longleftrightarrow b\in\GL(\A),
 \qquad
 \spec_\B(b)=\spec_\A(b).
\]
\end{cthproposicion}
\origen{Proposición 2.17}

\begin{proof}[Demostración]
Si $b\in\GL(\B)$, su inverso $b^{-1}\in\B\subset\A$ satisface
$bb^{-1}=b^{-1}b=1_\B=1_\A$; así $b\in\GL(\A)$.
Para la otra dirección, sea primero $h\in\B$ autoadjunto e invertible en $\A$.
Como $\spec_\B(h)\subset\mathbb R$, para cada $n\ge1$ el elemento
$h+i n^{-1}1_\B$ es invertible en $\B$. Estos elementos convergen a $h$;
la inversión es continua en $\GL(\A)$, por lo que sus inversos convergen en $\A$
a $h^{-1}$. Como $\B$ es cerrada, $h^{-1}\in\B$.
Para $b\in\B\cap\GL(\A)$ general, aplicamos este argumento a $b^*b$ y $bb^*$,
invertibles en $\A$ cuando $b$ lo es, y obtenemos que ambos son invertibles en $\B$.
La proposición anterior concluye que $b$ es invertible en $\B$.
Finalmente, para cada $\lambda\in\mathbb C$ aplicamos la equivalencia de
invertibilidad a $b-\lambda1_\B=b-\lambda1_\A$; resulta
$\lambda\notin\spec_\B(b)$ si y sólo si
$\lambda\notin\spec_\A(b)$, es decir,
$\spec_\B(b)=\spec_\A(b)$.
\end{proof}

\proofdiagrambranch
{$b\in\B\cap\GL(\A)$}
{$h=b^*b$ o $bb^*$: autoadjunto e invertible en $\A$}
{$\spec_\B(h)\subset\mathbb R$; $h+i n^{-1}1_\B\in\GL(\B)$}
{Inversión continua en $\A$ y clausura de $\B$: $h^{-1}\in\B$}
{$b^*b,bb^*\in\GL(\B)$; luego $b\in\GL(\B)$}

\begin{cthejemplo}{(Contraejemplo: el álgebra del disco no preserva espectros)}
La Tarea~1.12 compara el álgebra del disco $A(\mathbb D)$ con $C(\mathbb T)$.
Tomemos la función coordenada $u(z)=z$ sobre $\mathbb T$. Como
$|u(z)|=1$ en la circunferencia, $u$ es invertible en $C(\mathbb T)$,
con inversa $\overline u$. Sin embargo, $\overline u$ no pertenece a
$A(\mathbb D)$: una inversa holomorfa de $z$ en el disco tendría un polo
en $0$. Por ello $0\in\spec_{A(\mathbb D)}(u)$ mientras
$0\notin\spec_{C(\mathbb T)}(u)$.
Esta subálgebra cerrada no es estable por la involución puntual y no cumple
la hipótesis $C^*$ de la Proposición~2.17.
\end{cthejemplo}

\begin{cthnota}{(Dónde intervienen las hipótesis)}
La estabilidad de $\B$ por involución permite usar $b^*b$ y $bb^*$ dentro de ella;
la unidad común permite comparar la invertibilidad, y que $\B$ sea cerrada
permite retener el límite de los inversos.
\end{cthnota}


\begin{cthproposicion}{(El espectro de un normal es el espacio de caracteres generado por él)}
Sea $\A$ una $C^*$-álgebra unital y $a\in\A$ normal. Denotemos por
$C^*(a,1_\A)$ la menor
$C^*$-subálgebra unital de $\A$ que contiene a $a$. Entonces
\[
 \Omega(C^*(a,1_\A))\longrightarrow\spec_\A(a),\qquad
 \varphi\longmapsto\varphi(a),
\]
es un homeomorfismo.
\end{cthproposicion}
\origen{Proposición 2.18}

\begin{proof}[Demostración]
La $C^*$-álgebra $C^*(a,1_\A)$ es conmutativa porque $a$ conmuta con $a^*$.
Por permanencia espectral,
el espectro de $a$ dentro de esta subálgebra es el mismo que en $\A$.
En una álgebra de Banach conmutativa unital, todo punto del espectro de $a$
es $\varphi(a)$ para algún carácter $\varphi$; de aquí se obtiene la
sobreyectividad. La aplicación es continua porque la topología de
$\Omega(C^*(a,1_\A))$ es la de convergencia puntual en los elementos del
álgebra. Si $\varphi,\psi\in\Omega(C^*(a,1_\A))$ y
$\varphi(a)=\psi(a)$, entonces también
$\varphi(a^*)=\overline{\varphi(a)}=\overline{\psi(a)}=\psi(a^*)$. Como los polinomios en
$a$ y $a^*$ son densos en $C^*(a,1_\A)$, la linealidad y multiplicatividad dan
igualdad en esos polinomios, y la continuidad de los caracteres extiende la
igualdad a toda la subálgebra. Así $\varphi=\psi$. El espacio maximal y el
espectro son compactos Hausdorff; una biyección continua entre ellos tiene
inversa continua y es un homeomorfismo.
\end{proof}

\proofdiagram
{$a$ normal}
{$C^*(a,1_\A)$ es conmutativa; permanencia espectral; caracteres}
{Sobreyectiva por espectro; inyectiva porque $a,a^*$ generan densamente}
{Biyectiva continua entre compactos Hausdorff: homeomorfismo}

\begin{cthejemplo}{(Interpolación en el espectro de una matriz diagonal)}
Sean $\lambda_1,\ldots,\lambda_n$ valores distintos y
$A=\operatorname{diag}(\lambda_1,\ldots,\lambda_n)$.
La Proposición~2.18 identifica los caracteres de $C^*(A,I_n)$ con esos
valores propios. En concreto, para números arbitrarios $c_1,\ldots,c_n$,
la interpolación de Lagrange produce un polinomio $p$ tal que
$p(\lambda_j)=c_j$; entonces
\[
 p(A)=\operatorname{diag}(c_1,\ldots,c_n).
\]
Así $C^*(A,I_n)$ es toda el álgebra diagonal y se identifica con
$C(\{\lambda_1,\ldots,\lambda_n\})$. Cada carácter selecciona una coordenada.
\end{cthejemplo}

\begin{cthteorema}{(Stone--Weierstrass complejo)}
Sea $X$ compacto Hausdorff y $\mathcal D\subset C(X)$ una $*$-subálgebra unital que separa puntos.
Entonces $\mathcal D$ es densa en $C(X)$ para la norma uniforme.
\end{cthteorema}
\origen{Teorema 2.19. Se usa como herramienta externa; no se reproduce aquí su demostración completa.}

\begin{cthejemplo}{(Polinomios de coordenada y conjugada en un compacto)}
Sea $K\subset\mathbb C$ compacto. Queremos aproximar funciones de $C(K)$
mediante expresiones algebraicas de la coordenada $z$. La subálgebra unital
de polinomios en $z$ y $\overline z$ separa puntos porque la función
$z\mapsto z$ ya los distingue; además es estable por conjugación.
El Teorema~2.19 de Stone--Weierstrass afirma entonces que es densa en
$C(K)$ para la norma uniforme.
Sobre $K=\mathbb T$ esto incluye los polinomios trigonométricos, pues
$\overline z=z^{-1}$. La Tarea~2.8 utiliza precisamente esta densidad
para el cálculo funcional.
\end{cthejemplo}

\begin{cthteorema}{(Gelfand--Naimark conmutativo)}
Sea $\A$ una $C^*$-álgebra conmutativa unital. La transformada de Gelfand
\[
 \Gamma:\A\longrightarrow C(\Omega(\A))
\]
es un $*$-isomorfismo isométrico.
\end{cthteorema}
\origen{Teorema 2.20}

\begin{proof}[Demostración]
Para cada $a\in\A$, la conmutatividad hace normal a $a$. La descripción del
espectro por caracteres, junto con el teorema de la norma de un normal, da
\[
 \|\widehat a\|_\infty=r(a)=\|a\|.
\]
Así $\Gamma$ es isométrica e inyectiva. Además su imagen es cerrada: una
sucesión convergente de imágenes proviene, por la isometría, de una sucesión
de Cauchy en $\A$, que converge por completitud. La imagen es una
$*$-subálgebra unital de $C(\Omega(\A))$ y separa puntos: si $\varphi\neq\psi$, existe $a$ con
$\varphi(a)\neq\psi(a)$. Stone--Weierstrass dice que la imagen es densa. Cerrada y densa implica
que es todo $C(\Omega(\A))$.
\end{proof}

\proofdiagrambranch
{$\A$ conmutativa $C^*$ unital; $a\in\A$}
{Gelfand es $*$-homomorfismo; $\|a\|=r(a)$; Stone--Weierstrass}
{Isometría: $\|\widehat a\|_\infty=r(a)=\|a\|$}
{La imagen es cerrada, unital, $*$-estable y separa puntos; luego es densa}
{$\Gamma$ es sobreyectiva e isométrica}

\begin{cthejemplo}{(Matrices diagonales como funciones en dos puntos)}
Sea $\A=\{\operatorname{diag}(\alpha,\beta):\alpha,\beta\in\mathbb C\}$.
Sus dos caracteres son las evaluaciones de las entradas
$\varphi_1(\operatorname{diag}(\alpha,\beta))=\alpha$ y
$\varphi_2(\operatorname{diag}(\alpha,\beta))=\beta$.
La transformada de Gelfand del Teorema~2.20 envía
$\operatorname{diag}(\alpha,\beta)$ a la función
$f:\{\varphi_1,\varphi_2\}\to\mathbb C$ con valores $\alpha,\beta$.
La igualdad $\|\operatorname{diag}(\alpha,\beta)\|=\max\{|\alpha|,|\beta|\}
=\|f\|_\infty$ permite ver la isometría del teorema en este modelo finito.
\end{cthejemplo}

\begin{cthcorolario}{(Clasificación de $C^*$-álgebras conmutativas unitales)}
Sea $\A$ una $C^*$-álgebra conmutativa unital. Existe un compacto Hausdorff $X$ tal que
$\A$ es isométricamente $*$-isomorfa a $C(X)$; ese espacio $X$ es único salvo homeomorfismo.
\end{cthcorolario}
\origen{Corolario 2.21}

\begin{proof}[Demostración]
La existencia se obtiene con $X=\Omega(\A)$ por Gelfand--Naimark.
Para la unicidad, sea $Y$ compacto Hausdorff y $\varphi\in\Omega(C(Y))$.
Si las funciones de $\ker\varphi$ no tuvieran un cero común, por compacidad
los abiertos $\{y\in Y:f(y)\neq0\}$, con $f\in\ker\varphi$, tendrían una
subcubierta finita. Existirían entonces $f_1,\ldots,f_m\in\ker\varphi$
sin cero común. El corolario sobre la
involución implica que $f_j^*\in\ker\varphi$; por tanto
$g=\sum_{j=1}^m f_j^*f_j\in\ker\varphi$. Pero $g(y)>0$ para todo $y\in Y$,
de modo que $g$ es invertible en $C(Y)$ con inversa $1/g$. Esto contradice
$1=\varphi(1)=\varphi(g)\varphi(g^{-1})=0$.
Existe, pues, $y_0\in Y$ en el que se anula todo $f\in\ker\varphi$.
Para $h\in C(Y)$, la función $h-\varphi(h)1$ pertenece al núcleo, y entonces
$h(y_0)=\varphi(h)$. Así $\varphi=\operatorname{ev}_{y_0}$.
La aplicación $y\mapsto\operatorname{ev}_y$ es continua porque, para cada
$h\in C(Y)$, su coordenada $y\mapsto\operatorname{ev}_y(h)=h(y)$ lo es.
Es inyectiva porque $C(Y)$ separa puntos y sobreyectiva por lo anterior.
Como $Y$ es compacto y $\Omega(C(Y))$ es Hausdorff, es un homeomorfismo.
Si $F:\A\to C(Y)$ es un $*$-isomorfismo, la composición
$\chi\mapsto\chi\circ F$ es un homeomorfismo de $\Omega(C(Y))$ a
$\Omega(\A)$, con inversa inducida por $F^{-1}$. Por tanto
$Y\cong\Omega(\A)=X$. La Tarea 2.5 desarrolla por separado esta identificación.
\end{proof}

\proofdiagram
{$\A$ conmutativa $C^*$ unital; Gelfand--Naimark}
{$X=\Omega(\A)$}
{$\A\cong C(\Omega(\A))$; los caracteres de $C(X)$ son evaluaciones}
{La representación existe y $X$ es único salvo homeomorfismo}

\begin{cthejemplo}{(Reconstruir el intervalo a partir de su álgebra)}
Consideremos $\A=C([0,1])$. La prueba del Corolario~2.21 muestra que cada
carácter de $\A$ es una evaluación $\operatorname{ev}_x$ en un único
$x\in[0,1]$. Así, la aplicación
$x\mapsto\operatorname{ev}_x$ identifica topológicamente
$[0,1]$ con $\Omega(\A)$.
Si además $\A$ fuese $*$-isomorfa a $C(Y)$ para un compacto Hausdorff $Y$,
el mismo corolario daría $Y\cong\Omega(\A)\cong[0,1]$.
El álgebra de funciones recupera, por tanto, el espacio del que proviene.
\end{cthejemplo}

\begin{investigacion}[unbreakable]
\textbf{Uso de construcción.} En \cite[Sec.~2, pp.~465--466]{Karlovich2017},
$\SO(\mathbb R_+)$ es una $C^*$-subálgebra conmutativa unital de
$L^\infty(\mathbb R_+)$: sus funciones son continuas y acotadas, y su oscilación
en $[r,2r]$ tiende a cero cuando $r\to0$ o $r\to\infty$.
Identificando $C(\overline{\mathbb R}_+)$ con sus
restricciones a $\mathbb R_+=(0,\infty)$, donde
$\overline{\mathbb R}_+=[0,\infty]$ es la compactificación, las fibras sobre
$s\in\{0,\infty\}$ son
\[
 M_s(\SO(\mathbb R_+))=
 \{\xi\in\Omega(\SO(\mathbb R_+)):
 \xi|_{C(\overline{\mathbb R}_+)}=\operatorname{ev}_s\}.
\]
La Proposición~2.2 de esa fuente relaciona la evaluación de un carácter en una fibra
con límites parciales de funciones lentamente oscilantes.
\textbf{Lo que conviene reconocer:} los caracteres y el espacio maximal estudiados
en el curso sirven aquí para codificar información en los extremos.
\end{investigacion}

\section{Cálculo funcional continuo}

\begin{cthteorema}{(Cálculo funcional continuo para un elemento normal)}
Sea $\A$ una $C^*$-álgebra unital y $a\in\A$ normal. Existe un único
$*$-homomorfismo isométrico
\[
 C(\spec_\A(a))\longrightarrow\A,
 \qquad f\longmapsto f(a),
\]
que envía la función identidad $z\mapsto z$ a $a$ y la función constante $1$ a $1_\A$.
Su imagen es $C^*(a,1_\A)$.
\end{cthteorema}
\origen{Teorema 2.22}

\begin{proof}[Demostración]
Sea $\B=C^*(a,1_\A)$. Por Gelfand--Naimark,
$\Gamma_\B:\B\to C(\Omega(\B))$ es un $*$-isomorfismo isométrico. La proposición anterior da el
homeomorfismo
$\Psi:\Omega(\B)\to\spec(a)$, $\Psi(\varphi)=\varphi(a)$. La composición
\[
 C(\spec(a))\xrightarrow{f\mapsto f\circ\Psi} C(\Omega(\B))
 \xrightarrow{\Gamma_\B^{-1}}\B\subset\A
\]
es un $*$-isomorfismo isométrico de $C(\spec(a))$ sobre $\B$.
Para la función identidad $\mathrm{id}(z)=z$, se tiene
$(\mathrm{id}\circ\Psi)(\varphi)=\varphi(a)=\widehat a(\varphi)$; por
tanto su imagen es $a$. La constante $1$ corresponde a
$\widehat{1_\A}$, así que su imagen es $1_\A$.
Si $F:C(\spec(a))\to\A$ es otro $*$-homomorfismo isométrico con esas
propiedades, entonces $F(\overline{\mathrm{id}})=a^*$ y coincide con la
construcción anterior en todos los polinomios en $z$ y $\overline z$.
Stone--Weierstrass hace densa esa $*$-subálgebra unital en
$C(\spec(a))$; la continuidad de ambos mapas extiende la igualdad a toda
función. Así se obtienen la unicidad y la imagen indicadas.
\end{proof}

\proofdiagrambranch
{$\A$ es $C^*$ unital; $a\in\A$ normal}
{$C^*(a,1_\A)$ conmutativa + Gelfand--Naimark}
{Identificar $\Omega(C^*(a,1_\A))$ con $\spec_\A(a)$}
{Transportar funciones por el homeomorfismo y por la transformada de Gelfand}
{Definir de forma única $f(a)$}

\begin{cthejemplo}{(Raíz espectral de una matriz diagonal)}
Tomemos la matriz normal positiva
$A=\operatorname{diag}(1,4)$ y la función
$f(t)=\sqrt t$ en $\spec(A)=\{1,4\}$. El Teorema~2.22 da
\[
 f(A)=\operatorname{diag}(f(1),f(4))
     =\operatorname{diag}(1,2),
 \qquad f(A)^2=A.
\]
En general, si
$A=U\operatorname{diag}(\lambda_1,\ldots,\lambda_n)U^*$ es normal,
$f(A)=U\operatorname{diag}(f(\lambda_1),\ldots,f(\lambda_n))U^*$.
El ejemplo muestra que el cálculo funcional abstracto conserva la receta
espectral familiar de matrices.
\end{cthejemplo}

\begin{cthnota}{(Dependencia con las tareas)}
La normalidad hace conmutativa a $C^*(a,1_\A)$; la permanencia espectral identifica
su espectro con el de $a$ en $\A$. La Tarea 2.4 combina este cálculo con el espectro
de funciones continuas de la Tarea 1.10 para obtener
$\spec(f(a))=f(\spec(a))$.
\end{cthnota}


\begin{cthdefinicion}{(Cálculo funcional por restricción desde un conjunto ambiente)}
Sea $\A$ una $C^*$-álgebra unital, $a\in\A$ normal y $E\subset\mathbb C$
localmente compacto con $\spec_\A(a)\subset E$. Para $f\in C_b(E)$ se define
\[
 f(a):=(f|_{\spec_\A(a)})(a).
\]
Así se obtiene un $*$-homomorfismo unital $C_b(E)\to\A$. La misma fórmula
define $f(a)$ si $f:E\to\mathbb C$ es continua y no acotada, pues su restricción
al compacto $\spec_\A(a)$ sí es acotada. Sólo importa dicha restricción.
\end{cthdefinicion}
\origen{Definición 2.23}

\begin{cthobservacion}{(Polinomios en $a$ y $a^*$)}
Sea $a$ normal en una $C^*$-álgebra unital $\A$. Si
$f(z)=z^2-\overline z^{\,3}+1$ para $z\in\mathbb C$, entonces
\[
 f(a)=a^2-(a^*)^3+1_\A.
\]
Si $f(z)=|z|^2=\overline z z$, entonces $f(a)=a^*a$.
\end{cthobservacion}
\origen{Observación 2.24}

\begin{cthobservacion}{(En $C(X)$, el cálculo funcional es composición)}
Si $X$ es compacto Hausdorff, $g\in C(X)$ y $f\in C(g(X))$, entonces
\[
 f(g)=f\circ g.
\]
\end{cthobservacion}
\origen{Observación 2.25}

\begin{proof}[Demostración]
Por la identificación $\Omega(C(X))\cong X$, la transformada de $g$ toma en
$x\in X$ el valor $g(x)$, y $\spec(g)=g(X)$. La construcción del cálculo
funcional transporta $f$ por esa transformada; en cada $x$ se obtiene
$f(g)(x)=f(g(x))$. Dos funciones de $C(X)$ que coinciden en cada punto son
iguales, de modo que $f(g)=f\circ g$.
\end{proof}

\begin{cthejemplo}{(La raíz de $t^2$ es $|t|$)}
En $X=[-1,1]$ tomemos la función positiva $g(t)=t^2$.
Como $\spec_{C(X)}(g)=g(X)=[0,1]$, la función continua
$f(s)=\sqrt s$ puede aplicarse a $g$ mediante el Teorema~2.22.
La Observación~2.25 identifica el resultado punto a punto:
\[
 \sqrt g(t)=f(g(t))=\sqrt{t^2}=|t|.
\]
La raíz positiva del elemento $g$ es $|t|$, no la función $t$,
que cambia de signo.
\end{cthejemplo}

\begin{cthlema}{(Naturalidad del cálculo funcional)}
Sean $\A,\B$ $C^*$-álgebras unitales, $\Phi:\A\to\B$ un $*$-homomorfismo unital,
$a\in\A$ normal y $f\in C(\spec_\A(a))$. Como
$\spec_\B(\Phi(a))\subset\spec_\A(a)$, se evalúa $f$ en $\Phi(a)$ por
restricción a $\spec_\B(\Phi(a))$. Entonces
\[
 \Phi(f(a))=f(\Phi(a)).
\]
\end{cthlema}
\origen{Lema 2.26}

\begin{proof}[Demostración]
El mapa $f\mapsto\Phi(f(a))$ es composición del cálculo funcional de $a$
con $\Phi$; es un $*$-homomorfismo unital y contractivo por el teorema de
contractividad anterior.
El mapa $f\mapsto f(\Phi(a))$ es la restricción
$C(\spec_\A(a))\to C(\spec_\B(\Phi(a)))$ seguida del cálculo funcional de
$\Phi(a)$, que es normal porque $\Phi$ preserva producto e involución.
También es un $*$-homomorfismo unital y contractivo. Ambos envían la función
identidad a $\Phi(a)$ y la constante $1$ a $1_\B$; al preservar la involución,
también envían $\overline z$ a $\Phi(a)^*$. Por linealidad y multiplicatividad
coinciden en los polinomios en $z$ y $\overline z$. Estos son densos por
Stone--Weierstrass, y la contractividad permite pasar al límite para todo
$f\in C(\spec_\A(a))$.
\end{proof}

\proofdiagram
{$\Phi:\A\to\B$ unital y $*$; $a\in\A$ normal}
{Inclusión espectral; cálculo funcional; contractividad}
{Los mapas $f\mapsto\Phi(f(a))$ y $f\mapsto f(\Phi(a))$ coinciden sobre la identidad y constantes}
{Por unicidad/densidad, coinciden en todo $C(\spec(a))$}

\begin{cthejemplo}{(Una raíz transportada a operadores de multiplicación)}
Sea $\A=C([0,1])$ y represéntese en $H=L^2([0,1])$ mediante
$\pi(f)=M_f$, el operador de multiplicación por $f$.
Para $a(t)=t^2$, la raíz positiva en $\A$ es
$\sqrt a(t)=t$. Por el Lema~2.26,
\[
 \pi(\sqrt a)=M_t=\sqrt{M_{t^2}}=\sqrt{\pi(a)}.
\]
En efecto, $M_t$ es positivo y $M_t^2=M_{t^2}$.
Se obtiene el mismo operador al calcular primero la raíz en funciones
o al calcularla después de representar el álgebra.
\end{cthejemplo}

\begin{cthteorema}{(Un $*$-monomorfismo unital es isométrico)}
Sean $\A,\B$ $C^*$-álgebras unitales. Si $\Phi:\A\to\B$ es un $*$-homomorfismo unital
e inyectivo, entonces
\[
 \|\Phi(a)\|=\|a\|\qquad(a\in\A).
\]
\end{cthteorema}
\origen{Teorema 2.27}

\begin{proof}[Demostración]
Ya sabemos que $\|\Phi(a)\|\le\|a\|$. Supongamos desigualdad estricta para algún $a\neq0$ y
pongamos $x=a^*a$. Entonces $x$ es autoadjunto,
$\|\Phi(x)\|=\|\Phi(a)\|^2<\|a\|^2=\|x\|$.
Como $x$ es autoadjunto y $\|x\|=r(x)$, alguno de los puntos $\pm\|x\|$ pertenece a
$\spec(x)\subset[-\|x\|,\|x\|]$.
Además, $\spec(\Phi(x))\subset[-\|\Phi(x)\|,\|\Phi(x)\|]$.
Escribimos $c=\|\Phi(x)\|$ y $d=\|x\|$, con $0\le c<d$, y definimos
\[
 f(t)=\max\!\left\{0,\min\!\left\{1,\frac{|t|-c}{d-c}\right\}\right\}
 \qquad(t\in\mathbb R).
\]
Esta función es continua, se anula en $[-c,c]$ y vale $1$ en ambos puntos
$\pm d$. Entonces
$f(\Phi(x))=0$, mientras $f(x)\neq0$ porque el cálculo funcional es isométrico
y $f$ no se anula en $\spec(x)$. Por naturalidad,
$0=f(\Phi(x))=\Phi(f(x))$, contradiciendo la inyectividad.
\end{proof}

\proofdiagrambranch
{$\Phi:\A\to\B$ unital, $*$ e inyectivo}
{$a\in\A$; $x=a^*a$ es autoadjunto}
{Si $\|\Phi(x)\|<\|x\|$, separar los dos espectros con una función continua}
{Naturalidad: $0=f(\Phi(x))=\Phi(f(x))$; inyectividad contradice $f(x)\neq0$}
{$\|\Phi(a)\|=\|a\|$}

\begin{cthejemplo}{(La representación diagonal de $C(\{1,2\})$)}
Sea $X=\{1,2\}$ y $H=\ell^2(X)\cong\mathbb C^2$.
La representación unital $\pi:C(X)\to\mathcal B(H)$ asigna a
$f$ la matriz $M_f=\operatorname{diag}(f(1),f(2))$.
Es fiel, pues $M_f=0$ implica $f(1)=f(2)=0$, y
\[
 \|\pi(f)\|=\max\{|f(1)|,|f(2)|\}=\|f\|_\infty.
\]
Este cálculo concreto coincide con la conclusión general del
Teorema~2.27: una representación fiel preserva la norma.
\end{cthejemplo}

\section{Reconocer autoadjuntos, unitarios, proyecciones y positivos por el espectro}

\begin{cthproposicion}{(Clasificación espectral de elementos normales)}
Sea $\A$ una $C^*$-álgebra unital y $a\in\A$ normal. Entonces
\begin{enumerate}[label=(\roman*)]
\item $a=a^*$ si y sólo si $\spec(a)\subset\mathbb R$;
\item $a$ es unitario si y sólo si $\spec(a)\subset\mathbb T$;
\item $a$ es una proyección si y sólo si $\spec(a)\subset\{0,1\}$.
\end{enumerate}
\end{cthproposicion}
\origen{Proposición 2.28; corresponde directamente a la Tarea 2.6}

\begin{proof}[Demostración]
Si $a=a^*$, ya sabemos $\spec(a)\subset\mathbb R$. Recíprocamente, si el espectro es real,
las funciones $z$ y $\overline z$ coinciden en $\spec(a)$, de modo que por cálculo funcional
$a=a^*$.

Si $a$ es unitario, la identidad $a^*a=1_\A$ se traduce, por el cálculo funcional
isométrico, en $|z|^2=1$ sobre $\spec(a)$; así $\spec(a)\subset\mathbb T$.
Recíprocamente, si
$\spec(a)\subset\mathbb T$, entonces $|z|^2=1$ sobre el espectro y
$a^*a=(\overline z z)(a)=1_\A$; por normalidad también $aa^*=1_\A$.

Si $a$ es proyección, $a^2-a=0$. La isometría del cálculo funcional aplicada
a $z\mapsto z^2-z$ da $z^2-z=0$ en todo $\spec(a)$; por tanto
$\spec(a)\subset\{0,1\}$. Recíprocamente, si el espectro
está contenido en $\{0,1\}$, entonces $z^2-z=0$ y $\overline z-z=0$ sobre el espectro; por cálculo
funcional $a^2=a$ y $a^*=a$.
\end{proof}

\proofdiagrambranch
{$\A$ es $C^*$ unital; $a\in\A$ normal}
{Cálculo funcional continuo}
{Traducir $a^*=a$, $a^*a=1$, $a^2=a$ a identidades de funciones sobre $\spec(a)$}
{$\overline z=z$, $|z|^2=1$, $z^2=z$ en el espectro}
{Las inclusiones espectrales caracterizan las tres clases}

\begin{cthejemplo}{(Contraejemplos: espectro pequeño sin normalidad)}
En $M_2(\mathbb C)$, tomemos
$N=\begin{pmatrix}0&1\\0&0\end{pmatrix}$ y
$J=\begin{pmatrix}1&1\\0&1\end{pmatrix}$.
Ambas son triangulares, de modo que
$\spec(N)=\{0\}\subset\mathbb R\cap\{0,1\}$ y
$\spec(J)=\{1\}\subset\mathbb T$.
Sin embargo, $N^*\neq N$ y $N^2=0\neq N$;
además, $J^*J=\begin{pmatrix}1&1\\1&2\end{pmatrix}\neq I$.
La Tarea~2.6(d) desarrolla estos casos. Muestran que las conclusiones
inversas de la Proposición~2.28 necesitan la normalidad.
\end{cthejemplo}

\begin{cthproposicion}{(Cuatro unitarios generan linealmente cada elemento)}
Sea $\A$ una $C^*$-álgebra unital. Cada $a\in\A$ es combinación lineal de a lo más cuatro elementos
unitarios.
\end{cthproposicion}
\origen{Proposición 2.29}

\begin{proof}[Demostración]
Escribimos $a=x_1+ix_2$ con $x_j$ autoadjuntos. Sean
$r_j=\max\{1,\|x_j\|\}$ y $h_j=x_j/r_j$, de modo que $\|h_j\|\le1$.
Como $h_j=h_j^*$, su espectro está en $[-1,1]$. La función continua
$f(t)=t+i\sqrt{1-t^2}$ satisface
$|f(t)|^2=t^2+(1-t^2)=1$ y $f(t)+\overline{f(t)}=2t$ en ese intervalo.
Por cálculo funcional, $u_j=f(h_j)$ cumple
$u_j^*u_j=u_ju_j^*=(|f|^2)(h_j)=1_\A$ y es unitario. Además,
$u_j+u_j^*=(f+\overline f)(h_j)=2h_j$. Por tanto
\[
 a=\frac{r_1}{2}(u_1+u_1^*)+\frac{i r_2}{2}(u_2+u_2^*),
\]
la combinación anunciada.
\end{proof}

\proofdiagram
{$\A$ es $C^*$ unital; $a=x_1+ix_2$ con $x_j=x_j^*$}
{Escalar $h_j=x_j/r_j$ con $r_j=\max\{1,\|x_j\|\}$}
{$u_j=f(h_j)$, $f(t)=t+i\sqrt{1-t^2}$; $x_j=r_j(u_j+u_j^*)/2$}
{$a$ es combinación de $u_1,u_1^*,u_2,u_2^*$}

\begin{cthejemplo}{(Un autoadjunto como media de dos unitarios)}
En $M_2(\mathbb C)$ consideremos
$H=\operatorname{diag}(1/2,-1/2)$.
Es autoadjunto y $\|H\|\le1$. La construcción de la
Proposición~2.29 produce
\[
 U=H+i\sqrt{I-H^2}
 =\operatorname{diag}\!\left(
 \frac12+i\frac{\sqrt3}{2},
 -\frac12+i\frac{\sqrt3}{2}\right).
\]
Cada entrada diagonal tiene módulo uno, así que $U$ es unitario, y
$H=(U+U^*)/2$. Este caso exhibe el paso que permite escribir
cualquier elemento como combinación de hasta cuatro unitarios.
\end{cthejemplo}

\begin{cthproposicion}{(Raíz cuadrada positiva)}
Sea $\A$ una $C^*$-álgebra unital y $a\in\A$ normal con
$\spec_\A(a)\subset[0,\infty)$. Entonces
existe un elemento autoadjunto positivo
$y\in C^*(a,1_\A)$ tal que
\[
 y^2=a.
\]
La raíz obtenida por cálculo funcional se denota $y=\sqrt a$.
\end{cthproposicion}
\origen{Proposición 2.30}

\begin{proof}[Demostración]
La función $t\mapsto\sqrt t$ es continua y real en $\spec(a)$. Definimos
$y=\sqrt{\,\cdot\,}(a)$. Entonces $y^*=y$ y $y^2=(t\mapsto t)(a)=a$.
La identificación
$C^*(a,1_\A)\cong C(\spec_\A(a))$ muestra que la transformada de $y$ es
$t\mapsto\sqrt t\ge0$; por la descripción del espectro en un álgebra de funciones
y la permanencia espectral, $\spec_\A(y)\subset[0,\infty)$.
\end{proof}

\proofdiagram
{$\A$ es $C^*$ unital; $a\in\A$ normal y $\spec_\A(a)\subset[0,\infty)$}
{Cálculo funcional continuo}
{Aplicar $f(t)=\sqrt t$}
{$y=f(a)$ es autoadjunto, $\spec(y)\subset[0,\infty)$ y $y^2=a$}

\begin{cthejemplo}{(Módulo de una matriz nilpotente)}
Sea $A=\begin{pmatrix}0&2\\0&0\end{pmatrix}$ en $M_2(\mathbb C)$.
Aunque $A^2=0$, el producto
\[
 A^*A=\operatorname{diag}(0,4)
\]
es positivo. La Proposición~2.30 permite tomar su raíz positiva, que aquí
se calcula directamente:
\[
 |A|=(A^*A)^{1/2}=\operatorname{diag}(0,2).
\]
Este $|A|$ conserva la magnitud de la acción de $A$ y es el objeto usado
en la Tarea~2.7 para la descomposición polar. En general, aplicar
la proposición a $a^*a$ requiere el hecho estándar $a^*a\ge0$.
\end{cthejemplo}

\begin{investigacion}
\textbf{Ejemplo directo.} En \cite[Sec.~1, p.~20]{BFK2006}, los operadores
$(U_z\varphi)(t)=\varphi(zt)$, con $z,t\in\mathbb T$ y
$\varphi\in L^2(\mathbb T)$, forman una representación unitaria del grupo
$\mathbb T$ sobre $L^2(\mathbb T)$; las rotaciones preservan la medida de longitud.
\textbf{Lo que conviene reconocer:} el unitario abstracto aquí implementa una
simetría geométrica concreta y participa en la extensión no local del álgebra
estudiada en esa fuente.
\end{investigacion}

\section*{Nota histórica}
\origen{Nota histórica del capítulo base, pp.~8--10}

\paragraph{Los axiomas originales.}
El trabajo de Gelfand y Naimark de 1943 buscaba representar $C^*$-álgebras,
también no conmutativas, mediante operadores acotados en espacios de Hilbert.
Las notas base recuerdan que la definición original incluía, además de
$\|a^*a\|=\|a\|^2$, la isometría de la involución y la invertibilidad de
$1+a^*a$. Después se demostró que estas condiciones adicionales se deducen de
la identidad $C^*$.

\paragraph{$B^*$ y $C^*$.}
Durante un tiempo, $B^*$ designó la estructura axiomática y $C^*$ las
subálgebras cerradas por adjuntos de operadores acotados en Hilbert. El teorema
general de representación de Gelfand--Naimark identifica ambas descripciones;
el teorema conmutativo demostrado aquí es el paso del curso hacia esa idea.

\paragraph{Rigidez y cálculo funcional.}
El Teorema~2.6, el Corolario~2.7 y el Teorema~2.27 de las notas base muestran tres formas de
rigidez de la norma. La permanencia espectral contrasta con el álgebra del disco
de la Tarea~1.12, que es una subálgebra cerrada de $C(\mathbb T)$ pero no una
$C^*$-subálgebra. El cálculo funcional transporta funciones continuas del
espectro a elementos del álgebra sin elegir primero una representación en
Hilbert; Stone--Weierstrass explica por qué las funciones de $a$ y $a^*$ bastan.

\paragraph{Topología y positividad.}
La clasificación conmutativa permite recuperar el compacto $X$ desde los
caracteres de $C(X)$. Esta lectura motivó el lenguaje posterior de topología no
conmutativa. La raíz cuadrada de la Proposición~2.30 de las notas base abre el estudio sistemático
de elementos positivos y del orden, mientras la Proposición~2.29 de esas notas anticipa la
relevancia de los unitarios para la geometría del álgebra.

\chapter{Tarea 1 resuelta y conectada con la teoría}

\begin{context}
Al terminar esta tarea, podrás justificar qué continuidad hace posibles las construcciones por
clausura, decidir en qué álgebra se calcula un espectro y reconstruir la subálgebra conmutativa que
permite estudiar un autoadjunto. Cada solución identifica los resultados que utiliza y el uso
posterior de su conclusión.
\end{context}

\section{Problema 1: continuidad de multiplicación e inversión}

Sea $\A$ un álgebra de Banach. En el inciso (b) se supone además que tiene unidad.

\subsection*{1(a) La multiplicación es conjuntamente continua}

\noindent\textbf{Enunciado.} Para $x_n\to x$ y $y_n\to y$ en $\A$, prueba que $x_ny_n\to xy$; equivale a la continuidad conjunta de $\A\times\A\to\A$.

\begin{proof}[Solución]
Como $x_n\to x$, la sucesión $(x_n)$ es acotada: existe $M$ con $\|x_n\|\le M$. Entonces
\[
\|x_ny_n-xy\|
\le \|x_n\|\,\|y_n-y\|+\|x_n-x\|\,\|y\|
\le M\|y_n-y\|+\|y\|\|x_n-x\|\to0.
\]
Así la multiplicación $\A\times\A\to\A$ es continua en ambas variables de manera conjunta.
\end{proof}

\proofdiagram
{$x_n\to x$, $y_n\to y$}
{$\|ab\|\le\|a\|\|b\|$}
{Escribir $x_ny_n-xy=x_n(y_n-y)+(x_n-x)y$}
{$\|x_ny_n-xy\|\to0$}

\begin{cthejemplo}{(Producto de operadores que convergen)}
En $\mathcal B(H)$ comprobamos T1.1(a) con $T_n=T+n^{-1}I$ y $S_n=S+n^{-1}I$.
Entonces $T_nS_n-TS=n^{-1}(T+S)+n^{-2}I$, de donde
$\|T_nS_n-TS\|\le n^{-1}(\|T\|+\|S\|)+n^{-2}\to0$.
El cálculo muestra que el producto converge aunque cambien simultáneamente los dos factores.
\end{cthejemplo}

\noindent\textbf{Dependencias.} La desigualdad submultiplicativa de la definición de álgebra de Banach y el hecho de que toda sucesión convergente está acotada. La estimación separa la variación de cada factor.

\noindent\textbf{Qué prepara este ejercicio.} Permite cerrar por límites las álgebras generadas en T1.3, pasar al límite en la normalidad en T1.7 y controlar los polinomios en $x_n,x_n^*$ de T2.8.

\subsection*{1(b) La inversión es continua en $\GL(\A)$}

Para este inciso, $\A$ es un álgebra de Banach unital y $\GL(\A)$ lleva la norma heredada de
$\A$.

\noindent\textbf{Enunciado.} Prueba que la aplicación $\GL(\A)\to\GL(\A)$, $a\mapsto a^{-1}$, es continua en la norma de $\A$.

\begin{proof}[Solución]
Primero, si $x\to1$ y $\|1-x\|<1$, la serie de Neumann da
\[
 x^{-1}=\sum_{n=0}^\infty(1-x)^n,
 \qquad
 \|x^{-1}-1\|
 \le \sum_{n=1}^\infty\|1-x\|^n
 =\frac{\|1-x\|}{1-\|1-x\|}\to0.
\]
Por tanto la inversión es continua en $1$. Si ahora $a_n\to a\in\GL(\A)$, entonces
$a^{-1}a_n\to1$ y, para $n$ grande, $a_n$ es invertible. Además
\[
 a_n^{-1}=(a^{-1}a_n)^{-1}a^{-1}\longrightarrow a^{-1}.
\]
\end{proof}

\proofdiagram
{$a_n\to a\in\GL(\A)$}
{Continuidad en $1$ por serie de Neumann}
{$a_n=a(a^{-1}a_n)$ y $a^{-1}a_n\to1$}
{$a_n^{-1}=(a^{-1}a_n)^{-1}a^{-1}\to a^{-1}$}

\begin{cthejemplo}{(Inversión de una perturbación pequeña)}
En $\mathcal B(H)$ interesa saber cuándo una perturbación de la identidad conserva inverso.
Si $E\in\mathcal B(H)$ cumple $\|E\|<1$, la serie de Neumann da
\[
 (I+E)^{-1}=\sum_{k=0}^{\infty}(-E)^k,
 \qquad \|(I+E)^{-1}-I\|\le\frac{\|E\|}{1-\|E\|}.
\]
La cota tiende a cero con $\|E\|$: éste es el mecanismo concreto de continuidad de T1.1(b).
\end{cthejemplo}

\noindent\textbf{Dependencias.} La serie de Neumann de la Parte 1 da un inverso para $1-h$ si $\|h\|<1$; la apertura de $\GL(\A)$ permite aplicar esa fórmula cerca de cualquier invertible.

\noindent\textbf{Qué prepara este ejercicio.} La demostración de permanencia espectral de la Sección 2 usa esta continuidad al pasar de $(b+i/n)^{-1}$ a $b^{-1}$ dentro de una subálgebra cerrada.

\Needspace{14\baselineskip}
\section{Problema 2: conmutación e inversos}

Sea $\A$ un álgebra compleja unital, $a,b\in\A$ y $a$ invertible. En el inciso (b), $\A$ tiene
además una involución.

Recordemos que $x$ y $y$ \emph{$*$-conmutan} cuando
\[
 xy=yx,
 \qquad x^*y=yx^*.
\]

\subsection*{2(a) Conmutar con $a$ equivale a conmutar con $a^{-1}$}

\noindent\textbf{Enunciado.} Prueba que $ab=ba$ si y sólo si $a^{-1}b=ba^{-1}$.

\begin{proof}[Solución]
Si $ab=ba$, multiplicando por $a^{-1}$ a la izquierda y luego a la derecha,
\[
 b=a^{-1}ba \quad\Longrightarrow\quad ba^{-1}=a^{-1}b.
\]
Recíprocamente, si $a^{-1}b=ba^{-1}$, multiplicamos primero a la izquierda y después a la
derecha por $a$: de $b=aba^{-1}$ resulta $ba=ab$. Así también $a$ conmuta con $b$.
\end{proof}

\proofdiagram
{$ab=ba$, $a$ invertible}
{Multiplicar por $a^{-1}$ a izquierda y derecha}
{$a^{-1}b=ba^{-1}$}
{Equivalencia, usando que $(a^{-1})^{-1}=a$}

\begin{cthejemplo}{(Conmutación de matrices diagonales e inversos)}
En $M_2(\mathbb C)$ podemos ver la equivalencia de T1.2(a) sin recurrir a límites.
Tomemos $A=\operatorname{diag}(2,3)$ y $B=\operatorname{diag}(5,7)$.
Como $A^{-1}=\operatorname{diag}(1/2,1/3)$,
\[
 AB=BA=\operatorname{diag}(10,21),\qquad
 A^{-1}B=BA^{-1}=\operatorname{diag}(5/2,7/3).
\]
El inverso cambia las entradas de $A$, pero conserva exactamente la conmutación con $B$.
\end{cthejemplo}

\noindent\textbf{Dependencias.} Sólo se necesitan la asociatividad, la unidad y la existencia de $a^{-1}$; no interviene la norma.

\noindent\textbf{Qué prepara este ejercicio.} En T1.9 se aplica esta identidad a $\lambda1-c$ para mostrar que su inverso conmuta con la subálgebra elegida.

\subsection*{2(b) $*$-conmutación e inversos}

\noindent\textbf{Enunciado.} Prueba que $a$ y $b$ $*$-conmutan si y sólo si $a^{-1}$ y $b$ $*$-conmutan.

\begin{proof}[Solución]
De $ab=ba$ y el inciso (a), $a^{-1}b=ba^{-1}$. Además $a^*b=ba^*$; como $a^*$ es invertible,
$(a^*)^{-1}b=b(a^*)^{-1}$. Pero $(a^*)^{-1}=(a^{-1})^*$, así que se satisfacen las dos
condiciones de $*$-conmutación. Si, recíprocamente, $a^{-1}$ y $b$ $*$-conmutan, aplicamos el
argumento ya probado al invertible $a^{-1}$: su inverso $a$ y su adjunto $a^*$ conmutan con $b$.
\end{proof}

\Needspace{14\baselineskip}
\proofdiagram
{$a$ y $b$ $*$-conmutan}
{Inciso (a) aplicado a $a$ y a $a^*$}
{$a^{-1}$ y $(a^{-1})^*=(a^*)^{-1}$ conmutan con $b$}
{$a^{-1}$ y $b$ $*$-conmutan}

\begin{cthejemplo}{(El inverso de una matriz normal)}
Para comprobar la $*$-conmutación en T1.2(b), usemos la matriz normal e invertible
$A=\operatorname{diag}(2,i)$ en $M_2(\mathbb C)$.
Se tiene $A^{-1}=\operatorname{diag}(1/2,-i)$ y
\[
 A^{-1}(A^{-1})^*=(A^{-1})^*A^{-1}=\operatorname{diag}(1/4,1).
\]
Así $A^{-1}$ sigue siendo normal: el paso decisivo es que su adjunto también es el inverso de $A^*$.
\end{cthejemplo}

\noindent\textbf{Dependencias.} Se aplica T1.2(a) tanto a $a$ como a $a^*$, junto con $(a^{-1})^*=(a^*)^{-1}$. La definición de $*$-conmutación fija las dos relaciones que hay que verificar.

\noindent\textbf{Qué prepara este ejercicio.} Permite controlar los inversos en la versión involutiva de T1.9 y explica por qué el inverso de un elemento normal invertible vuelve a ser normal.

\section{Problema 3: álgebras generadas por un conjunto}

Sea $\A$ un álgebra de Banach y $S\subset\A$. Definimos
\[
\mathcal B(S)=\overline{\spn\{s_1\cdots s_k:k\ge1,\ s_j\in S\}}.
\]
Si $\A$ tiene una involución, ponemos $S^*=\{s^*:s\in S\}$ y
$\mathcal B^*(S)=\mathcal B(S\cup S^*)$.

\subsection*{3(a) Minimalidad de $\mathcal B(S)$ y $\mathcal B^*(S)$}

\noindent\textbf{Enunciado.} Prueba que $\mathcal B(S)$ es la mínima subálgebra de Banach de $\A$ que contiene $S$ y, cuando $\A$ es involutiva, que $\mathcal B^*(S)$ es la mínima subálgebra de Banach involutiva que contiene $S$.

\begin{proof}[Solución]
El conjunto de combinaciones lineales de palabras en $S$ es una subálgebra; su clausura sigue
siendo subálgebra por continuidad de la multiplicación. Así $\mathcal B(S)$ es una subálgebra de
Banach que contiene $S$. Si $C\subset\A$ es cualquier subálgebra de Banach con $S\subset C$,
entonces $C$ contiene todas las palabras $s_1\cdots s_k$, sus combinaciones lineales y, por ser
cerrada, su clausura. Por tanto $\mathcal B(S)\subset C$.

Para $\mathcal B^*(S)$ se aplica lo mismo a $S\cup S^*$. Ese conjunto generado es estable bajo
$*$: el adjunto de una palabra es otra palabra en $S\cup S^*$, y la involución es continua. Por ello
es la mínima subálgebra de Banach involutiva que contiene $S$.
\end{proof}

\pagebreak
\proofdiagram
{$S\subset\A$}
{Productos finitos, combinaciones lineales y clausura}
{Toda subálgebra cerrada que contiene $S$ contiene cada palabra y su límite}
{$\mathcal B(S)$ es la mínima subálgebra de Banach que contiene $S$}

\begin{cthejemplo}{(Añadir la unidad cambia el álgebra generada)}
En $M_2(\mathbb C)$, el operador $T=\operatorname{diag}(0,2)$ permite distinguir
$\mathcal B(\{T\})$ de la álgebra unital generada por $T$.
Como $T^k=\operatorname{diag}(0,2^k)$ para $k\ge1$,
\[
 \mathcal B(\{T\})=\{\operatorname{diag}(0,c):c\in\mathbb C\}.
\]
Al añadir $I$ se obtienen todas las matrices diagonales. La distinción responde a la definición
por palabras de longitud positiva en T1.3(a).
\end{cthejemplo}

\noindent\textbf{Dependencias.} La clausura de combinaciones lineales de palabras es completa por ser cerrada en $\A$; T1.1(a) garantiza que el producto pasa a la clausura. La continuidad de $*$ da el caso involutivo.

\noindent\textbf{Qué prepara este ejercicio.} Construye de modo concreto $C^*(a,1)$, la subálgebra conmutativa que se emplea en la Sección 2 para Gelfand--Naimark y el cálculo funcional. La misma operación de generación aparece en álgebras de operadores singulares.

\begin{investigacion}[title={Generación de un álgebra real de operadores}]
Para reconocer T1.3(a) en una construcción de investigación, consideremos la álgebra
\emph{real} de operadores en $\mathcal B(\mathbb R L^p(\Gamma,w))$ estudiada por Karlovich.
Sus generadores incluyen $M_a$ con $a\in\SO^\diamond(\Gamma)$, $S_\Gamma$ y la
conjugación $C$, que es real-lineal. Por ejemplo, una palabra $M_aS_\Gamma C$ pertenece
a la álgebra generada; también pertenecen las sumas reales de palabras y sus límites en
norma \cite[ec.~(1.4), p.~3]{Karlovich2025}.
La operación de generación es la de T1.3(a) sobre escalares reales. La presencia de $C$
sitúa este ejemplo fuera del cálculo funcional complejo de la Sección 2.
\end{investigacion}

\subsection*{3(b) La conmutación se propaga al álgebra generada}

\noindent\textbf{Enunciado.} Si $a\in\A$ conmuta con cada $s\in S$, prueba que conmuta con todo $x\in\mathcal B(S)$. En el caso involutivo, si $a$ $*$-conmuta con cada $s\in S$, prueba la misma conclusión para $\mathcal B^*(S)$ y la $*$-conmutación.

\begin{proof}[Solución]
Si $a$ conmuta con cada $s\in S$, entonces, por inducción, conmuta con cada palabra
$s_1\cdots s_k$. Por linealidad conmuta con toda combinación lineal de palabras. Si
$x_n\to x$ y $ax_n=x_na$, la continuidad del producto da $ax=x a$. Así conmuta con toda
$\mathcal B(S)$.

En el caso involutivo, si $a$ $*$-conmuta con cada $s\in S$, entonces $a$ y $a^*$ conmutan con
$s$ y, tomando adjuntos, también con $s^*$. Por el argumento anterior, $a$ y $a^*$ conmutan con
toda $\mathcal B^*(S)$, que es exactamente decir que $a$ $*$-conmuta con todos sus elementos.
\end{proof}

\proofdiagram
{$as=sa$ para todo $s\in S$}
{La conmutación se preserva por productos y sumas}
{Pasar a límites por continuidad}
{$ax=xa$ para todo $x\in\mathcal B(S)$}

\begin{cthejemplo}{(Un conmutante de matrices diagonales)}
En $M_2(\mathbb C)$ queremos seguir la conmutación desde un generador hasta su álgebra.
Sean $D=\operatorname{diag}(1,2)$ y $A=\operatorname{diag}(3,4)$.
Para cualquier polinomio $p$,
\[
 Ap(D)=\operatorname{diag}(3p(1),4p(2))=p(D)A.
\]
La igualdad persiste en límites de polinomios en norma, tal como establece T1.3(b).
\end{cthejemplo}

\noindent\textbf{Dependencias.} La identidad $as=sa$ se conserva al formar palabras y combinaciones lineales; T1.1(a) permite pasar al límite. Para $*$-conmutación se aplican además las relaciones adjuntas.

\noindent\textbf{Qué prepara este ejercicio.} La propagación de conmutación justifica en T1.9 que se pueda añadir un inverso sin perder conmutatividad y permite analizar el conmutante de familias de operadores.

\subsection*{3(c) Si $S$ conmuta dos a dos, $\mathcal B(S)$ es conmutativa}

\noindent\textbf{Enunciado.} Si $st=ts$ para todo $s,t\in S$, prueba que $\mathcal B(S)$ es conmutativa.

\begin{proof}[Solución]
Dos palabras $s_1\cdots s_k$ y $t_1\cdots t_m$ pueden reordenarse factor a factor porque todos
los generadores conmutan. Por tanto las palabras conmutan. El resultado se extiende por bilinealidad
a sus combinaciones lineales y por continuidad a la clausura.
\end{proof}

\proofdiagram
{Los elementos de $S$ conmutan dos a dos}
{Reordenar factores en dos palabras}
{Las palabras conmutan; luego sus combinaciones lineales también}
{La clausura $\mathcal B(S)$ es conmutativa}

\begin{cthejemplo}{(Álgebra de un nilpotente)}
La conmutatividad de una álgebra generada por un solo elemento no exige que éste sea normal.
En $M_2(\mathbb C)$, tomemos $N=\begin{pmatrix}0&1\\0&0\end{pmatrix}$.
Como $N^2=0$, toda palabra de longitud positiva es $N$ o $0$, y
$\mathcal B(\{N\})=\mathbb C N$.
Dos elementos $cN,dN$ tienen producto $0$ en ambos órdenes. Esto concreta T1.3(c).
\end{cthejemplo}

\noindent\textbf{Dependencias.} Se usa la construcción de T1.3(a): si los generadores conmutan dos a dos, dos palabras cualesquiera se pueden reordenar. T1.1(a) extiende la identidad a sus límites.

\noindent\textbf{Qué prepara este ejercicio.} Permite reconocer la conmutatividad de $C^*(a,1)$ cuando $a$ es normal, paso necesario para representar esa álgebra como funciones continuas en la Sección 2.

\subsection*{3(d) $*$-conmutación dos a dos}

\noindent\textbf{Enunciado.} En el caso involutivo, si cada par de elementos de $S$ $*$-conmuta, prueba que $\mathcal B^*(S)$ es una subálgebra de Banach involutiva y conmutativa.

\begin{proof}[Solución]
De $st=ts$ y $s^*t=ts^*$ se obtienen por adjunción también
$st^*=t^*s$ y $s^*t^*=t^*s^*$. Por tanto todos los elementos de $S\cup S^*$ conmutan dos a dos.
El inciso (c) aplicado a $S\cup S^*$ muestra que $\mathcal B^*(S)$ es conmutativa; por construcción
es una álgebra de Banach involutiva.
\end{proof}

\proofdiagram
{Cada par $s,t\in S$ $*$-conmuta}
{Tomar adjuntos de las relaciones}
{$S\cup S^*$ conmuta dos a dos}
{$\mathcal B^*(S)$ es conmutativa e involutiva}

\begin{cthejemplo}{(Un normal genera una álgebra involutiva conmutativa)}
En $M_2(\mathbb C)$, la matriz $A=\operatorname{diag}(1,i)$ muestra por qué
la normalidad permite usar cálculo funcional.
Su adjunto es $A^*=\operatorname{diag}(1,-i)$, de modo que $AA^*=A^*A=I$.
Así las palabras en $A$, $A^*$ e $I$ son diagonales y conmutan; también lo hace su clausura.
Por T1.3(d), $C^*(A,I)$ es conmutativa, condición de entrada de la Sección 2.
\end{cthejemplo}

\noindent\textbf{Dependencias.} La definición de $*$-conmutación aplicada a cada par de generadores da conmutación en $S\cup S^*$; luego se aplica T1.3(c).

\noindent\textbf{Qué prepara este ejercicio.} Si $a$ es normal, al aplicar el inciso a $S=\{a,1\}$ se obtiene la $C^*$-subálgebra conmutativa $C^*(a,1)$, punto de partida del cálculo funcional continuo.

\section{Problema 4: el cociente por un ideal es una álgebra de Banach}

Sea $\A$ un álgebra de Banach y $J\subset\A$ un ideal bilateral cerrado.

\noindent\textbf{Enunciado.} Prueba que $(x+J)(y+J)=xy+J$ está bien definido y que $\A/J$, con $\|x+J\|=\inf_{j\in J}\|x-j\|$, es una álgebra de Banach.

\begin{proof}[Solución]
Si $x'=x+j_1$ y $y'=y+j_2$ con $j_1,j_2\in J$, entonces
\[
 x'y'-xy=xj_2+j_1y+j_1j_2\in J,
\]
de modo que $(x+J)(y+J)=xy+J$ está bien definido.

Para la norma, dados $j_1,j_2\in J$,
\[
(x-j_1)(y-j_2)=xy-\underbrace{(xj_2+j_1y-j_1j_2)}_{\in J},
\]
por lo que
\[
\|xy+J\|\le\|x-j_1\|\,\|y-j_2\|.
\]
Tomando ínfimo en $j_1$ y $j_2$ se obtiene
$\|(x+J)(y+J)\|\le\|x+J\|\|y+J\|$. Finalmente, como $J$ es cerrado, el espacio cociente
$\A/J$ es completo. Así es una álgebra de Banach.
\end{proof}

\proofdiagrambranch
{$J$ ideal bilateral cerrado}
{$\|x+J\|=\inf_{b\in J}\|x-b\|$}
{Comprobar independencia de representantes y submultiplicatividad}
{Usar que el cociente por un subespacio cerrado es Banach}
{$\A/J$ es una álgebra de Banach}

\begin{cthejemplo}{(Una clase en el álgebra de Calkin)}
Para ver qué identifica el cociente de T1.4, trabajemos en $X=\ell^2(\mathbb N)$.
Sea $P$ la proyección de rango uno sobre $\mathbb C e_1$; entonces $P\in\K(X)$ y
\[
 (I+P)+\K(X)=I+\K(X).
\]
Dos operadores distintos representan la misma clase porque su diferencia es compacta.
Ésta es la igualdad básica que permite trabajar en el álgebra de Calkin.
\end{cthejemplo}

\noindent\textbf{Dependencias.} Un ideal bilateral asegura que el producto de clases no depende de representantes; que $J$ sea cerrado asegura la completitud del cociente. Se usa la norma cociente y la submultiplicatividad de $\A$.

\noindent\textbf{Qué prepara este ejercicio.} Al tomar $\A=\mathcal B(X)$ y $J=\K(X)$ se obtiene el álgebra de Calkin. La equivalencia entre Fredholm e invertibilidad en ese cociente es el teorema de Atkinson, un resultado adicional que este ejercicio permite formular pero no demuestra.

\begin{investigacion}[title={Fredholm en el álgebra de Calkin}]
Para situar T1.4 en una aplicación, sea $X$ un Banach de dimensión infinita y
$\pi:\mathcal B(X)\to\mathcal B(X)/\K(X)$ el mapa cociente.
La clase $\pi(A)$ es invertible exactamente cuando existe $B\in\mathcal B(X)$ con
$AB-I,BA-I\in\K(X)$; ambas igualdades se leen multiplicando clases en el cociente.
El teorema de Atkinson añade que esa condición equivale a que $A$ sea Fredholm: rango
cerrado, núcleo finito dimensional y cociente $X/\operatorname{Ran}(A)$ finito dimensional
\cite[p.~2]{Karlovich2025}. T1.4 proporciona el álgebra donde se formula esta
equivalencia; el teorema de Atkinson es un resultado adicional.
\end{investigacion}

\section{Problema 5: $C_0(X)$ es una subálgebra de Banach de $C_b(X)$}

Sea $X$ un espacio localmente compacto y Hausdorff.

\noindent\textbf{Enunciado.} Definimos $C_0(X)=\{f\in C_b(X):\text{ para todo }\varepsilon>0\text{ existe un compacto }K\subset X\text{ con }|f(x)|<\varepsilon\text{ fuera de }K\}$. Prueba que es una subálgebra de Banach de $C_b(X)$.

\begin{proof}[Solución]
Toda función continua que se desvanece al infinito es acotada: para $\varepsilon=1$ es menor
que $1$ fuera de un compacto, y sobre ese compacto alcanza un máximo finito.
Además, si $f,h\in C_0(X)$, elegimos compactos $K_f,K_h$ fuera de los cuales
$|f|,|h|<\varepsilon/2$; fuera de $K_f\cup K_h$ se tiene $|f+h|<\varepsilon$.
La multiplicación por escalares se verifica escogiendo el umbral
$\varepsilon/(1+|c|)$ para $cf$. Así $C_0(X)$ es un subespacio lineal de $C_b(X)$.

Si $f\in C_0(X)$ y $g\in C_b(X)$, dado $\varepsilon>0$ elegimos un compacto $K$ tal que
$|f(x)|<\varepsilon/(1+\|g\|_\infty)$ fuera de $K$. Entonces
$|f(x)g(x)|<\varepsilon$ fuera de $K$, así que $fg\in C_0(X)$. Esto muestra incluso que $C_0(X)$
es un ideal de $C_b(X)$.

Para ver que es cerrado, sea $f_n\in C_0(X)$ y $f_n\to f$ uniformemente. Dado $\varepsilon>0$,
elegimos $n$ con $\|f-f_n\|_\infty<\varepsilon/2$ y un compacto $K$ con
$|f_n(x)|<\varepsilon/2$ fuera de $K$. Entonces $|f(x)|<\varepsilon$ fuera de $K$.
Por tanto $f\in C_0(X)$. Un subespacio cerrado de un Banach es Banach.

La conjugación puntual preserva $C_0(X)$ y
$\|f^*f\|_\infty=\sup_x|f(x)|^2=\|f\|_\infty^2$; por ello $C_0(X)$ es además una
$C^*$-álgebra, aunque no tenga unidad cuando $X$ no es compacto.
\end{proof}

\proofdiagram
{$f\in C_0(X)$, $g\in C_b(X)$}
{Definición de desvanecimiento al infinito}
{$fg$ se hace pequeño fuera de un compacto; límites uniformes preservan la propiedad}
{$C_0(X)$ es un ideal cerrado de $C_b(X)$, luego Banach}

\begin{cthejemplo}{(Una función que se desvanece en la recta)}
En $X=\mathbb R$ queremos comparar una función que se desvanece con la unidad de $C_b(X)$.
Para $f(x)=1/(1+x^2)$ y $\varepsilon>0$, si $|x|>\varepsilon^{-1/2}$ entonces
$|f(x)|<\varepsilon$; por ello $f\in C_0(\mathbb R)$.
La función constante $1$ no cumple esa condición para $\varepsilon=1/2$ fuera de ningún
compacto. T1.5 sitúa a $f$ en el ideal cerrado $C_0(\mathbb R)$ y deja la unidad fuera de él.
\end{cthejemplo}

\noindent\textbf{Dependencias.} La definición de desvanecimiento al infinito permite escoger un compacto para cada tolerancia. La norma uniforme de $C_b(X)$ controla tanto productos como límites.

\noindent\textbf{Qué prepara este ejercicio.} Si $X$ no es compacto, $C_0(X)$ es el modelo conmutativo no unital; para $X$ compacto coincide con $C(X)$. La clasificación general de $C_0(X)$ requiere la versión no unital de Gelfand--Naimark, mientras la Sección 2 demuestra la versión unital.

\section{Problema 6: espectro e involución}

Sea $\A$ un álgebra de Banach involutiva unital y $x\in\A$.

\noindent\textbf{Enunciado.} Para cada $x\in\A$, prueba $\spec(x^*)=\overline{\spec(x)}$.

\begin{proof}[Solución]
Un elemento $a$ es invertible si y sólo si $a^*$ lo es, pues de $ab=ba=1$ se obtiene
$b^*a^*=a^*b^*=1$. Ahora
\[
(\lambda1-x^*)^*=\overline\lambda1-x.
\]
Por tanto
\[
\lambda1-x^*\notin\GL(\A)
\Longleftrightarrow
\overline\lambda1-x\notin\GL(\A),
\]
lo que equivale a
\[
 \spec(x^*)=\overline{\spec(x)}.
\]
\end{proof}

\proofdiagram
{$\lambda1-x^*$}
{La involución preserva invertibilidad: $a^{-1}\mapsto(a^*)^{-1}$}
{$(\lambda1-x^*)^*=\overline\lambda1-x$}
{$\lambda\in\spec(x^*)\Longleftrightarrow\overline\lambda\in\spec(x)$}

\begin{cthejemplo}{(Espectro del adjunto no normal)}
La simetría de T1.6 no necesita normalidad. Consideremos
$A=\begin{pmatrix}1&1\\0&i\end{pmatrix}$ en $M_2(\mathbb C)$.
Por triangularidad, $\spec(A)=\{1,i\}$; su adjunto es triangular inferior con diagonal
$(1,-i)$, así que $\spec(A^*)=\{1,-i\}=\overline{\spec(A)}$.
La conjugación del espectro proviene de la involución, aunque $A$ no conmute con $A^*$.
\end{cthejemplo}

\noindent\textbf{Dependencias.} La definición de espectro de la Parte 1 y la identidad $(a^*)^{-1}=(a^{-1})^*$ permiten transportar invertibilidad a través de la involución.

\noindent\textbf{Qué prepara este ejercicio.} Muestra una simetría espectral válida antes de imponer la identidad $C^*$. Para concluir que todo el espectro de un autoadjunto es real, la Sección 2 añade argumentos propios de $C^*$-álgebras.

\section{Problema 7: los elementos normales forman un conjunto cerrado}

Sea $\A$ una $C^*$-álgebra y $x_n\to x$ en norma, con cada $x_n$ normal.

\noindent\textbf{Enunciado.} Prueba que el conjunto $\{x\in\A:x^*x=xx^*\}$ de elementos normales es cerrado en norma.

\begin{proof}[Solución]
Como $\|x_n^*-x^*\|=\|x_n-x\|\to0$ y la multiplicación es conjuntamente continua,
\[
 x_n^*x_n\to x^*x,
 \qquad
 x_nx_n^*\to xx^*.
\]
Pero ambos términos son iguales para cada $n$. Por unicidad del límite, $x^*x=xx^*$.
\end{proof}

\proofdiagram
{$x_n\to x$, $x_n^*x_n=x_nx_n^*$}
{La involución es isométrica; producto conjuntamente continuo}
{$x_n^*\to x^*$ y ambos productos convergen}
{$x^*x=xx^*$}

\begin{cthejemplo}{(Límite de matrices normales)}
En $M_2(\mathbb C)$ podemos observar el cierre del conjunto de normales sin cambiar de dimensión.
Sea $A_n=\operatorname{diag}(1+1/n,i-1/n)$; cada $A_n$ es normal y
$\|A_n-\operatorname{diag}(1,i)\|=1/n$.
El límite $A=\operatorname{diag}(1,i)$ también conmuta con su adjunto.
Éste es un caso concreto del paso al límite demostrado en T1.7.
\end{cthejemplo}

\noindent\textbf{Dependencias.} La involución isométrica da $x_n^*\to x^*$ y T1.1(a) permite pasar al límite en los dos productos de la ecuación de normalidad.

\noindent\textbf{Qué prepara este ejercicio.} Es el primer paso de T2.8: el límite $x$ sigue siendo normal, de modo que $f(x)$ está definido por cálculo funcional continuo.

\section{Problema 8: inclusión espectral bajo homomorfismos}

Sean $\A,\B$ álgebras de Banach unitales y $\psi:\A\to\B$ un homomorfismo complejo unital.

\noindent\textbf{Enunciado.} Para cada $x\in\A$, prueba $\spec_\B(\psi(x))\subset\spec_\A(x)$ y $r_\B(\psi(x))\le r_\A(x)$.

\begin{proof}[Solución]
Si $\lambda\notin\spec_\A(x)$, entonces existe $(\lambda1-x)^{-1}$. Aplicando $\psi$,
\[
 \psi(\lambda1-x)\,\psi((\lambda1-x)^{-1})
 =\psi((\lambda1-x)^{-1})\,\psi(\lambda1-x)=1_\B,
\]
así que $\lambda1-\psi(x)$ es invertible. Esto prueba
$\spec_\B(\psi(x))\subset\spec_\A(x)$. Tomando máximos de $|\lambda|$ se obtiene
$r(\psi(x))\le r(x)$.
\end{proof}

\proofdiagram
{$\psi:\A\to\B$ homomorfismo unital}
{Los homomorfismos unitales mandan inversos en inversos}
{$\lambda1-\psi(x)=\psi(\lambda1-x)$}
{$\spec(\psi(x))\subset\spec(x)$ y $r(\psi(x))\le r(x)$}

\begin{cthejemplo}{(Inclusión espectral estricta en un cociente)}
Para mostrar que la inclusión de T1.8 puede ser estricta, tomemos
$\A=\mathbb C\oplus\mathbb C$ con norma supremo y $J=\{0\}\oplus\mathbb C$.
La clase de $a=(1,2)$ en $\A/J\cong\mathbb C$ es $1$.
Como $\lambda1-a=(\lambda-1,\lambda-2)$ es invertible exactamente cuando
$\lambda\notin\{1,2\}$,
\[
 \spec_{\A/J}(a+J)=\{1\}\subsetneq\{1,2\}=\spec_\A(a).
\]
El cociente puede perder puntos espectrales, de acuerdo con T1.8.
\end{cthejemplo}

\noindent\textbf{Dependencias.} Un homomorfismo unital conserva el inverso de $\lambda1-x$ cuando existe; la definición de espectro convierte esta observación algebraica en una inclusión.

\noindent\textbf{Qué prepara este ejercicio.} En la Sección 2 se aplica la inclusión a $a^*a$ para probar que los $*$-homomorfismos unitarios son contractivos.

\section{Problema 9: una subálgebra conmutativa que preserva el espectro}

Sea $\A$ un álgebra de Banach unital y $S\subset\A$ una familia que conmuta dos a dos. Para la
versión involutiva, suponemos además que $\A$ es involutiva y los elementos de $S$
$*$-conmutan dos a dos.

\noindent\textbf{Enunciado.} Prueba que existe una subálgebra de Banach unital y conmutativa $C\subset\A$ que contiene $S$ y satisface $\spec_C(s)=\spec_\A(s)$ para todo $s\in S$. En el caso involutivo con $*$-conmutación dos a dos, exige además que $C$ sea involutiva.

\begin{proof}[Solución]
Sea $\mathscr F$ la familia de subálgebras de Banach conmutativas unitales $C\subset\A$ que contienen
a $S$, ordenada por inclusión. Es no vacía: contiene la clausura de la álgebra unital generada por
$S$. Si $\{C_i\}$ es una cadena, la clausura de $\bigcup_i C_i$ vuelve a ser una subálgebra de
Banach conmutativa unital que contiene $S$. Por el lema de Zorn existe un elemento maximal $C$.

Afirmamos que $C$ es inversamente cerrada en $\A$. Si $c\in C$ es invertible en $\A$, entonces
de $cd=dc$ para $d\in C$ se deduce $c^{-1}d=dc^{-1}$. Por tanto la subálgebra cerrada unital
generada por $C\cup\{c^{-1}\}$ sigue siendo conmutativa. La maximalidad obliga a
$c^{-1}\in C$. En consecuencia, para $s\in S$,
\[
\lambda1-s\in\GL(\A)\Longleftrightarrow\lambda1-s\in\GL(C),
\]
y los espectros coinciden.

En el caso involutivo, partimos de la familia de subálgebras de Banach involutivas, conmutativas y
unitales que contienen $S$. La hipótesis de $*$-conmutación garantiza que la familia es no vacía por
el Problema 3(d), aplicado a $S\cup\{1\}$. Si $c\in C$ es invertible en $\A$, tanto $c^{-1}$ como $(c^{-1})^*$ conmutan con
$C$ y entre sí: $c,c^*\in C$ conmutan y sus inversos también. Se pueden añadir conservando
conmutatividad e involución. La maximalidad vuelve a
forzar $c^{-1}\in C$.
\end{proof}

\proofdiagrambranch
{$S$ conmuta dos a dos}
{Familia de subálgebras cerradas, conmutativas, unitales que contienen $S$}
{Zorn produce una maximal; los inversos de sus elementos también conmutan con ella}
{Maximalidad fuerza que los inversos pertenezcan a la subálgebra}
{$\spec_\A(s)=\spec_C(s)$ para $s\in S$}

\begin{cthejemplo}{(Una subálgebra espectralmente fiel de matrices)}
En $M_2(\mathbb C)$, el autoadjunto $A=\operatorname{diag}(1,2)$ permite construir
explícitamente la subálgebra de T1.9.
Tomemos $C$ como todas las matrices diagonales. El inverso de
$\lambda I-A=\operatorname{diag}(\lambda-1,\lambda-2)$, cuando existe, sigue en $C$:
\[
 (\lambda I-A)^{-1}=\operatorname{diag}((\lambda-1)^{-1},(\lambda-2)^{-1}).
\]
Por ello $\spec_C(A)=\{1,2\}=\spec_{M_2(\mathbb C)}(A)$; el espectro se conserva porque
los inversos necesarios pertenecen a $C$.
\end{cthejemplo}

\noindent\textbf{Dependencias.} Se usan T1.3 para generar una subálgebra conmutativa, T1.2 para conmutar con inversos y el lema de Zorn para escoger una subálgebra maximal. La versión involutiva exige $*$-conmutación.

\noindent\textbf{Qué prepara este ejercicio.} La Sección 2 emplea esta reducción para estudiar un autoadjunto mediante caracteres sin cambiar su espectro. En el entorno $C^*$, la permanencia espectral ofrece después una vía más directa para $C^*(a,1)$.

\pagebreak
\section{Problema 10: espectro de una función continua}

Sea $X$ compacto Hausdorff y $f\in C(X)$. Después compararemos con $C_b(X)$ para un espacio
localmente compacto que no sea compacto.

\noindent\textbf{Enunciado.} Prueba que $\spec_{C(X)}(f)=f(X)$ para cada $f\in C(X)$ y decide, mediante un ejemplo, si la igualdad sigue siendo válida en $C_b(X)$ para $X$ sólo localmente compacto.

\begin{proof}[Solución]
Si $\lambda\notin f(X)$, la función $\lambda-f$ no se anula. Como $X$ es compacto,
$|\lambda-f|$ tiene mínimo positivo; por tanto
$(\lambda-f)^{-1}\in C(X)$ y $\lambda\notin\spec(f)$. Si $\lambda=f(x_0)$, entonces
$\lambda-f$ se anula en $x_0$ y no puede ser invertible. Por tanto $\spec(f)=f(X)$.

Si $X$ sólo es localmente compacto y usamos $C_b(X)$, el enunciado puede fallar porque una función
no nula puede acercarse a cero sin alcanzar cero. Por ejemplo, en $X=(0,1)$, $f(x)=x$ tiene
$f(X)=(0,1)$, pero
\[
 \spec_{C_b((0,1))}(f)=[0,1].
\]
En efecto, para $\lambda\in(0,1)$ la función $\lambda-f$ se anula; para $\lambda=0,1$,
sus recíprocos no son acotados. Si $\lambda\notin[0,1]$, en cambio,
$|\lambda-f(x)|\ge\operatorname{dist}(\lambda,[0,1])>0$ para todo $x$, de modo que el
recíproco pertenece a $C_b((0,1))$. Esto prueba la igualdad anunciada.
\end{proof}

\proofdiagram
{$f\in C(X)$, $X$ compacto}
{Invertibilidad puntual en una álgebra de funciones}
{$\lambda1-f$ invertible $\Longleftrightarrow$ $\lambda-f$ no se anula}
{$\spec(f)=f(X)$}

\begin{cthejemplo}{(Un espectro con forma de circunferencia)}
Para leer geométricamente T1.10, consideremos $f(z)=2+z$ en $C(\mathbb T)$.
La función $\lambda-f$ se anula en $\mathbb T$ exactamente cuando
$\lambda=2+z$ para algún $|z|=1$; fuera de esos valores tiene recíproco continuo.
Por tanto
\[
 \spec_{C(\mathbb T)}(f)=\{\lambda\in\mathbb C:|\lambda-2|=1\}.
\]
El espectro es la imagen de $f$, una circunferencia de centro $2$ y radio $1$.
\end{cthejemplo}

\noindent\textbf{Dependencias.} En $C(X)$, $\lambda1-f$ es invertible exactamente cuando $\lambda-f$ no se anula; la compacidad de $X$ asegura que su recíproco es continuo y acotado.

\noindent\textbf{Qué prepara este ejercicio.} Identifica el espectro con un rango de valores en la representación de Gelfand--Naimark y es el paso concreto que usa T2.4 para demostrar el mapeo espectral.

\pagebreak
\section{Problema 11: $\mathcal B(H)$ y el shift unilateral}

Sea $H$ un espacio de Hilbert y $T\in\mathcal B(H)$. Los incisos (b)--(d) cambian las hipótesis
sobre $H$ según se indica en cada uno.

\subsection*{11(a) $\mathcal B(H)$ es una $C^*$-álgebra}

\noindent\textbf{Enunciado.} Prueba que $\mathcal B(H)$, con la norma operatorial y la involución $T\mapsto T^*$, es una $C^*$-álgebra.

\begin{proof}[Solución]
La Parte 1 ya estableció que los operadores acotados sobre un espacio de Banach forman una álgebra
de Banach; por tanto $\mathcal B(H)$ es completa. Las reglas del adjunto dan una involución.
La norma del adjunto coincide con la de $T$: por la caracterización de la norma de un vector de
Hilbert mediante productos interiores,
\[
 \|T^*\|=\sup_{\|\eta\|=\|\xi\|=1}|\langle T^*\eta,\xi\rangle|
 =\sup_{\|\eta\|=\|\xi\|=1}|\langle\eta,T\xi\rangle|=\|T\|.
\]
La submultiplicatividad da entonces $\|T^*T\|\le\|T^*\|\|T\|=\|T\|^2$.
Para la desigualdad inversa, si $\|\xi\|=1$,
\[
\|T\xi\|^2=\langle T^*T\xi,\xi\rangle\le\|T^*T\|.
\]
Tomando supremo en $\xi$ se obtiene $\|T\|^2\le\|T^*T\|$.
\end{proof}

\proofdiagram
{$T\in\mathcal B(H)$}
{Propiedades del adjunto}
{$\|T^*T\|\le\|T\|^2$ y $\|T\xi\|^2=\langle T^*T\xi,\xi\rangle$}
{$\|T^*T\|=\|T\|^2$}

\begin{cthejemplo}{(La identidad $C^*$ en una matriz diagonal)}
En $H=\mathbb C^2$ queremos comprobar la igualdad de normas que distingue a
$\mathcal B(H)$ como $C^*$-álgebra.
Para $T=\operatorname{diag}(2,i)$, se tiene
$T^*T=\operatorname{diag}(4,1)$, luego $\|T^*T\|=4$ y $\|T\|=2$.
Así $\|T^*T\|=\|T\|^2$, la identidad verificada de modo general en T1.11(a).
\end{cthejemplo}

\noindent\textbf{Dependencias.} La Parte 1 presenta $\mathcal B(H)$ como álgebra de Banach; las propiedades del adjunto y la norma operatorial permiten comprobar la identidad $C^*$.

\noindent\textbf{Qué prepara este ejercicio.} Suministra el modelo concreto de la Sección 2 y el espacio de operadores en el que T2.1 representa isométricamente a $C_b(X)$.

\subsection*{11(b) En dimensión finita, espectro = valores propios}

\noindent\textbf{Enunciado.} Si $\dim H<\infty$, prueba que $\spec(T)$ coincide con el conjunto de valores propios de $T$.

\begin{proof}[Solución]
En dimensión finita un operador lineal es invertible si y sólo si su núcleo es cero. Así
$\lambda I-T$ no es invertible exactamente cuando existe $\xi\neq0$ con
$(\lambda I-T)\xi=0$, es decir, $T\xi=\lambda\xi$.
\end{proof}

\proofdiagram
{$\dim H<\infty$}
{En dimensión finita: inyectivo $\Longleftrightarrow$ sobreyectivo}
{$\lambda I-T$ no invertible $\Longleftrightarrow\Ker(\lambda I-T)\neq0$}
{$\lambda\in\spec(T)\Longleftrightarrow\lambda$ valor propio}

\begin{cthejemplo}{(Valores propios de una matriz triangular)}
En dimensión dos, el determinante permite contrastar directamente el espectro y los
vectores propios de T1.11(b).
Para $T=\begin{pmatrix}1&1\\0&2\end{pmatrix}$,
$\det(\lambda I-T)=(\lambda-1)(\lambda-2)$.
Además $Te_1=e_1$ y $T(1,1)=(2,2)$, así que $1$ y $2$ son valores propios.
No hay otros puntos espectrales: en dimensión finita la singularidad de $\lambda I-T$
coincide con la existencia de un vector propio.
\end{cthejemplo}

\noindent\textbf{Dependencias.} En dimensión finita, inyectividad y sobreyectividad de $\lambda I-T$ son equivalentes por rango y nulidad. Se aplica la definición de espectro de la Parte 1.

\noindent\textbf{Qué prepara este ejercicio.} Permite interpretar los ejemplos matriciales de la Sección 2 como cálculos de espectro mediante valores propios; T1.11(c)--(d) muestran qué cambia en dimensión infinita.

\subsection*{11(c) Vectores propios aproximados producen puntos espectrales}

\noindent\textbf{Enunciado.} Si $\lambda\in\mathbb C$ y existen vectores unitarios $\xi_n\in H$ tales que $\|(\lambda I-T)\xi_n\|\to0$, prueba que $\lambda\in\spec(T)$.

\begin{proof}[Solución]
Si $\lambda\notin\spec(T)$, entonces $R=(\lambda I-T)^{-1}$ es acotado. Por tanto
\[
1=\|\xi_n\|=\|R(\lambda I-T)\xi_n\|
\le\|R\|\,\|(\lambda I-T)\xi_n\|\to0,
\]
contradicción. Luego $\lambda\in\spec(T)$.
\end{proof}

\proofdiagram
{$\|\xi_n\|=1$, $\|(\lambda I-T)\xi_n\|\to0$}
{Suponer $\lambda I-T$ invertible}
{$\xi_n=(\lambda I-T)^{-1}(\lambda I-T)\xi_n$}
{$1\le\|(\lambda I-T)^{-1}\|\|(\lambda I-T)\xi_n\|\to0$, contradicción}

\begin{cthejemplo}{(Espectro aproximado sin vector propio)}
En $\ell^2(\mathbb N)$ queremos ver que un punto espectral aproximado puede carecer de
vector propio. Definamos $T(\xi_n)=(\xi_n/n)$ y tomemos $\lambda=0$.
Cada $e_n$ es unitario y $\|Te_n\|=1/n\to0$, pero $T\xi=0$ implica
$\xi_n/n=0$ para todo $n$, es decir, $\xi=0$.
Por T1.11(c), $0\in\spec(T)$; el cálculo distingue espectro aproximado de espectro puntual.
\end{cthejemplo}

\noindent\textbf{Dependencias.} Si $\lambda I-T$ es invertible, su inverso es acotado. La estimación de la norma de la composición contradice la existencia de vectores unitarios casi anulados.

\noindent\textbf{Qué prepara este ejercicio.} Define el espectro puntual aproximado y deja preparada la comparación con el shift unilateral de T1.11(d), cuyo cero espectral no es un valor propio aproximado.

\subsection*{11(d) El shift unilateral muestra la diferencia}

\noindent\textbf{Enunciado.} Sea $S(x_1,x_2,\ldots)=(0,x_1,x_2,\ldots)$ en
$\ell^2(\mathbb N)$. Prueba que $0\in\spec(S)$ y que no existen vectores unitarios
$\xi_n$ con $\|S\xi_n\|\to0$.

\begin{proof}[Solución]
$S$ no es sobreyectivo: ningún vector con primera coordenada no nula está en su rango. Así $S$ no
es invertible y $0\in\spec(S)$. Pero
\[
 \|S\xi\|^2=0+\sum_{n\ge1}|\xi_n|^2=\|\xi\|^2.
\]
Por tanto, para todo vector unitario $\xi$, $\|S\xi\|=1$; no puede existir una sucesión unitaria con
$\|S\xi_n\|\to0$.
\end{proof}

\proofdiagram
{$S(x_1,x_2,\dots)=(0,x_1,x_2,\dots)$}
{$S$ es isometría pero no es sobreyectivo}
{$0\in\spec(S)$; sin embargo $\|S\xi\|=\|\xi\|$}
{No hay vectores unitarios con $\|S\xi_n\|\to0$}

\begin{cthejemplo}{(Una isometría que no es sobreyectiva)}
Para distinguir conservación de norma e invertibilidad en T1.11(d), tomemos el primer
vector de la base de $\ell^2(\mathbb N)$.
El shift satisface $Se_1=e_2$ y $\|Se_1\|=\|e_1\|=1$, pero ningún vector puede
transformarse en $e_1$, porque la primera coordenada de $S\xi$ siempre es cero.
La isometría no es sobreyectiva; por eso $0$ pertenece a su espectro sin haber vectores
unitarios casi anulados.
\end{cthejemplo}

\noindent\textbf{Dependencias.} El shift unilateral es una isometría, pero su rango omite vectores con primera coordenada no nula. La definición de espectro detecta esta falta de sobreyectividad.

\noindent\textbf{Qué prepara este ejercicio.} T2.10 determinará todo el disco espectral usando vectores propios del adjunto. La isometría sin sobreyectividad explica por qué $0$ puede ser espectral sin ser valor propio aproximado.

\begin{investigacion}[title={Dos shifts en espacios distintos}]
Para reconocer el alcance de T1.11(d), comparemos el shift $S$ de $\ell^2(\mathbb N)$ con
$V_\alpha$ en la formulación de Haseman. El primero satisface $Se_1=e_2$; el segundo
compone funciones de frontera, $(V_\alpha f)(t)=f(\alpha(t))$
\cite[pp.~463--464, ecs.~(1.1)--(1.2)]{Karlovich2017}.
Ambos desplazan datos, pero actúan en espacios y geometrías diferentes. T1.11(d)
determina $0\in\spec(S)$; no determina por analogía ningún punto de $\spec(V_\alpha)$.
\end{investigacion}

\section{Problema 12: el álgebra del disco y la no permanencia espectral general}

Sea $\mathbb D=\{z:|z|<1\}$, $\mathbb T=\{z:|z|=1\}$ y
\[
 A(\mathbb D)=\{f\in C(\overline{\mathbb D}): f\text{ es holomorfa en }\mathbb D\}.
\]

\subsection*{12(a) $A(\mathbb D)$ es una álgebra de Banach}

\noindent\textbf{Enunciado.} Prueba que $A(\mathbb D)$ es una álgebra de Banach unital con la norma del supremo en $\overline{\mathbb D}$.

\begin{proof}[Solución]
Sea $(f_n)$ de Cauchy en norma supremo. Converge uniformemente en $\overline{\mathbb D}$ a una
función continua $f$. En cada compacto $K\subset\mathbb D$, la convergencia es uniforme; por el
teorema de Weierstrass, $f$ es holomorfa en $\mathbb D$. Así $f\in A(\mathbb D)$ y el espacio es
completo. Las sumas y productos de funciones holomorfas son holomorfos, y la función constante
$1$ pertenece a $A(\mathbb D)$; las operaciones y la norma se heredan de
$C(\overline{\mathbb D})$.
\end{proof}

\proofdiagram
{Sucesión de Cauchy en norma supremo}
{$C(\overline{\mathbb D})$ es Banach; teorema de Weierstrass para límites holomorfos}
{El límite uniforme es continuo y el límite en compactos interiores es holomorfo}
{$A(\mathbb D)$ es cerrado, unital y Banach}

\begin{cthejemplo}{(Aproximación polinómica en el álgebra del disco)}
En $A(\mathbb D)$ interesa aproximar funciones holomorfas sin salir de la norma uniforme.
Para $f\in A(\mathbb D)$, las dilataciones $f_r(z)=f(rz)$ convergen uniformemente a $f$
en $\overline{\mathbb D}$ cuando $r\uparrow1$, por continuidad uniforme de $f$.
Cada $f_r$ es holomorfa en un disco de radio $1/r>1$, así que sus polinomios de Taylor
convergen uniformemente en $\overline{\mathbb D}$.
Esto explica por qué los polinomios forman una subálgebra densa del Banach de T1.12(a).
\end{cthejemplo}

\noindent\textbf{Dependencias.} La completitud de $C(\overline{\mathbb D})$ y el teorema de Weierstrass sobre límites uniformes locales de funciones holomorfas garantizan que $A(\mathbb D)$ es cerrado.

\noindent\textbf{Qué prepara este ejercicio.} Hace posible tratar la restricción de frontera como una inclusión de álgebras de Banach en T1.12(b) y comparar espectros en T1.12(c).

\subsection*{12(b) La restricción a la frontera es isométrica}

\noindent\textbf{Enunciado.} Prueba que la restricción $f\mapsto f|_{\mathbb T}$ es un homomorfismo unital e isométrico de $A(\mathbb D)$ en $C(\mathbb T)$.

\begin{proof}[Solución]
Por el principio del máximo,
\[
 \|f\|_{C(\overline{\mathbb D})}=\max_{|z|=1}|f(z)|=\|f|_{\mathbb T}\|_{C(\mathbb T)}.
\]
Así la restricción es isométrica e inyectiva. Para $f,g\in A(\mathbb D)$ y $z\in\mathbb T$,
\[
 (f+g)|_{\mathbb T}(z)=f(z)+g(z),\qquad
 (fg)|_{\mathbb T}(z)=f(z)g(z),\qquad
 1|_{\mathbb T}(z)=1.
\]
Por tanto es un homomorfismo unital de álgebras de Banach.
\end{proof}

\proofdiagram
{$f\in A(\mathbb D)$}
{Principio del máximo para funciones holomorfas}
{$\max_{\overline{\mathbb D}}|f|=\max_{\mathbb T}|f|$}
{$f\mapsto f|_{\mathbb T}$ es un encaje isométrico}

\begin{cthejemplo}{(La norma del disco se alcanza en la frontera)}
En $A(\mathbb D)$ comprobamos T1.12(b) con $f(z)=1+z^2$.
Para $|z|\le1$, $|f(z)|\le2$, y la igualdad se alcanza en $z=\pm1\in\mathbb T$.
Así $\|f\|_{C(\overline{\mathbb D})}=\|f|_{\mathbb T}\|_{C(\mathbb T)}=2$:
la frontera ya determina la norma de esta función.
\end{cthejemplo}

\noindent\textbf{Dependencias.} El principio del máximo identifica la norma en el disco cerrado con la norma sobre la frontera; por ello la restricción es isométrica e inyectiva.

\noindent\textbf{Qué prepara este ejercicio.} Proporciona una inclusión isométrica y unital en $C(\mathbb T)$ para T1.12(c), donde la misma función tendrá distintos inversos disponibles según el álgebra.

\subsection*{12(c) Un mismo elemento cambia de espectro al cambiar de álgebra}

\noindent\textbf{Enunciado.} Para $h(z)=z$, prueba $\spec_{A(\mathbb D)}(h)=\overline{\mathbb D}$ y $\spec_{C(\mathbb T)}(h|_{\mathbb T})=\mathbb T$.

\begin{proof}[Solución]
Si $|\lambda|>1$, la función $(\lambda-z)^{-1}$ es continua en $\overline{\mathbb D}$ y holomorfa
en $\mathbb D$, así que $\lambda-h$ es invertible en $A(\mathbb D)$. Si $|\lambda|<1$,
$\lambda-z$ se anula en $z=\lambda$; si $|\lambda|=1$, cualquier supuesto inverso coincidiría en
$\mathbb D$ con $(\lambda-z)^{-1}$, que se hace no acotado al acercarse a $\lambda$. Por tanto
\[
 \spec_{A(\mathbb D)}(h)=\overline{\mathbb D}.
\]

En $C(\mathbb T)$, $\lambda-h$ es invertible exactamente cuando no se anula sobre $\mathbb T$; por
tanto
\[
 \spec_{C(\mathbb T)}(h)=\mathbb T.
\]
\end{proof}

\proofdiagrambranch
{$h(z)=z$ en $A(\mathbb D)$}
{$h|_{\mathbb T}$ en $C(\mathbb T)$}
{Invertibilidad de $\lambda-h$ depende de si el recíproco pertenece al álgebra considerada}
{En $A(\mathbb D)$ falla para $\lambda\in\overline{\mathbb D}$; en $C(\mathbb T)$ sólo falla en $\mathbb T$}
{$\spec_{A(\mathbb D)}(h)=\overline{\mathbb D}$, $\spec_{C(\mathbb T)}(h)=\mathbb T$}

\begin{cthejemplo}{(El inverso depende del álgebra ambiente)}
Para comparar los dos espectros de T1.12(c), tomemos $h(z)=z$ en $A(\mathbb D)$ y
su restricción en $C(\mathbb T)$.
En el primer álgebra, $0-h$ no tiene inverso holomorfo porque $1/z$ tiene un polo en
$0$; en el segundo, $h^{-1}(z)=\overline z$ es continuo sobre $\mathbb T$.
Así $0\in\spec_{A(\mathbb D)}(h)$ y $0\notin\spec_{C(\mathbb T)}(h|_{\mathbb T})$.
La misma función de frontera dispone de inverso sólo en el álgebra mayor. Además,
$\overline z$ no proviene de $A(\mathbb D)$: si una función $g$ tuviera esos valores de
frontera, $zg(z)=1$ en $\mathbb T$ y el principio del máximo daría $zg(z)=1$ en todo
$\mathbb D$, contradicción en $z=0$. Por eso la inclusión no es de $C^*$-subálgebras.
\end{cthejemplo}

\noindent\textbf{Dependencias.} Se usa la definición de espectro y la inclusión de T1.12(b); T1.10 calcula el espectro en $C(\mathbb T)$. En $A(\mathbb D)$ el recíproco de $\lambda-z$ debe además ser holomorfo en el disco.

\noindent\textbf{Qué prepara este ejercicio.} Es el contraejemplo que motiva la hipótesis de $C^*$-subálgebra en la permanencia espectral de la Sección 2: una subálgebra de Banach cerrada y unital puede cambiar el espectro.

\chapter{Tarea 2 resuelta: consecuencias del cálculo funcional}

\begin{context}
Al terminar esta tarea, podrás verificar una representación concreta de $C_b(X)$, reconstruir
caracterizaciones espectrales de elementos normales y distinguir la convergencia en norma de la
convergencia fuerte de operadores. Las dependencias indicadas después de cada solución permiten
comprobar qué hipótesis hacen válido cada paso.
\end{context}

\section{Problema 1: $C_b(X)$ como $C^*$-álgebra concreta}

Sea $X$ localmente compacto y Hausdorff.


Sea $H=\ell^2(X)$ y, para $f\in C_b(X)$, definamos
\[
 (M_f\xi)(x)=f(x)\xi(x).
\]
Para un conjunto de índices arbitrario, la suma en $\|\xi\|_2^2$ se entiende como
\[
 \|\xi\|_2^2=\sup_{F\subset X\text{ finito}}\sum_{x\in F}|\xi(x)|^2.
\]
La estimación inferior se hará cuando $X\ne\varnothing$.

\noindent\textbf{Enunciado.} Prueba que $M:C_b(X)\to\mathcal B(\ell^2(X))$, $f\mapsto M_f$, es un $*$-homomorfismo isométrico.

\begin{proof}[Solución]
Si $X=\varnothing$, $C_b(X)$ contiene sólo la función vacía y $\ell^2(X)=\{0\}$; por tanto
$\mathcal B(\ell^2(X))=\{0\}$ y la única aplicación $f\mapsto M_f$ preserva operaciones y
normas. Supongamos en adelante $X\ne\varnothing$.

Para $\xi\in\ell^2(X)$,
\[
\|M_f\xi\|_2^2
=\sum_{x\in X}|f(x)|^2|\xi(x)|^2
\le\|f\|_\infty^2\|\xi\|_2^2,
\]
de modo que $\|M_f\|\le\|f\|_\infty$. Para la desigualdad inversa, dado
$\varepsilon>0$ elegimos $x_0$ con $|f(x_0)|>\|f\|_\infty-\varepsilon$ y tomamos el vector
$e_{x_0}$ concentrado en $x_0$. Entonces
$\|M_fe_{x_0}\|=|f(x_0)|$, de donde $\|M_f\|=\|f\|_\infty$.

Además,
\[
M_{f+g}=M_f+M_g,
\qquad M_{fg}=M_fM_g,
\qquad M_1=I.
\]
Finalmente,
\[
\langle M_f\xi,\eta\rangle
=\sum_x f(x)\xi(x)\overline{\eta(x)}
=\langle \xi,M_{\bar f}\eta\rangle,
\]
así que $M_f^*=M_{\bar f}$. Por tanto $M$ es un $*$-homomorfismo isométrico.
\end{proof}

\proofdiagram
{$f\in C_b(X)$}
{Norma en $\ell^2(X)$; conjugación puntual}
{$\|M_f\|=\|f\|_\infty$, $M_{fg}=M_fM_g$, $M_{\bar f}=M_f^*$}
{$f\mapsto M_f$ es un $*$-homomorfismo isométrico}

\begin{cthejemplo}{(Un operador diagonal en el espacio de sucesiones)}
Para ver la representación de T2.1 sobre un espacio conocido, tomemos $X=\mathbb N$ y
$f(n)=1/(n+1)$.
En $H=\ell^2(\mathbb N)$ se tiene $M_fe_n=(n+1)^{-1}e_n$; además,
\[
 \|M_f\|=\sup_{n\ge1}\|M_fe_n\|=\frac12=\|f\|_\infty.
\]
La función acotada aparece como operador diagonal, y su norma no cambia bajo el
$*$-homomorfismo de T2.1.
\end{cthejemplo}

\noindent\textbf{Dependencias.} Se usan la definición de $C_b(X)$ de la Parte 1, la norma de $\ell^2(X)$ y la caracterización del adjunto en Hilbert. Aquí $X$ sirve como conjunto de índices, sin usar su topología para definir $\ell^2(X)$.

\noindent\textbf{Qué prepara este ejercicio.} Exhibe un $*$-monomorfismo isométrico hacia $\mathcal B(H)$, modelo concreto de los homomorfismos y representaciones que estudia la Sección 2.

\section{Problema 2: espectros de idempotentes, productos y unitarios}

Sea $\A$ un álgebra de Banach unital. En la conclusión sobre unitarios del inciso (c), $\A$
es además una $C^*$-álgebra.

\subsection*{2(a) Espectro de un idempotente}

\noindent\textbf{Enunciado.} Si $e\in\A$ y $e^2=e$, prueba $\spec(e)\subset\{0,1\}$.

\begin{proof}[Solución]
Si $\lambda\neq0,1$, pongamos
\[
 b=\frac1\lambda 1+\frac{1}{\lambda(\lambda-1)}e.
\]
Como $e(1-e)=(1-e)e=0$ y $e^2=e$, podemos escribir
\[
 \lambda1-e=\lambda(1-e)+(\lambda-1)e,
 \qquad b=\lambda^{-1}(1-e)+(\lambda-1)^{-1}e.
\]
Al multiplicar, los términos cruzados se anulan y queda
$(\lambda1-e)b=b(\lambda1-e)=(1-e)+e=1$. Así todo
$\lambda\notin\{0,1\}$ está en el resolvente y
$\spec(e)\subset\{0,1\}$.
\end{proof}

\pagebreak
\proofdiagram
{$e^2=e$}
{Resolver algebraicamente $(\lambda1-e)^{-1}$}
{Para $\lambda\notin\{0,1\}$ construir $\lambda^{-1}1+[\lambda(\lambda-1)]^{-1}e$}
{$\spec(e)\subset\{0,1\}$}

\begin{cthejemplo}{(Espectro de una proyección matricial)}
En $M_2(\mathbb C)$ queremos identificar los dos puntos permitidos por T2.2(a).
Para $P=\operatorname{diag}(1,0)$, se cumple $P^2=P$ y
$\lambda I-P=\operatorname{diag}(\lambda-1,\lambda)$.
Esta matriz es invertible exactamente cuando $\lambda\ne0,1$; por ello
$\spec(P)=\{0,1\}$. Si $P$ se sustituye por $0$ o $I$, el mismo cálculo deja sólo
$\{0\}$ o $\{1\}$, respectivamente.
\end{cthejemplo}

\noindent\textbf{Dependencias.} La ecuación $e^2=e$ permite verificar un inverso explícito de $\lambda1-e$ para $\lambda\ne0,1$; se aplica la definición de espectro de la Parte 1.

\noindent\textbf{Qué prepara este ejercicio.} Da la implicación directa de la caracterización espectral de proyecciones en T2.6(c) y en la Proposición 2.28.

\subsection*{2(b) Radio espectral de un producto conmutativo}

\noindent\textbf{Enunciado.} Si $x,y\in\A$ conmutan, prueba $r(xy)\le r(x)r(y)$.

\begin{proof}[Solución]
La conmutatividad permite escribir $(xy)^n=x^ny^n$. Entonces
\[
\|(xy)^n\|^{1/n}
\le\|x^n\|^{1/n}\|y^n\|^{1/n}.
\]
Al pasar al límite y usar la fórmula del radio espectral,
$r(xy)\le r(x)r(y)$.
\end{proof}

\proofdiagram
{$xy=yx$}
{Fórmula del radio espectral}
{$(xy)^n=x^ny^n$ y $\|(xy)^n\|\le\|x^n\|\|y^n\|$}
{$r(xy)\le r(x)r(y)$}

\begin{cthejemplo}{(La desigualdad del radio puede ser estricta)}
En $C([0,1])$ podemos comprobar cuánto puede mejorar la cota de T2.2(b).
Tomemos $f(t)=t$ y $g(t)=1-t$; conmutan porque el producto es puntual.
Por T1.10, $r(f)=\|f\|_\infty=1$ y $r(g)=1$, mientras que
\[
 r(fg)=\max_{0\le t\le1}t(1-t)=\frac14<r(f)r(g)=1.
\]
La desigualdad es estricta porque los máximos de $f$ y $g$ ocurren en extremos distintos.
\end{cthejemplo}

\noindent\textbf{Dependencias.} La conmutación permite $(xy)^n=x^ny^n$ y la fórmula del radio espectral de la Parte 1 convierte la estimación de normas en la desigualdad buscada.

\noindent\textbf{Qué prepara este ejercicio.} Permite estimar productos dentro de subálgebras conmutativas como $C^*(x,1)$. La clasificación de T2.6 usa cálculo funcional y no depende de esta desigualdad.

\subsection*{2(c) Espectro de un invertible con norma controlada}

\noindent\textbf{Enunciado.} Si $z\in\GL(\A)$ y $\|z\|,\|z^{-1}\|\le1$, prueba $\spec(z)\subset\mathbb T$; deduce la misma inclusión para todo unitario de una $C^*$-álgebra unital.

\begin{proof}[Solución]
Si $\lambda\in\spec(z)$, entonces $|\lambda|\le r(z)\le\|z\|\le1$. Como $z$ es invertible,
$0\notin\spec(z)$. Además
\[
 \lambda\in\spec(z)\quad\Longleftrightarrow\quad
 \lambda^{-1}\in\spec(z^{-1}),
\]
porque
$\lambda1-z=-\lambda z(\lambda^{-1}1-z^{-1})$. Por tanto
$|\lambda|^{-1}\le r(z^{-1})\le1$, es decir $|\lambda|\ge1$. Luego $|\lambda|=1$.

Si $u$ es unitario en una $C^*$-álgebra, $u^{-1}=u^*$ y
$\|u\|=\|u^*\|=1$, así que el resultado anterior da
$\spec(u)\subset\mathbb T$.
\end{proof}

\proofdiagrambranch
{$z\in\GL(\A)$ y $\|z\|\le1$}
{$\|z^{-1}\|\le1$}
{$\lambda\in\spec(z)\Rightarrow |\lambda|\le1$ y $\lambda^{-1}\in\spec(z^{-1})$}
{$|\lambda|^{-1}\le1$}
{$|\lambda|=1$, luego $\spec(z)\subset\mathbb T$}

\begin{cthejemplo}{(Espectro de una matriz unitaria diagonal)}
Para ver la conclusión unitaria de T2.2(c) con dos valores espectrales, tomemos
$U=\operatorname{diag}(1,i)$ en $M_2(\mathbb C)$.
Se tiene $U^*U=I$, $U^{-1}=U^*$ y $\|U\|=\|U^{-1}\|=1$.
Además $\det(\lambda I-U)=(\lambda-1)(\lambda-i)$, de modo que
$\spec(U)=\{1,i\}\subset\mathbb T$.
La cota de las normas fuerza que ambos puntos queden en el círculo unitario.
\end{cthejemplo}

\noindent\textbf{Dependencias.} Se usan la cota $r(z)\le\|z\|$, la relación entre los espectros de un invertible y su inverso, y ambas hipótesis $\|z\|,\|z^{-1}\|\le1$.

\noindent\textbf{Qué prepara este ejercicio.} Da la implicación unitaria de T2.6(b) y de la Proposición 2.28; para un unitario $u$, la identidad $C^*$ asegura $\|u\|=\|u^{-1}\|=1$.

\begin{investigacion}[title={Rotaciones unitarias en $L^2(\mathbb T)$}]
Para aplicar T2.2(c) en el modelo Hilbertiano de \cite{BFK2006}, fijemos $z\in\mathbb T$ y
definamos $(U_z\varphi)(t)=\varphi(zt)$ sobre $L^2(\mathbb T)$.
La invariancia de la medida de Haar bajo rotaciones da $\|U_z\varphi\|_2=\|\varphi\|_2$;
además $U_{z^{-1}}U_z=U_zU_{z^{-1}}=I$. Así $U_z$ es unitario y
$\|U_z\|=\|U_z^{-1}\|=1$.
Por T2.2(c), $\spec(U_z)\subset\mathbb T$. La conclusión usa esas propiedades
verificadas en $L^2(\mathbb T)$.
\end{investigacion}

\pagebreak
\section{Problema 3: productos de $C^*$-álgebras y el ideal $c_0$}

Sea $I$ un conjunto de índices y $(A_i)_{i\in I}$ una familia de $C^*$-álgebras. En el inciso
(b) se toma $I=\mathbb N$.

\subsection*{3(a) Producto acotado de una familia de $C^*$-álgebras}

\noindent\textbf{Enunciado.} Prueba que el producto de familias $(a_i)_{i\in I}$ con $a_i\in A_i$ y $\sup_i\|a_i\|<\infty$ es una $C^*$-álgebra con operaciones puntuales y norma supremo.

\begin{proof}[Solución]
Sea
\[
\prod_{i\in I}^{\ell^\infty} A_i
=\{(a_i): a_i\in A_i,\ \sup_i\|a_i\|<\infty\}.
\]
Las operaciones puntuales preservan la acotación uniforme; en particular,
\[
 \sup_i\|a_i+b_i\|\le\sup_i\|a_i\|+\sup_i\|b_i\|,
 \qquad
 \sup_i\|a_ib_i\|\le\bigl(\sup_i\|a_i\|\bigr)
                        \bigl(\sup_i\|b_i\|\bigr).
\]
Con la norma $\|(a_i)\|=\sup_i\|a_i\|$ es, por tanto, una álgebra normada. Si
$(a^{(n)})$ es de Cauchy, cada coordenada $a_i^{(n)}$ converge en $A_i$ a un elemento $a_i$.
Fijado $\varepsilon>0$, existe $N$ tal que
$\sup_i\|a_i^{(n)}-a_i^{(m)}\|<\varepsilon$ para $n,m\ge N$. Al hacer $m\to\infty$
coordenada a coordenada, obtenemos $\|a_i^{(n)}-a_i\|\le\varepsilon$ para todo $i$ y $n\ge N$.
En particular $\sup_i\|a_i\|\le\|a^{(N)}\|+\varepsilon<\infty$, y la misma estimación muestra
$\|a^{(n)}-a\|\le\varepsilon$ para $n\ge N$. Por tanto el producto acotado es completo.
Definiendo $(a_i)^*=(a_i^*)$,
\[
\|(a_i)^*(a_i)\|
=\sup_i\|a_i^*a_i\|
=\sup_i\|a_i\|^2
=\|(a_i)\|^2.
\]
Así es una $C^*$-álgebra.
\end{proof}

\proofdiagram
{$\{A_i\}_{i\in I}$ $C^*$-álgebras}
{Operaciones coordenada a coordenada; norma supremo}
{Completitud coordenada y $\|(a_i)^*(a_i)\|=\sup_i\|a_i^*a_i\|$}
{$\prod_i A_i$ acotado es una $C^*$-álgebra}

\begin{cthejemplo}{(Un producto acotado de matrices)}
En el producto de T2.3(a), cada coordenada puede ser una matriz distinta.
Tomemos $A_n=M_2(\mathbb C)$ y $a_n=\operatorname{diag}(1,1/n)$ para $n\ge1$.
La familia $a=(a_n)$ pertenece al producto acotado porque $\sup_n\|a_n\|=1$; además
\[
 a_n^*a_n=\operatorname{diag}(1,1/n^2),
 \qquad \|a^*a\|=1=\|a\|^2.
\]
La identidad $C^*$ del producto se verifica coordenada a coordenada, como en T2.3(a).
\end{cthejemplo}

\noindent\textbf{Dependencias.} Se aplican operaciones coordenada a coordenada, completitud de cada $A_i$ y la identidad $C^*$ en cada coordenada. La norma supremo restringe el producto a familias uniformemente acotadas.

\noindent\textbf{Qué prepara este ejercicio.} Construye el álgebra ambiente de T2.3(b), donde las familias que tienden a cero forman un ideal cerrado.

\subsection*{3(b) La suma $c_0$ es un ideal cerrado}

\noindent\textbf{Enunciado.} Para $I=\mathbb N$, prueba que las sucesiones $(a_n)$ del producto acotado que cumplen $\|a_n\|\to0$ forman un ideal bilateral cerrado.

\begin{proof}[Solución]
Sea
\[
\bigoplus_{n\in\mathbb N}^{c_0}A_n
=\{(a_n)\in\prod^{\ell^\infty}A_n:\|a_n\|\to0\}.
\]
Si $\|a_n\|,\|c_n\|\to0$, entonces
$\|a_n+c_n\|\le\|a_n\|+\|c_n\|\to0$ y
$\|\alpha a_n\|=|\alpha|\|a_n\|\to0$; por ello la suma $c_0$ es un subespacio lineal.
Si $b=(b_n)$ está en el producto acotado y $a=(a_n)$ en la suma $c_0$, entonces
\[
\|b_na_n\|\le\|b\|_\infty\|a_n\|\to0,
\quad
\|a_nb_n\|\le\|a_n\|\|b\|_\infty\to0,
\]
así que es un ideal bilateral. Si $a^{(k)}\to a$ en norma supremo y cada $a^{(k)}$ tiende a cero en
coordenadas, fijado $\varepsilon$ elegimos $k$ con
$\|a-a^{(k)}\|_\infty<\varepsilon/2$ y después $N$ con
$\|a_n^{(k)}\|<\varepsilon/2$ para $n\ge N$. Entonces $\|a_n\|<\varepsilon$ para $n\ge N$.
Por tanto el ideal es cerrado.
Además $(a_n^*)$ pertenece al mismo ideal porque $\|a_n^*\|=\|a_n\|\to0$.
\end{proof}

\proofdiagram
{$a=(a_n)$ con $\|a_n\|\to0$}
{Producto coordenada y norma supremo}
{$\|b_na_n\|\le\|b\|_\infty\|a_n\|\to0$; límite uniforme preserva la condición}
{$\bigoplus_n A_n$ es ideal bilateral cerrado}

\begin{cthejemplo}{(El ideal clásico de sucesiones nulas)}
Para reconocer la construcción de T2.3(b), pongamos $A_n=\mathbb C$ para cada $n$.
La sucesión $a_n=1/n$ pertenece a $c_0$, mientras que $b_n=(-1)^n$ está en $\ell^\infty$.
Su producto satisface $|b_na_n|=1/n\to0$, y lo mismo ocurre con $a_nb_n$.
Así el ideal $c_0\subset\ell^\infty$ es el caso escalar de la propiedad bilateral probada en T2.3(b).
\end{cthejemplo}

\noindent\textbf{Dependencias.} El producto acotado de T2.3(a) aporta la norma supremo. Las estimaciones coordenada a coordenada prueban la propiedad de ideal y el criterio uniforme prueba el cierre.

\noindent\textbf{Qué prepara este ejercicio.} Permite formar el cociente del producto por el ideal $c_0$ mediante T1.4; interpretar propiedades de ese cociente requiere resultados adicionales sobre cocientes de $C^*$-álgebras.

\section{Problema 4: teorema del mapeo espectral continuo}

Sea $\A$ una $C^*$-álgebra unital, $x\in\A$ normal y $f:\spec(x)\to\mathbb C$ continua.

\noindent\textbf{Enunciado.} Prueba el teorema del mapeo espectral continuo: $\spec_\A(f(x))=f(\spec_\A(x))$.

\begin{proof}[Solución]
Trabajemos en $B=C^*(x,1)$, que es conmutativa. Por permanencia espectral, calcular el espectro de
$f(x)$ en $B$ o en $\A$ da lo mismo. El cálculo funcional es un $*$-isomorfismo
\[
 C(\spec(x))\longrightarrow B.
\]
Bajo este isomorfismo, $f(x)$ corresponde a la función $f$. Por el Problema 10 de la Tarea 1, el
espectro de $f$ en $C(\spec(x))$ es su rango. Por tanto
\[
 \spec_\A(f(x))=\spec_B(f(x))=f(\spec(x)).
\]
\end{proof}

\proofdiagrambranch
{$x$ normal}
{$f\in C(\spec(x))$}
{Cálculo funcional + permanencia espectral en $C^*(x,1)$}
{Identificar $C^*(x,1)\cong C(\spec(x))$}
{$\spec(f(x))=f(\spec(x))$}

\begin{cthejemplo}{(Cuadrar un autoadjunto mediante el espectro)}
En $M_2(\mathbb C)$, tomemos $x=\operatorname{diag}(-2,3)$ para comprobar T2.4.
El elemento es autoadjunto y $\spec(x)=\{-2,3\}$.
Como $x^2=\operatorname{diag}(4,9)$,
\[
 \spec(x^2)=\{4,9\}=\{\lambda^2:\lambda\in\spec(x)\}.
\]
El cuadrado transforma cada valor espectral mediante $t\mapsto t^2$ y produce espectro no negativo.
\end{cthejemplo}

\noindent\textbf{Dependencias.} El cálculo funcional para $x$ normal identifica $C^*(x,1)$ con $C(\spec(x))$; la permanencia espectral de la Sección 2 y T1.10 permiten calcular el espectro de $f(x)$ en la álgebra ambiente.

\noindent\textbf{Qué prepara este ejercicio.} Transporta igualdades y restricciones espectrales de funciones a operadores; T2.6 lo aplica para reconocer autoadjuntos, unitarios y proyecciones normales.

\section{Problema 5: $X\cong\Omega(C(X))$}

Sea $X$ un espacio compacto Hausdorff y, para $x\in X$, sea
$\operatorname{ev}_x(f)=f(x)$ la aplicación de evaluación sobre $C(X)$.

\noindent\textbf{Enunciado.} Prueba que $x\mapsto\operatorname{ev}_x$ define un homeomorfismo $X\to\Omega(C(X))$.

\begin{proof}[Solución]
Para $f,g\in C(X)$, se tiene
$\operatorname{ev}_x(fg)=f(x)g(x)$ y $\operatorname{ev}_x(1)=1$;
además $|\operatorname{ev}_x(f)|\le\|f\|_\infty$. Así cada evaluación es un carácter continuo
y el mapa está bien definido. Es inyectivo porque
si $x\neq y$, existe $f\in C(X)$ con $f(x)\neq f(y)$.

Sea ahora $\varphi\in\Omega(C(X))$ y $I=\ker\varphi$, que es maximal. Afirmamos que existe
$x_0\in X$ tal que $I\subset I_{x_0}:=\{f:f(x_0)=0\}$. Si no, para cada $x\in X$ existiría
$f_x\in I$ con $f_x(x)\neq0$. Los abiertos
$U_x=\{y:|f_x(y)|>0\}$ cubrirían $X$; por compacidad escogeríamos
$f_1,\dots,f_n\in I$ con $\sum_j|f_j|^2>0$ en todo $X$. Entonces
\[
 g:=\sum_{j=1}^n f_j^*f_j\in I
\]
sería una función estrictamente positiva y por tanto invertible en $C(X)$, imposible para un ideal
propio. Así existe $x_0$ con $I\subset I_{x_0}$. Como $I_{x_0}$ es propio, la maximalidad da
$I=I_{x_0}$. Para cada $f\in C(X)$, tenemos $f-f(x_0)1\in I=\ker\varphi$; como
$\varphi(1)=1$, resulta $\varphi(f)=f(x_0)$. Por tanto
$\varphi=\operatorname{ev}_{x_0}$.

El mapa $x\mapsto\operatorname{ev}_x$ es continuo para la topología débil-$*$ porque, para cada
$f\in C(X)$, la coordenada $x\mapsto\operatorname{ev}_x(f)=f(x)$ es continua. Como $X$ es compacto
y $\Omega(C(X))$ es Hausdorff, una biyección continua es un homeomorfismo.
\end{proof}

\proofdiagrambranch
{$x\mapsto\operatorname{ev}_x$}
{Funciones continuas separan puntos en compactos Hausdorff}
{Para sobreyectividad, $\ker\varphi$ es ideal maximal y debe ser $I_{x_0}$ para algún $x_0$}
{Biyectividad continua entre compactos Hausdorff}
{$X\cong\Omega(C(X))$}

\begin{cthejemplo}{(Un carácter dado por un punto del círculo)}
En $C(\mathbb T)$ queremos hacer visible el homeomorfismo de T2.5.
Fijemos $w=e^{i\theta}\in\mathbb T$ y consideremos la evaluación
$\operatorname{ev}_w(f)=f(w)$. Para las funciones $h(z)=z$ y $g(z)=1+z$,
\[
 \operatorname{ev}_w(h)=w,\qquad \operatorname{ev}_w(g)=1+w.
\]
El valor del carácter sobre $h$ identifica el punto $w$; T2.5 asegura que todo carácter de
$C(\mathbb T)$ surge de una evaluación de este tipo.
\end{cthejemplo}

\noindent\textbf{Dependencias.} Se usan caracteres e ideales maximales de la Parte 1, separación de puntos de $C(X)$ y compacidad de $X$. La demostración del conjunto común de ceros usa $\sum_j|f_j|^2$.

\noindent\textbf{Qué prepara este ejercicio.} Completa la unicidad del espacio compacto en Gelfand--Naimark conmutativo: un $*$-isomorfismo $C(X)\cong C(Y)$ induce $X\cong Y$.

\section{Problema 6: caracterizaciones espectrales y necesidad de normalidad}

Sea $\A$ una $C^*$-álgebra unital y $x\in\A$ normal. El inciso (d) pide comprobar que la
normalidad no puede omitirse en las tres recíprocas.

\subsection*{6(a) Autoadjunto $\Longleftrightarrow$ espectro real para normales}

\noindent\textbf{Enunciado.} Prueba que $x=x^*$ si y sólo si $\spec(x)\subset\mathbb R$.

\begin{proof}[Solución]
La implicación directa ya se demostró en la Sección 2. Para la recíproca, por cálculo funcional,
$x$ corresponde a la función $z$ y $x^*$ a $\overline z$ sobre $\spec(x)$. Si el espectro es real,
ambas funciones coinciden, luego $x=x^*$.
\end{proof}

\proofdiagram
{$x$ normal}
{Espectro de autoadjunto real; cálculo funcional}
{Si $\spec(x)\subset\mathbb R$, entonces $z=\overline z$ sobre el espectro}
{$x=x^*$}

\begin{cthejemplo}{(Espectro real sin autoadjunción)}
La recíproca de T2.6(a) necesita normalidad; una matriz triangular permite verlo.
Para $N=\begin{pmatrix}0&1\\0&0\end{pmatrix}$,
$\det(\lambda I-N)=\lambda^2$, luego $\spec(N)=\{0\}\subset\mathbb R$.
Pero $N^*=\begin{pmatrix}0&0\\1&0\end{pmatrix}\ne N$.
Además $N^*N=\operatorname{diag}(0,1)\ne\operatorname{diag}(1,0)=NN^*$.
El espectro real no obliga a ser autoadjunto sin normalidad, como se precisa en T2.6(d).
\end{cthejemplo}

\noindent\textbf{Dependencias.} La realidad del espectro de un autoadjunto se probó en la Sección 2; la recíproca usa cálculo funcional en $C^*(x,1)$, disponible porque $x$ es normal.

\noindent\textbf{Qué prepara este ejercicio.} Identifica exactamente qué información espectral reconoce la autoadjunción. T2.6(d) muestra con una matriz nilpotente por qué la hipótesis de normalidad es necesaria.

\subsection*{6(b) Unitario $\Longleftrightarrow$ espectro en $\mathbb T$ para normales}

\noindent\textbf{Enunciado.} Prueba que $x$ es unitario si y sólo si $\spec(x)\subset\mathbb T$.

\begin{proof}[Solución]
Si $x$ es unitario, el Problema 2(c) da $\spec(x)\subset\mathbb T$. Recíprocamente, si $x$ es
normal y el espectro está en $\mathbb T$, entonces
\[
 x^*x=(\overline z z)(x)=(|z|^2)(x)=1.
\]
Como $x$ es normal, $xx^*=x^*x=1$, así que es unitario.
\end{proof}

\proofdiagram
{$x$ normal}
{Problema 2(c); cálculo funcional}
{$\spec(x)\subset\mathbb T\Rightarrow |z|^2=1$ en el espectro}
{$x^*x=(\overline z z)(x)=1$ y, por normalidad, $xx^*=1$}

\begin{cthejemplo}{(Espectro unitario sin unitariedad)}
Para probar el límite de la recíproca en T2.6(b), tomemos
$J=\begin{pmatrix}1&1\\0&1\end{pmatrix}$.
Su polinomio característico es $(\lambda-1)^2$, de modo que
$\spec(J)=\{1\}\subset\mathbb T$, pero
$J^*J=\begin{pmatrix}1&1\\1&2\end{pmatrix}\ne I$.
El espectro está en el círculo y la matriz no es unitaria: T2.6(d) explica el papel de la normalidad.
\end{cthejemplo}

\noindent\textbf{Dependencias.} T2.2(c) da la dirección directa. En la recíproca, la normalidad permite el cálculo funcional y la identidad $x^*x=(|z|^2)(x)$.

\noindent\textbf{Qué prepara este ejercicio.} Permite reconocer unitarios a partir de su espectro, por ejemplo en álgebras de operadores con shifts unitarios cuando se han verificado sus hipótesis.

\subsection*{6(c) Proyección $\Longleftrightarrow$ espectro en $\{0,1\}$ para normales}

\noindent\textbf{Enunciado.} Prueba que $x$ es una proyección si y sólo si $\spec(x)\subset\{0,1\}$.

\begin{proof}[Solución]
Si $x$ es una proyección, es idempotente y el Problema 2(a) da
$\spec(x)\subset\{0,1\}$. Para la recíproca, si $x$ es normal y su espectro está contenido en
$\{0,1\}$, el cálculo funcional da
\[
 x^2-x=(z^2-z)(x)=0,
 \qquad
 x^*-x=(\overline z-z)(x)=0.
\]
Así $x$ es una proyección.
\end{proof}

\proofdiagram
{$x$ normal}
{Espectro de idempotente; cálculo funcional}
{$z^2-z=0$ y $\overline z-z=0$ sobre $\{0,1\}$}
{$x^2=x$ y $x^*=x$}

\begin{cthejemplo}{(Idempotente con espectro correcto que no es proyección)}
En T2.6(c) la autoadjunción forma parte de ser proyección. Consideremos
$E=\begin{pmatrix}1&1\\0&0\end{pmatrix}$.
Se calcula $E^2=E$ y, por triangularidad, $\spec(E)=\{0,1\}$; no obstante,
$E^*=\begin{pmatrix}1&0\\1&0\end{pmatrix}\ne E$.
Así el espectro permitido y la idempotencia no bastan sin la normalidad que usa la
recíproca de T2.6(c).
\end{cthejemplo}

\noindent\textbf{Dependencias.} T2.2(a) da la dirección directa. Para la recíproca, la normalidad permite transportar $z^2=z$ y $\overline z=z$ desde el espectro a identidades de $x$.

\noindent\textbf{Qué prepara este ejercicio.} Conecta la clasificación espectral de la Sección 2 con proyecciones de Cauchy sólo cuando, además de idempotencia, se verifica la autoadjunción de esas proyecciones.

\subsection*{6(d) ¿Dónde se usa la normalidad?}

\noindent\textbf{Enunciado.} Exhibe en $M_2(\mathbb C)$ tres elementos no normales que invaliden las recíprocas de los incisos (a)--(c) sin normalidad, y localiza el paso de cada prueba que usa esa hipótesis.

\begin{proof}[Solución]
Tomemos
\[
 N=\begin{pmatrix}0&1\\0&0\end{pmatrix},
 \qquad J=I+N=\begin{pmatrix}1&1\\0&1\end{pmatrix},
 \qquad E=\begin{pmatrix}1&1\\0&0\end{pmatrix}.
\]
Como son triangulares, $\spec(N)=\{0\}\subset\mathbb R$,
$\spec(J)=\{1\}\subset\mathbb T$ y $\spec(E)=\{0,1\}$.
Sin embargo, $N\ne N^*$,
\[
 J^*J=\begin{pmatrix}1&1\\1&2\end{pmatrix}\ne I.
\]
$E$ no es una proyección porque, aunque $E^2=E$, se tiene $E^*\ne E$.
Las tres son no normales: en efecto,
$N^*N=\operatorname{diag}(0,1)\ne\operatorname{diag}(1,0)=NN^*$ y
$J^*J-JJ^*=N^*N-NN^*\ne0$; además,
$E^*E=\begin{pmatrix}1&1\\1&1\end{pmatrix}$ y
$EE^*=\operatorname{diag}(2,0)$ son distintos.

La normalidad se usó en los incisos (a)--(c) para que $x$ y $x^*$ generen la
$C^*$-álgebra conmutativa $C^*(x,1)$ y se pueda aplicar el cálculo funcional continuo.
En (b) se usó además la igualdad $xx^*=x^*x$ tras obtener $x^*x=1$.
\end{proof}

\proofdiagram
{$N$, $J=I+N$ y $E$ son no normales}
{El cálculo funcional continuo sólo se construyó para normales}
{La igualdad de funciones sobre el espectro se transporta a igualdad de elementos sólo vía ese cálculo}
{Sin normalidad, el espectro no determina esas relaciones algebraicas}

\begin{cthejemplo}{(Dos operadores con el mismo espectro)}
Para ver qué pierde el espectro sin normalidad, comparemos $I$ con
$J=\begin{pmatrix}1&1\\0&1\end{pmatrix}$ en $M_2(\mathbb C)$.
Ambos tienen espectro $\{1\}$, pero $I^*I=I$, mientras que
$J^*J=\begin{pmatrix}1&1\\1&2\end{pmatrix}\ne I$.
Sólo $I$ es unitario. La comparación concreta explica por qué T2.6(d) exige normalidad
antes de recuperar propiedades algebraicas desde el espectro.
\end{cthejemplo}

\noindent\textbf{Dependencias.} Los incisos T2.6(a)--(c) usan cálculo funcional sobre $C^*(x,1)$; esa subálgebra es conmutativa porque $x$ es normal. Los bloques de Jordan proporcionan contraejemplos.

\noindent\textbf{Qué prepara este ejercicio.} Delimita el alcance de la Proposición 2.28 y evita inferir autoadjunción, unitariedad o proyección a partir del espectro de un operador no normal.

\pagebreak
\section{Problema 7: parte unitaria de la descomposición polar}

Sea $\A$ una $C^*$-álgebra unital. Se define $|a|$ para todo $a\in\A$ y se supone que $a$ es
invertible al construir el factor unitario.


Definimos $|a|=(a^*a)^{1/2}$ por cálculo funcional positivo.

\noindent\textbf{Enunciado.} Para $a\in\GL(\A)$, prueba que $u=a|a|^{-1}$ es unitario, donde $|a|=(a^*a)^{1/2}$.

\begin{proof}[Solución]
Si $a$ es invertible, entonces $a^*a$ tiene inverso $a^{-1}(a^*)^{-1}$.
Como $|a|^2=a^*a$ es invertible, también lo es $|a|$: los elementos $|a|$ y
$(|a|^2)^{-1}$ conmutan, y $|a|(|a|^2)^{-1}$ es un inverso bilateral de $|a|$.
Podemos, por tanto, definir $u=a|a|^{-1}$.
Como $|a|$ es positivo, $|a|^*=|a|$ y
\[
 u^*=|a|^{-1}a^*.
\]
Entonces
\[
 u^*u=|a|^{-1}a^*a|a|^{-1}
 =|a|^{-1}|a|^2|a|^{-1}=1.
\]
Además
\[
 uu^*=a|a|^{-2}a^*
 =a(a^*a)^{-1}a^*
 =a\,a^{-1}(a^*)^{-1}a^*=1.
\]
Así $u$ es unitario y $a=u|a|$.
\end{proof}

\proofdiagram
{$a\in\GL(\A)$}
{Raíz positiva $|a|=(a^*a)^{1/2}$}
{Definir $u=a|a|^{-1}$ y calcular $u^*u$, $uu^*$}
{$u^*u=uu^*=1$; $u$ es unitario}

\begin{cthejemplo}{(Descomposición polar de una matriz diagonal)}
Para aplicar la fórmula de T2.7, tomemos la matriz invertible
$A=\operatorname{diag}(2i,3)$.
Entonces $A^*A=\operatorname{diag}(4,9)$,
$|A|=\operatorname{diag}(2,3)$ y
\[
 U=A|A|^{-1}=\operatorname{diag}(i,1),\qquad U^*U=UU^*=I.
\]
La igualdad $A=U|A|$ separa una parte unitaria de una parte positiva, como afirma T2.7.
\end{cthejemplo}

\noindent\textbf{Dependencias.} La raíz positiva $|a|=(a^*a)^{1/2}$ procede del cálculo funcional de la Sección 2. La invertibilidad de $a$ hace invertibles a $a^*a$ y $|a|$.

\noindent\textbf{Qué prepara este ejercicio.} Construye la descomposición polar $a=u|a|$ de un invertible y aclara dónde se usa la invertibilidad para obtener un factor unitario en la misma álgebra.

\Needspace{15\baselineskip}
\section{Problema 8: continuidad del cálculo funcional bajo convergencia de normales}

Sea $\A$ una $C^*$-álgebra unital y sea $(x_n)$ una sucesión de elementos
normales que converge en norma a $x$. Fijamos un compacto $K\subset\mathbb C$
que contiene $\spec(x_n)$ para todo $n$ y $\spec(x)$, y tomamos $f\in C(K)$.

\noindent\textbf{Enunciado.} Prueba que $x$ es normal y que $f(x_n)\to f(x)$ en norma para toda $f\in C(K)$.

\begin{proof}[Solución]
Por la Tarea 1.7, $x$ es normal. Sea $f\in C(K)$ y $\varepsilon>0$. Por Stone--Weierstrass existe
un polinomio $p(z,\overline z)$ tal que
\[
 \sup_{z\in K}|f(z)-p(z,\overline z)|<\varepsilon.
\]
El cálculo funcional es contractivo, de modo que para todo $n$,
\[
\|f(x_n)-p(x_n,x_n^*)\|<\varepsilon,
\qquad
\|f(x)-p(x,x^*)\|<\varepsilon.
\]
Como $x_n\to x$, también $x_n^*\to x^*$, y por continuidad de suma y producto
$p(x_n,x_n^*)\to p(x,x^*)$. Por tanto, para $n$ grande,
\[
\|f(x_n)-f(x)\|
\le 2\varepsilon+\|p(x_n,x_n^*)-p(x,x^*)\|<3\varepsilon.
\]
Como $\varepsilon$ es arbitrario, $f(x_n)\to f(x)$.
\end{proof}

\proofdiagrambranch
{$x_n\to x$ y cada $x_n$ normal}
{$\spec(x_n),\spec(x)\subset K$ compacto}
{Normalidad cerrada + Stone--Weierstrass aproxima $f$ por polinomios en $z,\bar z$}
{Polinomios son continuos en $x_n,x_n^*$; controlar el error uniformemente}
{$f(x_n)\to f(x)$}

\begin{cthejemplo}{(Continuidad de la raíz positiva)}
Para ver T2.8 en matrices positivas, fijemos
$x_n=\operatorname{diag}(1+1/n,4)$ y $x=\operatorname{diag}(1,4)$.
Todos sus espectros están en $K=[1,5]$ y $x_n\to x$ en norma.
Con $f(t)=\sqrt t$,
\[
 f(x_n)=\operatorname{diag}(\sqrt{1+1/n},2)
 \longrightarrow\operatorname{diag}(1,2)=f(x).
\]
El cálculo muestra que la raíz positiva varía continuamente bajo las hipótesis de T2.8.
\end{cthejemplo}

\noindent\textbf{Dependencias.} T1.7 garantiza que $x$ es normal; Stone--Weierstrass aproxima $f$ uniformemente por polinomios en $z,\overline z$, y el cálculo funcional controla el error.

\noindent\textbf{Qué prepara este ejercicio.} Muestra la estabilidad en norma de $f(x_n)$, en particular de raíces positivas cuando los espectros están en un mismo compacto de $[0,\infty)$.

\Needspace{17\baselineskip}
\section{Problema 9: convergencia monótona de operadores autoadjuntos}

Sea $H$ un espacio de Hilbert. Una red está indexada por un conjunto dirigido $\Lambda$.
Para autoadjuntos $R,U\in\mathcal B(H)$, escribimos $R\le U$ cuando $U-R$ es positivo,
es decir, cuando $\langle(U-R)\xi,\xi\rangle\ge0$ para todo $\xi\in H$.


Sea $(T_\lambda)_{\lambda\in\Lambda}$ una red acotada de autoadjuntos con
$T_{\lambda_1}\le T_{\lambda_2}$ si $\lambda_1\le\lambda_2$.
Aquí ``convergencia puntual'' significa que, para cada $\xi\in H$,
$\|T_\lambda\xi-T\xi\|\to0$; es la convergencia fuerte de operadores.

\noindent\textbf{Enunciado.} Prueba que existe un autoadjunto $T\in\mathcal B(H)$ tal que $\|T_\lambda\xi-T\xi\|\to0$ para cada $\xi\in H$.

\begin{proof}[Solución]
Sea $M=\sup_\lambda\|T_\lambda\|$. Para cada $\xi\in H$, los números
$q_\lambda(\xi)=\langle T_\lambda\xi,\xi\rangle$ forman una red creciente y acotada por
$M\|\xi\|^2$, por lo que tienen límite $q(\xi)$. Mediante la identidad de polarización obtenemos
una forma sesquilineal acotada
\[
 B(\xi,\eta)=\lim_\lambda\langle T_\lambda\xi,\eta\rangle.
\]
En efecto, $|\langle T_\lambda\xi,\eta\rangle|\le M\|\xi\|\|\eta\|$, y la misma cota pasa al
límite.
Con la convención de que el producto interior es lineal en la primera variable, la identidad de
polarización expresa este límite mediante cuatro valores de $q$:
\[
 B(\xi,\eta)=\frac14\bigl(q(\xi+\eta)-q(\xi-\eta)
   +i q(\xi+i\eta)-i q(\xi-i\eta)\bigr).
\]
Por Riesz existe un único $T\in\mathcal B(H)$ con
$B(\xi,\eta)=\langle T\xi,\eta\rangle$, y la cota de $B$ implica $\|T\|\le M$.
Como cada $T_\lambda$ es autoadjunto, $T$ también lo es, y
$T_\lambda\le T$.

Pongamos $A_\lambda=T-T_\lambda\ge0$. Entonces
$\|A_\lambda\|\le 2M$. Como
$\spec(A_\lambda)\subset[0,\|A_\lambda\|]$, la función
$t\mapsto t(\|A_\lambda\|-t)$ es no negativa en ese espectro. El cálculo funcional positivo
da $0\le A_\lambda^2\le\|A_\lambda\|A_\lambda\le2M A_\lambda$. Por tanto
\[
\|(T-T_\lambda)\xi\|^2
=\langle A_\lambda^2\xi,\xi\rangle
\le2M\langle A_\lambda\xi,\xi\rangle.
\]
El último término tiende a cero por definición de $T$. Así
$T_\lambda\xi\to T\xi$ en norma para cada $\xi$.
\end{proof}

\proofdiagrambranch
{Monotonía: $\langle T_\lambda\xi,\xi\rangle$ creciente}
{Acotación uniforme $\sup\|T_\lambda\|<\infty$}
{Definir una forma límite y usar Riesz para obtener $T$}
{Usar $0\le(T-T_\lambda)^2\le\|T-T_\lambda\|(T-T_\lambda)$}
{$T_\lambda\xi\to T\xi$ para cada $\xi$}

\begin{cthejemplo}{(Convergencia fuerte de matrices crecientes)}
Para observar el límite de T2.9, tomemos en $\mathbb C^2$
$T_n=\operatorname{diag}(1-1/n,2-1/n)$, $n\ge1$.
Se tiene $T_{n+1}-T_n=(1/n-1/(n+1))I\ge0$ y $\|T_n\|\le2$.
Con $T=\operatorname{diag}(1,2)$,
\[
 \|(T-T_n)\xi\|=\frac1n\|\xi\|\longrightarrow0
 \qquad(\xi\in\mathbb C^2).
\]
La convergencia aquí incluso es en norma operatorial; T2.9 garantiza la convergencia fuerte
para redes acotadas en cualquier espacio de Hilbert.
\end{cthejemplo}

\noindent\textbf{Dependencias.} Se usan la convergencia de redes monótonas acotadas de números reales, la identidad de polarización, la representación de Riesz y el cálculo funcional positivo para justificar $0\le A^2\le\|A\|A$ cuando $A\ge0$.

\noindent\textbf{Qué prepara este ejercicio.} Obtiene convergencia fuerte: $\|T_\lambda\xi-T\xi\|\to0$ para cada $\xi\in H$. Este resultado prepara técnicas de topologías de operadores, pero no es una consecuencia de la convergencia en norma de T2.8.

\section{Problema 10: espectro del shift unilateral}


Sea $S:\ell^2(\mathbb N)\to\ell^2(\mathbb N)$,
\[
S(x_1,x_2,\dots)=(0,x_1,x_2,\dots).
\]

\noindent\textbf{Enunciado.} Prueba que el espectro del shift unilateral es $\spec(S)=\overline{\mathbb D}$.

\begin{proof}[Solución]
Como $S$ es una isometría, $\|S\|=1$ y por la cota espectral
$\spec(S)\subset\overline{\mathbb D}$. Ahora sea $|\lambda|<1$. El vector
\[
 v_\lambda=(1,\overline\lambda,\overline\lambda^2,\ldots)
\]
pertenece a $\ell^2$ y satisface
\[
 S^*v_\lambda=\overline\lambda\,v_\lambda.
\]
Así $\overline\lambda I-S^*$ no es invertible. Si $\lambda I-S$ fuese invertible, entonces su
adjunto $\overline\lambda I-S^*$ también sería invertible, contradicción. Por tanto
$\lambda\in\spec(S)$ para todo $|\lambda|<1$. Como el espectro es cerrado,
$\overline{\mathbb D}\subset\spec(S)$. Junto con la inclusión opuesta,
\[
 \spec(S)=\overline{\mathbb D}.
\]
\end{proof}

\proofdiagrambranch
{$\|S\|=1\Rightarrow\spec(S)\subset\overline{\mathbb D}$}
{Para $|\lambda|<1$, $S^*$ tiene vector propio de valor $\overline\lambda$}
{Si $\lambda I-S$ fuera invertible, su adjunto también lo sería}
{$\mathbb D\subset\spec(S)$; el espectro es cerrado}
{$\spec(S)=\overline{\mathbb D}$}

\begin{cthejemplo}{(Un disco espectral sin valores propios)}
El shift unilateral de T2.10 ilustra la diferencia entre espectro y valores propios.
Si $S\xi=\lambda\xi$, la primera coordenada da $0=\lambda\xi_1$; para
$\lambda\ne0$, las coordenadas siguientes obligan sucesivamente a $\xi_n=0$.
Para $\lambda=0$, la inyectividad de $S$ da igualmente $\xi=0$.
Así $S$ no tiene valores propios, aunque T2.10 demuestra
$\spec(S)=\overline{\mathbb D}$.
\end{cthejemplo}

\noindent\textbf{Dependencias.} La cota espectral de la Parte 1 da $\spec(S)\subset\overline{\mathbb D}$. T1.11(d) ya muestra $0\in\spec(S)$; los vectores propios de $S^*$ dan todos los puntos interiores y el cierre del espectro añade la frontera.

\noindent\textbf{Qué prepara este ejercicio.} Contrasta el espectro de un operador no normal con la caracterización de unitarios de T2.6(b). No permite atribuir el mismo disco espectral a los shifts de composición del problema de Haseman.

\chapter{Ejemplos de investigación: cómo reaparece la teoría en Haseman y operadores integrales singulares}

\section{De la condición de frontera al operador}

En Karlovich \cite[pp.~463--464, ecs.~(1.1)--(1.2)]{Karlovich2017}, $\Gamma$ es una
estrella abierta de espirales logarítmicas y $1<p<\infty$.
Se toman $G\in\SO(\Gamma)$, $g\in L^p(\Gamma)$ y un difeomorfismo
$\alpha:\Gamma\to\Gamma$ que preserva la orientación de cada rama y
satisface las hipótesis de la fuente. La condición de Haseman es
\[
 \Phi^+(\alpha(t))=G(t)\Phi^-(t)+g(t),\qquad t\in\Gamma.
\]
Aquí $S_\Gamma\in\mathcal B(L^p(\Gamma))$ es el operador integral singular de Cauchy,
$V_\alpha f=f\circ\alpha$ y $M_Gf=Gf$. Si $\varphi$ es la densidad de Cauchy, las
fórmulas de Sokhotski--Plemelj dan
\[
 \Phi^+=P_+\varphi,\qquad \Phi^-=-P_-\varphi,\qquad
 P_\pm=\tfrac12(I\pm S_\Gamma).
\]
El signo de la traza negativa explica el signo más en la ecuación operatorial:
\[
 \bigl(V_\alpha P_++M_GP_-\bigr)\varphi=g,\qquad
 T:=V_\alpha P_++M_GP_-\in\mathcal B(L^p(\Gamma)).
\]
El curso proporciona el lenguaje para preguntar si $T$ es invertible y estudiar su
espectro. La reducción de frontera y el criterio de Fredholm son resultados posteriores
que dependen de las hipótesis geométricas y de símbolos del artículo.

\begin{cthejemplo}{(La identidad como caso de la ecuación de Haseman)}
En $X=L^p(\Gamma)$ consideremos la fórmula $T=V_\alpha P_++M_GP_-$
obtenida arriba. Queremos comprobar qué operador resulta cuando
$\alpha=\mathrm{id}_\Gamma$ y $G=1$. En ese caso $V_\alpha=I_X$ y
$M_G=I_X$, mientras que la definición de $P_\pm$ da
\[
 P_++P_- =\tfrac12(I_X+S_\Gamma)+\tfrac12(I_X-S_\Gamma)=I_X.
\]
Por tanto $T=P_++P_-=I_X$. Este cálculo docente muestra que la
reducción por Plemelj se puede leer sin suponer $S_\Gamma^2=I_X$.
\end{cthejemplo}

\begin{center}
\begin{tikzpicture}[node distance=5mm and 8mm]
\node[theorybox,text width=3.35cm] (b) {Frontera de Haseman\\$\Phi^+\circ\alpha=G\Phi^-+g$};
\node[theorybox,text width=3.35cm,right=of b] (p) {Sokhotski--Plemelj\\$\Phi^+=P_+\varphi$, $\Phi^-=-P_-\varphi$};
\node[researchbox,text width=3.35cm,right=of p] (t) {$T=V_\alpha P_++M_GP_-$\\$T\varphi=g$};
\draw[arr] (b)--(p);
\draw[arr] (p)--(t);
\end{tikzpicture}
\end{center}

\section{Del ideal compacto a la propiedad de Fredholm}

Sea $X$ un espacio de Banach complejo infinito-dimensional. El ideal bilateral cerrado
$\mathcal K(X)\triangleleft\mathcal B(X)$ produce la álgebra de Calkin
$\mathcal Q(X)=\mathcal B(X)/\mathcal K(X)$. Para $A\in\mathcal B(X)$,
\[
 A\text{ Fredholm}
 \quad\Longleftrightarrow\quad
 A+\mathcal K(X)\in\GL(\mathcal Q(X)).
\]
Karlovich enuncia esta equivalencia en la introducción de su artículo de 2025
\cite[p.~2]{Karlovich2025}. La Tarea 1.4 explica por qué el cociente por un ideal
cerrado sigue siendo una álgebra de Banach; la equivalencia de Fredholm añade un
teorema propio de operadores.

\begin{cthejemplo}{(Invertible módulo compactos sin inversa exacta)}
El criterio de Atkinson anterior distingue la inversa exacta de una
inversa en Calkin. Para verlo en $X=\ell^2$, tome $P$ la proyección ortogonal de rango uno sobre
$\mathbb Ce_1$ y $A=I_X-P$. Se tiene
$\ker A=\operatorname{Ran}P\ne\{0\}$, así que $A$ no es invertible. Sin embargo,
$P\in\mathcal K(X)$ y $A+\mathcal K(X)=I_X+\mathcal K(X)$: su clase es la
unidad de $\mathcal Q(X)$. Además,
$\operatorname{Ran}(A)=(\mathbb Ce_1)^\perp$ es cerrado, mientras que el núcleo
y el conúcleo tienen dimensión uno. Así $A$ es Fredholm. El cálculo es una
elaboración docente del criterio anterior.
\end{cthejemplo}

Si se define el \emph{espectro esencial de Fredholm} por
\[
 \sigma_{\mathrm{ess}}(A)
 :=\sigma_{\mathcal Q(X)}\bigl(A+\mathcal K(X)\bigr),
\]
entonces $0\notin\sigma_{\mathrm{ess}}(A)$ equivale a que $A$ sea Fredholm.
Esto no dice que $0\notin\sigma_{\mathcal B(X)}(A)$; el ejemplo $I_X-P$
muestra la diferencia.

\begin{center}
\begin{tikzpicture}[node distance=5mm and 7mm]
\node[theorybox,text width=2.65cm] (ideal) {$\mathcal K(X)\triangleleft\mathcal B(X)$\\ideal cerrado};
\node[theorybox,text width=2.65cm,right=of ideal] (quot) {$\mathcal Q(X)$\\cociente};
\node[theorybox,text width=2.65cm,right=of quot] (class) {$A+\mathcal K(X)$\\clase invertible};
\node[researchbox,text width=2.65cm,right=of class] (fred) {$A$ Fredholm\\criterio de Calkin};
\draw[arr] (ideal)--(quot);
\draw[arr] (quot)--(class);
\draw[arr] (class)--(fred);
\end{tikzpicture}
\end{center}

En un espacio de Hilbert $H$, $\mathcal B(H)$ y su cociente por
$\mathcal K(H)$ son $C^*$-álgebras. Para un Banach general, el argumento
anterior sólo necesita la estructura de álgebra de Banach. Los cálculos de
símbolos y la localización intentan convertir la invertibilidad de la clase
en condiciones verificables para una \emph{álgebra de operadores específica};
no basta con escribir el cociente para obtener un criterio.

\Needspace{17\baselineskip}
\section{Proyecciones de Cauchy: el caso modelo y el caso Haseman}

La comparación precisa qué hipótesis convierte la fórmula
$(I\pm S_\Gamma)/2$ en una proyección y qué parte de ella se utiliza en
la ecuación de Haseman.

En el modelo $\Gamma=\mathbb R$ o $\mathbb T$ sobre $L^2(\Gamma)$,
sea $S_\Gamma$ el operador integral singular de Cauchy.
Bastos--Fernandes--Karlovich recuerdan $S_\Gamma^2=I$ y que
\[
 P_\Gamma^\pm=\tfrac12(I\pm S_\Gamma)
\]
son proyecciones ortogonales \cite[§2, p.~22]{BFK2006}. La identidad
$S_\Gamma^2=I$ da $P_\Gamma^\pm{}^2=P_\Gamma^\pm$ y
$P_\Gamma^+P_\Gamma^-=0$; la autoadjunción de $S_\Gamma$ da además
$(P_\Gamma^\pm)^*=P_\Gamma^\pm$. Es ésta la conexión literal con las
proyecciones $p=p^*=p^2$ de una $C^*$-álgebra.

\begin{center}
\begin{tikzpicture}[node distance=5mm and 7mm]
\node[theorybox,text width=2.6cm] (s) {$S_\Gamma^2=I$\\$S_\Gamma^*=S_\Gamma$};
\node[theorybox,text width=2.6cm,right=of s] (def) {$P_\Gamma^\pm=(I\pm S_\Gamma)/2$};
\node[theorybox,text width=2.6cm,right=of def] (proj) {$P_\Gamma^\pm{}^2=P_\Gamma^\pm$\\$(P_\Gamma^\pm)^*=P_\Gamma^\pm$};
\node[researchbox,text width=2.6cm,right=of proj] (split) {$L^2(\Gamma)=\operatorname{Ran}P_\Gamma^+\oplus\operatorname{Ran}P_\Gamma^-$};
\draw[arr] (s)--(def);
\draw[arr] (def)--(proj);
\draw[arr] (proj)--(split);
\end{tikzpicture}
\end{center}

En la estrella abierta de Karlovich 2017, la misma \emph{fórmula}
$P_\pm=(I\pm S_\Gamma)/2$ aparece en Sokhotski--Plemelj y en
$T=V_\alpha P_++M_GP_-$, pero no se puede importar la idempotencia del
modelo anterior. Si $X$ es un espacio de Banach, $S\in\mathcal B(X)$
y $P_\pm=(I_X\pm S)/2$, entonces
\[
 P_\pm^2-P_\pm=\tfrac14(S^2-I_X),\qquad
 P_+P_-=\tfrac14(I_X-S^2),\qquad P_++P_-=I_X.
\]
De hecho, el símbolo Mellin de $S_\Gamma$ en el caso de una sola rama,
$p=2$ y parámetro geométrico $\delta=0$ es
$\varepsilon_1\tanh(\pi x)$, donde $\varepsilon_1\in\{-1,1\}$
registra la orientación; su cuadrado
no es idénticamente uno \cite[teorema~5.3, p.~476]{Karlovich2017}.
Por tanto, en ese caso $P_\pm$ no son proyecciones algebraicas.
La conexión correcta con Haseman es la \emph{descomposición de las trazas}
por Plemelj, no una descomposición ortogonal del espacio.

\begin{cthnota}{(Dónde intervienen las hipótesis)}
La igualdad $S_\Gamma^2=I$ es exactamente la que da idempotencia;
$S_\Gamma^*=S_\Gamma$ es necesaria para ortogonalidad. El modelo
$L^2(\mathbb R)$ o $L^2(\mathbb T)$ satisface ambas propiedades en la
fuente citada. El problema de Haseman de 2017 actúa en $L^p(\Gamma)$
sobre una estrella abierta y sólo usa aquí la fórmula de Plemelj.
\end{cthnota}

\Needspace{25\baselineskip}
\section{Shifts unitarios y shifts de composición}

Comparamos una rotación que es unitaria con el operador de composición
que aparece en Haseman. El espacio de acción y la medida deciden qué
propiedad se puede afirmar.

En \cite[§1, p.~20]{BFK2006}, el círculo $\mathbb T$ actúa mediante una representación unitaria;
para $z,t\in\mathbb T$ y $\varphi\in L^2(\mathbb T)$ se escribe
\[
 U:\mathbb T\to\mathcal B(L^2(\mathbb T)),
 \qquad
 (U_z\varphi)(t)=\varphi(zt).
\]
La rotación preserva la medida de Haar $m$: al sustituir $u=zt$,
\[
 \|U_z\varphi\|_2^2=\int_{\mathbb T}|\varphi(zt)|^2\,dm(t)
 =\int_{\mathbb T}|\varphi(u)|^2\,dm(u)=\|\varphi\|_2^2.
\]
Como $U_z^{-1}=U_{z^{-1}}$, $U_z$ es unitario. En el problema de Haseman aparece en cambio
$V_\alpha f=f\circ\alpha$ sobre $L^p(\Gamma)$; Karlovich impone condiciones
de regularidad y control de la derivada del shift en cada rama
\cite[§2, p.~467]{Karlovich2017}. Estas condiciones permiten trabajar con
un operador acotado, pero no afirman que sea unitario.

\begin{center}
\begin{tikzpicture}[node distance=10mm]
\node[theorybox] (shift) {Idea común\\operador de composición / shift};
\node[researchbox,above right=of shift] (unit) {$U_z$ en $L^2(\mathbb T)$\\unitario};
\node[researchbox,below right=of shift] (hasa) {$V_\alpha$ en $L^p(\Gamma)$\\acotado bajo hipótesis del shift};
\node[theorybox,right=22mm of unit] (uspec) {Unitario $\Rightarrow\spec(U_z)\subset\mathbb T$};
\node[theorybox,right=22mm of hasa] (fred) {Se combina con $P_\pm$ y $M_G$ en el operador de Haseman};
\draw[arr] (shift)--(unit);
\draw[arr] (shift)--(hasa);
\draw[arr] (unit)--(uspec);
\draw[arr] (hasa)--(fred);
\end{tikzpicture}
\end{center}

\section{Funciones lentamente oscilantes y fibras de caracteres}

Para $a\in C_b(\mathbb R_+)$, escriba
$\operatorname{osc}(a,[r,2r])=\sup_{s,t\in[r,2r]}|a(s)-a(t)|$.
Karlovich define $\SO(\mathbb R_+)$ exigiendo que esta oscilación tienda a
cero tanto cuando $r\to0$ como cuando $r\to\infty$
\cite[§2, pp.~465--466]{Karlovich2017}. Para $I=[r,2r]$ y funciones
$a,b\in C_b(\mathbb R_+)$ se tiene
\[
\begin{aligned}
 \operatorname{osc}(a+b,I)&\le\operatorname{osc}(a,I)+\operatorname{osc}(b,I),\\
 \operatorname{osc}(ab,I)&\le\|a\|_\infty\operatorname{osc}(b,I)
                         +\|b\|_\infty\operatorname{osc}(a,I),\\
 \operatorname{osc}(\overline a,I)&=\operatorname{osc}(a,I).
\end{aligned}
\]
Si $a_n\to a$ uniformemente, entonces
$\operatorname{osc}(a,I)\le2\|a-a_n\|_\infty+
\operatorname{osc}(a_n,I)$; primero se fija $n$ grande y después se
deja tender $r$ al extremo. Así la clase es una subálgebra cerrada,
contiene las constantes y es estable bajo la involución puntual.
La identidad $\|\overline a a\|_\infty=\|a\|_\infty^2$ muestra que
es una $C^*$-álgebra conmutativa unital.

\begin{cthejemplo}{(Oscilación lenta sin límite ordinario)}
Queremos comprobar que la oscilación lenta en los extremos no obliga a
tener un límite ordinario. En $\mathbb R_+$ consideremos la función
\[
 a(r)=\sin\!\bigl(\log(1+|\log r|)\bigr),\qquad r>0,
\]
pertenece a $\SO(\mathbb R_+)$. Si $s,t\in[r,2r]$, entonces
$|\log s-\log t|\le\log2$. Ponga $h(u)=\log(1+|u|)$.
En el intervalo $I_r=[\log r,\log r+\log2]$, cuando está alejado de $0$,
el teorema del valor medio da
\[
 |h(\log s)-h(\log t)|
 \le \frac{\log2}{1+\min_{u\in I_r}|u|}.
\]
El denominador tiende a infinito tanto si $r\to0$ como si $r\to\infty$.
Como $|\sin v-\sin w|\le|v-w|$, la oscilación de $a$ en $[r,2r]$
tiende a cero en ambos extremos. Para ver que no hay límite ordinario,
defina
\[
 r_n^+=\exp\!\bigl(e^{\pi/2+2\pi n}-1\bigr),\qquad
 r_n^-=\exp\!\bigl(e^{3\pi/2+2\pi n}-1\bigr).
\]
Ambas sucesiones tienden a infinito, mientras que
$a(r_n^+)=1$ y $a(r_n^-)=-1$. Sus recíprocas tienden a cero y dan
los mismos dos valores, porque $a(r)=a(1/r)$.
Esta función, elaborada a partir de la definición de la fuente, muestra por
qué los caracteres de frontera de $\SO(\mathbb R_+)$ no se pueden sustituir
simplemente por evaluaciones $a(0)$ y $a(\infty)$.
\end{cthejemplo}

El espacio maximal $\Omega(\SO(\mathbb R_+))$ contiene las evaluaciones
en $r>0$ y fibras $M_0,M_\infty$ sobre los extremos
\cite[§2, p.~466]{Karlovich2017}. La transformada de Gelfand
$a\mapsto\widehat a$, $\widehat a(\xi)=\xi(a)$, permite describir
valores de estos caracteres sin fingir que existen $a(0)$ y
$a(\infty)$. En los cálculos de símbolos, esas fibras son parte del
espacio de parámetros; el mapa de Gelfand de los coeficientes y el
símbolo completo de un operador son mapas distintos.

\section{Generadores, cocientes y mapas de símbolos}

Karlovich \cite[ec.~(1.4), p.~3]{Karlovich2025} fija una curva estrellada
con cúspides $\Gamma$ de (1.1)--(1.2), un peso de potencia admisible $w$
de su §2.1, con $1<p<\infty$ y exponente
$-1/p<\alpha_{\mathrm{peso}}<1-1/p$, y considera la álgebra
de Banach generada por
\[
 \mathfrak B=
 \operatorname{alg}_{\mathbb R}
 \{M_a,S_\Gamma,C:\ a\in\SO^\diamond(\Gamma)\}
 \subset\mathcal B(\mathbb RL^p(\Gamma,w)).
\]
Aquí $M_a$ multiplica por $a$, $S_\Gamma$ es el operador singular de
Cauchy y $Cf=\overline f$ es conjugación compleja. El símbolo
$\mathbb RL^p(\Gamma,w)$ significa el espacio $L^p$ considerado sobre
los escalares reales: para $f\in L^p(\Gamma,w)$ se tiene
$C(if)=\overline{if}=-i\overline f=-iCf$, mientras que
$C(rf)=rCf$ para $r\in\mathbb R$. Así $C$ es lineal sobre
$\mathbb R$, pero no sobre $\mathbb C$. Por eso esta álgebra no se convierte automáticamente en una
$C^*$-álgebra compleja ni se puede aplicar sin más el teorema espectral
complejo de la Parte 1. La notación $\SO^\diamond(\Gamma)$ pertenece
al artículo de 2025 y no designa la misma álgebra
$\SO(\mathbb R_+)$ del artículo de 2017.
El artículo de 2025 llama $\alpha$ al exponente del peso; aquí el
subíndice evita confundirlo con el shift $\alpha$ de Karlovich 2017.

\begin{center}
\begin{tikzpicture}[node distance=5mm and 8mm]
\node[researchbox,text width=3.25cm] (gen) {$M_a,\ S_\Gamma,\ C$\\generadores tipados};
\node[theorybox,text width=3.25cm,right=of gen] (alg) {$\mathfrak B$\\cierre real bajo suma,\\producto y norma};
\node[researchbox,text width=3.25cm,right=of alg] (sym) {Cociente y símbolos\\criterio en una álgebra fijada};
\draw[arr] (gen)--(alg);
\draw[arr] (alg)--(sym);
\end{tikzpicture}
\end{center}

La Tarea 1.3 enseña qué significa generar y cerrar una álgebra. La
Tarea 1.4 enseña a pasar a un cociente. En el artículo de 2025, el
paso adicional es construir mapas matriciales de símbolo y probar
que detectan Fredholmness. El teorema~6.2 trata la subálgebra compleja
$\mathfrak A=\operatorname{alg}_{\mathbb C}
\{M_a,S_\Gamma,CS_\Gamma C:a\in\SO^\diamond(\Gamma)\}$;
la representación matricial compleja del teorema~7.1 permite después formular el criterio
para $N\in\mathfrak B$ mediante el mapa de símbolo $Y$ del teorema~7.2
\cite[pp.~25--26]{Karlovich2025}. Este último exige pertenencia a
$\mathfrak B$ y las hipótesis de curva, peso y coeficientes del artículo;
no certifica por sí solo un operador de Haseman con shift.

\begin{cthnota}{(Cuatro mapas relacionados y diferentes)}
Un homomorfismo de álgebras preserva suma y producto. Un
$*$-homomorfismo preserva además la involución. La transformada de
Gelfand $a\mapsto\widehat a$ de una álgebra conmutativa evalúa en sus
caracteres. Un mapa de símbolo de Fredholm se construye para una
álgebra concreta de operadores, frecuentemente después de pasar
módulo compactos, y puede tomar valores matriciales. Estas
construcciones comparten operaciones algebraicas, pero sus dominios,
codominios e hipótesis no son intercambiables.
\end{cthnota}

\Needspace{20\baselineskip}
\section{Ruta completa desde el problema hasta un criterio}

En el diagrama, azul marca el lenguaje del curso, verde el objeto de
investigación y la flecha hacia símbolos presupone que se ha construido
el cálculo correspondiente para la álgebra concreta. Aquí
$X=L^p(\Gamma)$, de modo que $T+\mathcal K(X)$ es una clase de la
álgebra de Calkin asociada al mismo espacio de acción. El criterio de
Karlovich 2017 exige además las hipótesis de su teorema 7.1
\cite[pp.~487--489]{Karlovich2017}.

\begin{center}
\begin{tikzpicture}[node distance=5mm and 7mm]
\node[researchbox,text width=3.25cm] (h) {[investigación]\\frontera de Haseman};
\node[researchbox,text width=3.25cm,right=of h] (p) {[investigación]\\Sokhotski--Plemelj};
\node[researchbox,text width=3.25cm,right=of p] (t) {[investigación]\\$T=V_\alpha P_++M_GP_-$};
\node[theorybox,text width=3.25cm,below=of t] (b) {[curso]\\$T\in\mathcal B(L^p(\Gamma))$};
\node[theorybox,text width=3.25cm,left=of b] (q) {[curso + puente]\\$T+\mathcal K(X)$ en $\mathcal Q(X)$};
\node[researchbox,text width=3.25cm,left=of q] (f) {[puente]\\invertibilidad de clase\\$\Longleftrightarrow$ Fredholm};
\node[researchbox,text width=3.25cm,below=of f] (s) {[investigación]\\símbolos y localización};
\node[researchbox,text width=3.25cm,right=of s] (c) {[investigación]\\criterio de Fredholm\\bajo hipótesis del artículo};
\draw[arr] (h)--(p);
\draw[arr] (p)--(t);
\draw[arr] (t)--(b);
\draw[arr] (b)--(q);
\draw[arr] (q)--(f);
\draw[arr] (f)--(s);
\draw[arr] (s)--(c);
\end{tikzpicture}
\end{center}

La última fila describe cómo se \emph{comprueba} Fredholmness en la
literatura. No afirma que todo operador Fredholm tenga automáticamente
el símbolo concreto del artículo, ni que el criterio de 2017 se aplique
a un operador con cúspides de 2025.

\Needspace{19\baselineskip}
\section{Una ruta de lectura para el manuscrito de investigación}

Cuando aparezca una construcción operatorial en el manuscrito, conviene hacer las siguientes
preguntas en este orden:
\begin{enumerate}
\item ¿En qué álgebra de Banach vive el operador?
\item ¿Cuál es la unidad y cuáles son los generadores relevantes?
\item ¿Se está trabajando exactamente o módulo un ideal, normalmente $\K(X)$?
\item ¿La subálgebra considerada es conmutativa o sólo posee una subálgebra central conmutativa?
\item ¿Aparecen caracteres, espacio maximal o transformada de Gelfand?
\item ¿Hay una estructura de $C^*$-álgebra compleja (típicamente en $L^2$) que permita usar normalidad, unitarios,
      proyecciones o cálculo funcional?
\item ¿Qué parte del criterio de Fredholm equivale a invertibilidad en un cociente o de un símbolo?
\end{enumerate}

Estas preguntas obligan a distinguir el lenguaje común de las hipótesis que autorizan
cada teorema. Por ejemplo, un criterio publicado para una álgebra de operadores con cúspides
no se aplica automáticamente a un operador con shift de otra geometría.

\chapter{Mapa de dependencias entre resultados, tareas e investigación}

\section{Qué significa una flecha en este mapa}

\begin{context}
Una flecha continua indica que el resultado de origen se usa en la demostración o en la
construcción del resultado de destino. Una flecha discontinua indica una aplicación o comparación
conceptual: no afirma que el resultado de destino se deduzca del primero. Las hipótesis escritas
junto a cada ruta son parte de la conexión.
\end{context}

\section{Rutas que intervienen en demostraciones}

\begin{dependencias}
\textbf{Continuidad e inversos.}
\[
\boxed{\text{Serie de Neumann (Parte 1)}}
\longrightarrow
\boxed{\text{T1.1(b): inversión continua}}
\longrightarrow
\boxed{\text{Prop. 2.17: permanencia espectral}}.
\]
En la última flecha se usa que $(b+i/n)^{-1}\to b^{-1}$ en la álgebra ambiente y que la
$C^*$-subálgebra es cerrada. T1.1(a) aporta, por otra vía, continuidad del producto para T1.3,
T1.7 y T2.8.
\end{dependencias}

\begin{dependencias}
\textbf{Reducción conmutativa.}
\[
\begin{gathered}
\boxed{\text{T1.2: inversos que conmutan}}
\quad+\quad
\boxed{\text{T1.3: álgebra generada}}\\[2mm]
\Downarrow\\[2mm]
\boxed{\text{T1.9: subálgebra inversamente cerrada}}\\[2mm]
\Downarrow\\[2mm]
\boxed{\text{Prop. 2.13: espectro real}}.
\end{gathered}
\]
La última prueba añade la identidad $C^*$, la realidad de los caracteres en autoadjuntos y la
descripción del espectro mediante caracteres. T1.9 aporta la subálgebra conmutativa que conserva
el espectro; por sí solo no demuestra que el espectro sea real.
\end{dependencias}

\begin{dependencias}
\textbf{De las funciones al cálculo funcional.}
\[
\begin{gathered}
\boxed{\text{T1.3(d): }C^*(a,1)\text{ conmutativa para $a$ normal}}\\[2mm]
\quad+\quad
\boxed{\text{Prop. 2.17: permanencia espectral}}
\quad+\quad
\boxed{\text{Teo. 2.20: Gelfand--Naimark}}\\[2mm]
\Downarrow\\[2mm]
\boxed{\text{Teo. 2.22: cálculo funcional}}.
\end{gathered}
\]
T1.10 calcula el espectro de una función continua y, combinado con el cálculo funcional y la
permanencia espectral, da T2.4. T2.2(a,c) proporciona las direcciones directas de T2.6(c,b);
las recíprocas usan normalidad y cálculo funcional.
\end{dependencias}

\begin{dependencias}
\textbf{Continuidad al variar el operador.}
\[
\begin{gathered}
\boxed{\text{T1.7: límite de normales normal}}
\quad+\quad
\boxed{\text{Stone--Weierstrass}}\\[2mm]
\quad+\quad
\boxed{\text{cálculo funcional isométrico}}\\[2mm]
\Downarrow\\[2mm]
\boxed{\text{T2.8: }f(x_n)\to f(x)}.
\end{gathered}
\]
La convergencia aquí es en norma. T2.9 estudia otra situación: una red creciente y acotada de
autoadjuntos converge fuertemente, es decir, vector por vector en norma de $H$.
\end{dependencias}

\section{Cobertura de los ejercicios}

\small
\begin{longtable}{>{\bfseries}p{1.3cm}p{4.3cm}p{8.2cm}}
\toprule
Tarea & Resultado demostrado & Uso exacto o alcance de la conexión\\
\midrule
\endhead
T1.1 & Producto e inversión continuos & El producto permite T1.3 y T1.7; la inversión se usa en la Prop. 2.17 y el producto de nuevo en T2.8.\\
T1.2 & Conmutación con inversos y adjuntos & Se usa al añadir inversos en la construcción de T1.9.\\
T1.3 & Álgebra generada y conmutación & Para un normal $a$, construye $C^*(a,1)$ conmutativa, usada en las Props. 2.18 y 2.22.\\
T1.4 & Cociente por ideal cerrado & Hace de $\mathcal B(X)/\K(X)$ una álgebra de Banach; la equivalencia Fredholm--invertibilidad es el teorema externo de Atkinson.\\
T1.5 & $C_0(X)$ cerrado en $C_b(X)$ & Si $X$ no es compacto, es no unital; la clasificación general de $C_0(X)$ requiere la versión no unital del teorema de la Sección 2.\\
T1.6 & $\spec(x^*)=\overline{\spec(x)}$ & Simetría válida en álgebras involutivas; la realidad del espectro autoadjunto exige además el argumento $C^*$ de la Prop. 2.13.\\
T1.7 & Normalidad cerrada & Primera hipótesis que T2.8 necesita para definir $f(x)$ en el límite.\\
T1.8 & Inclusión espectral por homomorfismos & Aplicada a $a^*a$, entra en la prueba de contractividad de $*$-homomorfismos (Teo. 2.8).\\
T1.9 & Subálgebra conmutativa inversamente cerrada & Se usa expresamente en la prueba de la Prop. 2.13 para preservar el espectro de un autoadjunto.\\
T1.10 & $\spec(f)=f(X)$ para $X$ compacto & Calcula el espectro de $f$ en la prueba de T2.4; ilustra el modelo $C(X)$ de Gelfand--Naimark.\\
T1.11 & $\mathcal B(H)$ y shift unilateral & El inciso (a) da el modelo concreto de T2.1; (c,d) separan espectro puntual aproximado y espectro, y T2.10 completa el cálculo de $\spec(S)$.\\
T1.12 & Álgebra del disco & Muestra que una inclusión unital y cerrada de álgebras de Banach puede cambiar el espectro; la inclusión no es de $C^*$-subálgebras.\\
\midrule
T2.1 & $C_b(X)\hookrightarrow\mathcal B(\ell^2(X))$ & Verifica directamente un $*$-homomorfismo isométrico; ejemplo concreto del Teo. 2.27.\\
T2.2 & Idempotentes, productos, unitarios & (a,c) dan direcciones directas de la Prop. 2.28; (b) estima radios de productos conmutativos y no se necesita para esa proposición.\\
T2.3 & Producto acotado y suma $c_0$ & Construye una $C^*$-álgebra y su ideal cerrado; T1.4 permite considerar el cociente como álgebra de Banach.\\
T2.4 & Mapeo espectral continuo & Depende del Teo. 2.22, la Prop. 2.17 y T1.10; se aplica al espectro de funciones de normales.\\
T2.5 & $X\cong\Omega(C(X))$ & Da la unicidad del espacio compacto en el Cor. 2.21.\\
T2.6 & Autoadjuntos, unitarios y proyecciones normales & Reconstruye la Prop. 2.28 mediante cálculo funcional y muestra con contraejemplos la necesidad de normalidad.\\
T2.7 & Factor unitario de un invertible & Usa la raíz positiva de la Prop. 2.30 y la invertibilidad de $a^*a$ para obtener $a=u|a|$.\\
T2.8 & Continuidad de $f(x_n)$ & Combina T1.7, Stone--Weierstrass y cálculo funcional; la convergencia es en norma.\\
T2.9 & Convergencia monótona de autoadjuntos & Usa orden positivo, polarización y Riesz; concluye convergencia fuerte, sin afirmar convergencia en norma.\\
T2.10 & Espectro del shift unilateral & Amplía T1.11(d) a $\spec(S)=\overline{\mathbb D}$; $S$ no es normal.\\
\bottomrule
\end{longtable}
\normalsize

\section{Dos puentes controlados hacia Haseman}

\begin{center}
\begin{tikzpicture}[node distance=8mm]
\node[taskbox,text width=2.6cm] (q) {T1.4\\cociente por ideal};
\node[theorybox,text width=2.7cm,right=of q] (calkin) {$\K(X)$ cerrado\\$\mathcal B(X)/\K(X)$};
\node[researchbox,text width=2.6cm,right=of calkin] (class) {Clase de $A$\\$A+\K(X)$};
\node[flowend,text width=2.7cm,right=of class] (fred) {Atkinson\\invertible $\Leftrightarrow$ Fredholm};
\draw[arr] (q)--(calkin);
\draw[arr] (calkin)--(class);
\draw[arr] (class)--(fred);
\end{tikzpicture}
\end{center}

\noindent Para $X$ de Banach e infinito-dimensional, la primera flecha usa que $\K(X)$ es un ideal cerrado. La última es el teorema de
Atkinson, no una consecuencia de T1.4. En los operadores de la línea Haseman, esta formulación
permite plantear el problema de Fredholm en un cociente; véase
\cite[p.~2]{Karlovich2025}.
Como comprobación concreta, en $X=\ell^2$ tome $P$ como proyección
sobre el primer vector canónico $e_1$ y ponga $A=I_X-P$.
Entonces $A e_1=0$, pero
$A+\K(X)=I_X+\K(X)$ porque $P$ es compacto. El cálculo muestra qué
información conserva el cociente y por qué T1.4 prepara la formulación
de Atkinson.

\begin{center}
\begin{tikzpicture}[node distance=9mm and 8mm]
\node[theorybox,text width=3.0cm] (model) {Modelo $L^2(\Gamma_0)$,\\$\Gamma_0=\mathbb R$ o $\mathbb T$\\$S_{\Gamma_0}^2=I$, $S_{\Gamma_0}^*=S_{\Gamma_0}$};
\node[taskbox,text width=3.0cm,right=of model] (orth) {$P_{\Gamma_0}^\pm=(I\pm S_{\Gamma_0})/2$\\proyecciones ortogonales};
\draw[arr] (model)--(orth);
\node[researchbox,text width=3.0cm,below=of model] (plem) {Estrella abierta de Karlovich\\Sokhotski--Plemelj};
\node[researchbox,text width=3.0cm,right=of plem] (trace) {$P_\pm=(I\pm S_\Gamma)/2$\\componentes de trazas};
\node[flowend,text width=3.0cm,right=of trace] (has) {$T=V_\alpha P_+ + GP_-$\\ecuación de Haseman};
\draw[arr] (plem)--(trace);
\draw[arr] (trace)--(has);
\draw[mergearr,dash pattern=on 4pt off 2pt] (orth.south)--(trace.north);
\end{tikzpicture}
\end{center}

\noindent La flecha discontinua sólo compara la forma de las fórmulas.
En el modelo de Bastos--Fernandes--Karlovich \cite[§2, p.~22]{BFK2006},
$P_{\Gamma_0}^\pm$ son proyecciones ortogonales. En la estrella abierta de
Karlovich 2017, la reducción de frontera a $T$ procede de
Sokhotski--Plemelj \cite[pp.~463--464]{Karlovich2017}; la identidad
$S_\Gamma^2=I$ no vale en general allí, según el teorema 5.3 de esa fuente.
Por tanto, los operadores $P_\pm$ de esa ecuación no se llaman aquí
idempotentes ni proyecciones.
En el caso de comprobación $\alpha=\mathrm{id}$ y $G=1$, la fórmula de
Haseman da $T=P_++P_-=I$ por cancelación de los dos términos $S_\Gamma$.
Este cálculo usa la fórmula de Plemelj, sin requerir $S_\Gamma^2=I$.

\begin{reflection}
Para comprobar una flecha, localiza la línea de la demostración que usa su origen. Por ejemplo,
identifica dónde entra T1.9 en la realidad del espectro y dónde entra T1.10 en T2.4. Después
verifica en el modelo $L^2(\Gamma_0)$ cómo $S_{\Gamma_0}^2=I$ produce
idempotentes y cómo $S_{\Gamma_0}^*=S_{\Gamma_0}$ los hace proyecciones
ortogonales. Contrasta esas hipótesis con las de la estrella abierta.
\end{reflection}

\chapter*{Índice de conexiones con la literatura de investigación}
\addcontentsline{toc}{chapter}{Índice de conexiones con la literatura de investigación}

Esta tabla permite regresar a una fuente después de estudiar el concepto del curso.
La palabra \emph{relacionada} no afirma que el objeto avanzado sea una
generalización del resultado elemental. En las filas de operadores,
$X$ denota un espacio de Banach complejo y $H$ uno de Hilbert.

\footnotesize
\begin{longtable}{>{\raggedright\arraybackslash}p{2.5cm}>{\raggedright\arraybackslash}p{5.5cm}>{\raggedright\arraybackslash}p{3.6cm}>{\raggedright\arraybackslash}p{2.0cm}}
\toprule
Concepto del curso & Aparición y qué reconocer & Fuente precisa & Tipo\\
\midrule
\endfirsthead
\toprule
Concepto del curso & Aparición y qué reconocer & Fuente precisa & Tipo\\
\midrule
\endhead
Álgebra de Banach &
$T=V_\alpha P_++M_GP_-\in\mathcal B(L^p(\Gamma))$:
suma y composición de operadores acotados. &
Karlovich 2017, pp.~463--464, (1.1)--(1.2) \cite{Karlovich2017}. &
Ejemplo directo\\
Álgebra $C^*$ conmutativa &
$C_b(\Gamma)$ y $\SO(\Gamma)$ son álgebras concretas de coeficientes,
con conjugación puntual. &
Karlovich 2017, p.~464 \cite{Karlovich2017}. &
Ejemplo directo\\
Ideal y cociente &
$J=\mathcal K(X)$ en $\A=\mathcal B(X)$ da la álgebra de Calkin;
la invertibilidad de la clase caracteriza Fredholmness. &
Karlovich 2025, p.~2 \cite{Karlovich2025}. &
Uso de propiedad\\
Espectro &
$\sigma_{\mathrm{ess}}(A)$ se define aquí mediante la clase en Calkin:
se distingue de $\sigma_{\mathcal B(X)}(A)$. &
Elaboración de estas notas a partir del cociente citado en Karlovich 2025, p.~2 \cite{Karlovich2025}. &
Construcción relacionada\\
Caracteres y Gelfand &
Las fibras $M_0,M_\infty$ de $\Omega(\SO(\mathbb R_+))$ codifican
comportamientos límite de coeficientes. &
Karlovich 2017, §2, pp.~465--466 \cite{Karlovich2017}. &
Uso de propiedad\\
Proyección en $C^*$ &
En $L^2(\Gamma)$, para $\Gamma=\mathbb R$ o $\mathbb T$,
$P_\Gamma^\pm$ son ortogonales.
La misma fórmula en la estrella abierta no garantiza idempotencia. &
BFK 2006, §2, p.~22 \cite{BFK2006}; Karlovich 2017, teorema~5.3, p.~476 \cite{Karlovich2017}. &
Ejemplo con hipótesis\\
Unitarios y shifts &
Las rotaciones $U_z$ son unitarias; el shift $V_\alpha$ de Haseman
requiere otras hipótesis de acotación. &
BFK 2006, §1, p.~20 \cite{BFK2006}; Karlovich 2017, §2, p.~467 \cite{Karlovich2017}. &
Comparación\\
Álgebra generada &
Multiplicadores, $S_\Gamma$ y conjugación generan una álgebra real
de operadores; tipar la conjugación evita llamarla compleja. &
Karlovich 2025, (1.4), p.~3 \cite{Karlovich2025}. &
Ejemplo directo\\
Homomorfismos &
Los mapas matriciales de símbolo ayudan a comprobar Fredholmness
en una álgebra concreta; no son automáticamente $*$-homomorfismos. &
Karlovich 2025, teoremas~6.2 y 7.2, pp.~25--26 \cite{Karlovich2025}. &
Construcción relacionada\\
Fredholmness &
El criterio para $T$ llega tras la representación Mellin módulo
compactos y sus hipótesis de símbolo. &
Karlovich 2017, teorema~7.1, pp.~487--489 \cite{Karlovich2017}. &
Para leer más\\
\bottomrule
\end{longtable}
\normalsize


\renewcommand{\bibname}{Referencias de trabajo}
\begin{thebibliography}{99}
\addcontentsline{toc}{chapter}{Referencias de trabajo}
\bibitem{CastillejosParte1} J. Castillejos, \emph{$C^*$-álgebras: Parte 1. Teoría espectral en álgebras de Banach}, notas de curso proporcionadas por el usuario.
\bibitem{CastillejosCap2} J. Castillejos, \emph{$C^*$-álgebras: Sección 2}, notas de curso proporcionadas por el usuario.
\bibitem{Karlovich2017} Yu. I. Karlovich, \emph{The Haseman Boundary Value Problem with Slowly Oscillating Coefficients and Shifts}, Operator Theory: Advances and Applications 259 (2017), 463--500.
\bibitem{Karlovich2025} Yu. I. Karlovich, \emph{Algebras of singular integral operators with complex conjugation on curves with cusps}, Bol. Soc. Mat. Mex. 31 (2025), 102.
\bibitem{BFK2006} M. A. Bastos, C. A. Fernandes, Yu. I. Karlovich, \emph{$C^*$-Algebras of Integral Operators with Piecewise Slowly Oscillating Coefficients and Shifts Acting Freely}, Integral Equations Operator Theory 55 (2006), 19--67.
\end{thebibliography}

\end{document}
